\documentclass[11pt]{amsart}
\usepackage[margin=2.5cm]{geometry}
\usepackage{amsmath,amssymb,amsthm,mathtools}
\usepackage{booktabs}
\usepackage{enumitem}
\usepackage[dvipsnames]{xcolor}
\usepackage[colorlinks=true,linkcolor=magenta,urlcolor=magenta,citecolor=magenta,filecolor=magenta,menucolor=magenta,runcolor=magenta]{hyperref}
\hypersetup{pdftitle={On the maximality conjecture for quantum permutation groups},pdfauthor={Claude and Teo}}
\setlist{itemsep=2pt,topsep=4pt}

\newtheorem{theorem}{Theorem}[section]
\newtheorem{proposition}[theorem]{Proposition}
\newtheorem{lemma}[theorem]{Lemma}
\newtheorem{corollary}[theorem]{Corollary}
\newtheorem{conjecture}[theorem]{Conjecture}
\newtheorem{question}[theorem]{Question}
\theoremstyle{definition}
\newtheorem{definition}[theorem]{Definition}
\newtheorem{example}[theorem]{Example}
\newtheorem{fact}[theorem]{Computational fact}
\theoremstyle{remark}
\newtheorem{remark}[theorem]{Remark}
\theoremstyle{plain}
\newtheorem{introtheorem}{Theorem}
\renewcommand{\theintrotheorem}{\Alph{introtheorem}}

\numberwithin{equation}{section}
\numberwithin{table}{section}
\setcounter{tocdepth}{1}
\emergencystretch=3em
\raggedbottom
\makeatletter
\renewcommand{\tocsection}[3]{\indentlabel{\@ifnotempty{#2}{\makebox[1.8em][l]{#2.}}}#3}
\makeatother

\newcommand{\C}{\mathbb C}
\newcommand{\Q}{\mathbb Q}
\newcommand{\Z}{\mathbb Z}
\newcommand{\F}{\mathbb F}
\newcommand{\T}{\mathbb T}
\newcommand{\TL}{\mathrm{TL}}
\newcommand{\HL}{\mathrm{HL}}
\newcommand{\Br}{\mathrm{Br}}
\newcommand{\Fix}{\mathrm{Fix}}
\newcommand{\NC}{\mathcal{NC}}
\newcommand{\Pp}{\mathcal{P}}
\DeclareMathOperator{\Aut}{Aut}
\DeclareMathOperator{\spn}{span}
\DeclareMathOperator{\supp}{supp}
\DeclareMathOperator{\Proj}{Proj}

\begin{document}

\title[Maximality conjecture]{On the maximality conjecture for quantum permutation groups}
\author{Claude and Teo}

\subjclass[2020]{46L65 (20G42, 05A18)}
\keywords{Quantum permutation group, maximality conjecture, noncrossing partition}

\begin{abstract}
A well-known conjecture, central in quantum permutation group theory, and going back to around 2006, states that $S_N\subset S_N^+$ is maximal. We revisit here this conjecture, and related questions, 20 years after, with 2026 knowledge and tools.
\end{abstract}

\maketitle

\tableofcontents

\section*{Introduction}

The quantum permutation group $S_N^+$ was introduced by Wang in 1998~\cite{wa}, as the universal compact quantum group acting on $N$ points. The algebra of functions on $S_N^+$ is generated by the entries of an $N\times N$ magic unitary, and its quotient by the commutator ideal is the algebra of functions on $S_N$, so that we have an inclusion $S_N\subset S_N^+$. This inclusion is an equality for $N\le3$, while for $N\ge4$ the quantum group $S_N^+$ is infinite, and for $N\ge5$ it is not even coamenable. Its representation theory is that of $SO_3$, with intertwiners given by the noncrossing partitions~\cite{ba}, instead of all partitions for $S_N$. When classifying the quantum subgroups of $S_4^+$, Banica and Bichon~\cite{bb} found nothing between $S_4$ and $S_4^+$, and this led, around 2006, to the following statement.

\smallskip
\noindent\textbf{Maximality conjecture.} \emph{For every $N$, the inclusion $S_N\subset S_N^+$ is maximal, in the sense that there is no intermediate quantum group $S_N\subsetneq G\subsetneq S_N^+$.}

\smallskip
The conjecture holds trivially for $N\le3$, it holds for $N=4$ by~\cite{bb}, and for $N=5$ by~\cite{bah}, the proof there relying on the classification of subfactors of index $\le5$~\cite{imps,jms}. For $N\ge6$ it is open. Since there is no easy quantum group strictly between $S_N$ and $S_N^+$~\cite{bcs}, the conjecture lies beyond easiness, and it is a natural first test case for the classification of compact quantum groups beyond easiness. We revisit here this conjecture, and some related questions, twenty years after, with 2026 knowledge and tools. After preliminaries in Section~\ref{sec:maximality}, the paper is organized around three approaches:

\smallskip
(1) \emph{Subgroups.} Given a quantum subgroup $G\subset S_N^+$ with $G\not\subset S_N$, the conjecture predicts $\langle S_N,G\rangle=S_N^+$. This is the natural source of counterexamples. It is also of independent interest, the equality $\langle S_N,G\rangle=S_N^+$ being a non-trivial result about $G$ itself. We test in this way group duals, which by a theorem of Bichon~\cite{bi} act through ``Fourier families'' of abelian groups, and reduce to just-non-abelian groups (Section~\ref{sec:duals}), quantum automorphism groups of graphs (Section~\ref{sec:graphs}), twists of $S_N$ and of $PU_n$ (Section~\ref{sec:twists}), and the quantum groups coming from complex Hadamard matrices (Section~\ref{sec:hadamard}). No counterexample appears.

\smallskip
(2) \emph{Exotic quantum groups.} Assume that $S_N\subsetneq G\subsetneq S_N^+$ exists, and develop the algebraic and analytic theory of $G$: representations and moments, orbitals, diagonal algebras, coamenability, and the finite case. The aim is twofold: having this theory available, should such a $G$ exist, and finding a contradiction coming from it. This is done in Section~\ref{sec:exotic}.

\smallskip
(3) \emph{Partitions.} By Tannakian duality, the intermediate quantum groups $S_N\subset G\subset S_N^+$ correspond to tensor categories $\spn(NC)\subset\mathcal C\subset\spn(P)$, made of linear combinations of partitions. Given such a $\mathcal C$, we look at the first level $k$ where it differs from $\spn(NC)$, and we write down linear, quadratic and cubic necessary conditions on the new vectors at level $k$, which we then solve by exact computer algebra, uniformly in $N$ whenever possible. This is our main contribution, developed in Sections~\ref{sec:method}, \ref{sec:six} and~\ref{sec:higher}.

\smallskip
The maximality conjecture has some interesting ramifications, which we discuss as well:

\smallskip
(A) \emph{Orthogonal case.} Inside Wang's free orthogonal group $O_N^+$ we have the half-liberation $O_N\subset O_N^*\subset O_N^+$, given by the relations $abc=cba$, and $O_N\subset O_N^*$ is known to be maximal~\cite{bbcc}. The conjecture here is that this chain is complete, in the sense that $O_N^*$ is the only intermediate quantum group $O_N\subsetneq G\subsetneq O_N^+$. We discuss this in Sections~\ref{sec:orthproblem} and~\ref{sec:pairings}, by adapting the above approaches (1) and (3), with group duals and twists as seeds, and with pairings in place of partitions.

\smallskip
(B) \emph{Projective versions.} For the projective versions $PO_N\subset H\subset(PO_N^+=PU_N^+)$, it is known that $PO_N^*=PU_N$ is the unique intermediate projectively easy quantum group~\cite{bme}, and conjecturally it is the unique intermediate quantum group. The new feature here is that $H$ is not necessarily the projective version of a quantum group $O_N\subset G\subset O_N^+$, and the level method must be adapted to this, with fewer symmetries. We discuss these questions in Section~\ref{sec:related}.

\smallskip
Our main results concern approach (3). The computations behind them are exact, and organized as certificates which can be checked independently, as explained in Section~\ref{sec:method}. First, the conjecture holds up to level $7$:

\begin{introtheorem}
Let $N\ge6$, and $S_N\subsetneq G\subset S_N^+$. Then for any $k\le7$ we have
\[
\Fix_G(u^{\otimes k})=\spn\big\{\xi_\pi:\pi\in NC(k)\big\}.
\]
Thus the first seven moments of the main character of $G$ are the Catalan numbers, as for $S_N^+$.
\end{introtheorem}

The study of the conjecture level by level goes back to~\cite{bah}, where the notion of easiness level was introduced, and was substantially developed by McCarthy, whose state-space approach to quantum permutations~\cite{mc0} led him to the program in~\cite{mc1,mc2}, with a number of results at levels $\le6$. Levels $k\le5$ were also announced by Freslon and Speicher~\cite{fs}, and we reprove them. Levels $6$ and $7$ are new, with a proof uniform in $N$. At level $8$ the method meets real resistance, but still gives a strong structural result:

\begin{introtheorem}
Let $7\le N\le37$, and let $S_N\subsetneq G\subset S_N^+$ first differ from $S_N^+$ at level $8$. Then the elements of $\Fix_G(u^{\otimes8})$ are, modulo $\spn(NC(8))$, invariant under the rotation of the $8$ legs by $4$ positions.
\end{introtheorem}

Pushing the rotation symmetry further, we can close level $8$ computationally, modulo primes, at each $8\le N\le37$: for these $N$ the conjecture holds up to level $8$ (Section~\ref{sec:higher}). Level $9$ is harder, with a kernel of dimension $307$ against $66$, and no half-turn. There a different tool, imported from the orthogonal problem of Section~\ref{sec:pairings}, does the job: the candidates split into nine rotation blocks, each of them has a real form, and a positive definite combination of the quadratic constraints on each real form excludes the whole block, by a certificate. This closes level $9$, numerically, at $N=9,\dots,13$ and $N=17,19,23$, so that at these $N$ the conjecture holds through level $9$ (Theorem~\ref{thm:nine}), the highest level reached so far.

Regarding approach (1), among our results is the following statement about graphs, where $G^+(X)$ denotes the quantum automorphism group of a graph $X$, and where the repetition lemma states that $\langle S_{4m},S_4^+\rangle=S_{4m}^+$, with $S_4^+$ acting on $4m$ points via $u\mapsto u\otimes1_m$:

\begin{introtheorem}
Let $X$ be a vertex-transitive graph with $N\le13$ vertices. Then
\[
\langle S_N,G^+(X)\rangle\in\{S_N,S_N^+\},
\]
except possibly for the cube, its complement, and four complementary pairs of graphs with $12$ vertices, where this holds provided that the repetition lemma mentioned above holds at $m=2,3$.
\end{introtheorem}

In the orthogonal case, a generation result of Brannan, Collins and Vergnioux~\cite{bcv}, namely $\langle O_N,O_{N-1}^+\rangle=O_N^+$ for $N\ge4$, allows an induction on $N$ which stops at $N=3$, where their method breaks down. Our main results concern the anticommutative twist $O_N^{-1}$, and the levels of $O_N^*\subset G\subset O_N^+$:

\begin{introtheorem}
$\langle O_N,O_N^{-1}\rangle$ agrees with $O_N^+$ up to level $\max(20,2N)$, and equals $O_N^+$ for all $N$, provided that $\langle O_3,O_2^+\rangle=O_3^+$.
\end{introtheorem}

So the twist problem reduces to the single generation statement at $N=3$ left open in~\cite{bcv}, which we verify up to level $20$. As for the levels themselves, the orthogonal version of approach (3), with uncappable vectors and positivity certificates, gives:

\begin{introtheorem}
Let $N\ge3$, and $O_N^*\subset G\subsetneq O_N^+$. If the fixed vectors of $G$ differ from those of $O_N^+$ at some level $\le10$, then $G=O_N^*$.
\end{introtheorem}

Many questions remain open, and they are stated along the way, at their natural places. The main problem, of course, is that of going beyond level $9$ in approach (3), and beyond level $7$ uniformly in $N$, and here the main difficulty is computational.

\smallskip
\noindent\emph{Acknowledgments.} We would like to thank JP McCarthy for useful discussions.


%======================================================================
% Section: maximality
%======================================================================

\section{Maximality and three approaches}\label{sec:maximality}

\subsection{Quantum permutation groups}

Let $N\in\mathbb N$. A \emph{magic unitary} is a square matrix $u=(u_{ij})\in M_N(A)$ over a $C^*$-algebra $A$ whose entries are projections summing to $1$ on each row and each column. Wang showed~\cite{wa} that the universal $C^*$-algebra
\[
C(S_N^+)=C^*\Big((u_{ij})_{i,j=1,\dots,N}\ \Big|\ u=\text{magic}\Big),
\]
with comultiplication $\Delta(u_{ij})=\sum_ku_{ik}\otimes u_{kj}$, counit $\varepsilon(u_{ij})=\delta_{ij}$ and antipode $S(u_{ij})=u_{ji}$, defines a compact quantum group $S_N^+$ in the sense of Woronowicz~\cite{wo1}, the \emph{quantum permutation group}. It is the universal compact quantum group acting on the set $\{1,\dots,N\}$, in the sense that it is the quantum automorphism group of the algebra $\C^N$ endowed with its uniform trace. The commutative quotient of $C(S_N^+)$ is $C(S_N)$, with $u_{ij}$ mapping to the characteristic function of $\{\sigma:\sigma(j)=i\}$. So we have an inclusion of compact quantum groups
\[
S_N\subset S_N^+ .
\]

\emph{Small $N$.} For $N=1,2,3$ the entries of a magic unitary automatically commute, and $S_N^+=S_N$~\cite{wa}. At $N=4$ things change: $C(S_4^+)$ is infinite-dimensional. A good model comes from the Pauli matrices. The magic unitary whose entries are the rank-one projections onto $c_ix c_j$, with $x\in SU_2$ and $c_1,\dots,c_4$ the Pauli matrices, gives a faithful representation $C(S_4^+)\to M_4(C(SU_2))$. This can be used to integrate over $S_4^+$ \cite{bc}. In fact $S_4^+=SO_3^{-1}$, a twist of $SO_3$, and all the quantum subgroups of $S_4^+$ can be classified this way~\cite{bb}. For $N\ge5$, $S_N^+$ is not coamenable. One way to see it is that it contains the dual of $\Z_2*\Z_3$, which is a non-amenable group, acting by Fourier matrices on $2+3$ of the $N$ points.

The liberation operation $S_N\to S_N^+$ is best understood in representation theory terms, and more specifically in the context of the easy quantum groups introduced in~\cite{bsp}. In order to explain this, let us start with:

\emph{Conventions.} We write $P(k)$ and $NC(k)$ for the sets of partitions and noncrossing partitions of $\{1,\dots,k\}$, and
\[
B_k=|P(k)|,\qquad C_k=|NC(k)|=\frac{1}{k+1}\binom{2k}{k}
\]
for the Bell and Catalan numbers. Thus $B_0,B_1,B_2,\dots$ are 1, 1, 2, 5, 15, 52, 203, 877, 4140, 21147, \dots, and $C_0,C_1,C_2,\dots$ are 1, 1, 2, 5, 14, 42, 132, 429, 1430, 4862, \dots, the two sequences starting to differ at $k=4$, because of the basic crossing. Scalar products are linear in the first variable: $\langle x,y\rangle=\sum_ix_i\bar y_i$ on $\C^N$, $\langle a,b\rangle=\mathrm{tr}(ab^*)$ on matrix algebras, and $\langle a,b\rangle=h(b^*a)$ on $L^2(G)$, with $h$ the Haar state. Thus the matrix of an operator $T$ in an orthonormal basis is $T_{ij}=\langle Te_j,e_i\rangle$, and $\Proj(\xi)=\langle\cdot,\xi\rangle\xi/\|\xi\|^2$ is the orthogonal projection onto $\C\xi$. A dimension sequence is called ``Catalan'' through level $k$ when it agrees with $C_l$ for all $l\le k$. Statements resting on computer calculations are marked as such in their headers (\emph{exact}, \emph{computational}, \emph{numerical}), and Section~\ref{sec:certify} explains how the exact ones are certified.

\medskip

\emph{Representation theory.} For any $N\ge4$, the irreducible representations of $S_N^+$ follow the same fusion rules as those of $SO_3$~\cite{ba}. Their intertwiners are described by partitions, as follows. Let $P(k,l)$ be the set of partitions of $k$ upper legs and $l$ lower legs, and $NC(k,l)\subset P(k,l)$ the set of noncrossing ones. Each $\pi\in P(k,l)$ gives a linear map $T_\pi:(\C^N)^{\otimes k}\to(\C^N)^{\otimes l}$, defined by
\[
T_\pi(e_{i_1}\otimes\dots\otimes e_{i_k})=\sum_{j_1,\dots,j_l}\delta_\pi\binom{i_1\dots i_k}{j_1\dots j_l}\,e_{j_1}\otimes\dots\otimes e_{j_l},
\]
where $e_1,\dots,e_N$ is the standard basis of $\C^N$, and where $\delta_\pi=1$ if $\pi\le\ker\binom{i}{j}$, and $\delta_\pi=0$ otherwise. Here $\ker\binom{i}{j}$ is the partition of the $k+l$ legs obtained by joining two legs when they carry the same index, and $\pi\le\ker\binom{i}{j}$ means that each block of $\pi$ is contained in a block of $\ker\binom{i}{j}$, that is, the indices are constant on the blocks of $\pi$. With this notation, we have
\[
\mathrm{Hom}_{S_N^+}(u^{\otimes k},u^{\otimes l})=\mathrm{span}\big\{T_\pi:\pi\in NC(k,l)\big\},
\]
while for $S_N$ one gets the span of all the $T_\pi$ with $\pi\in P(k,l)$. By Woronowicz's Tannaka--Krein duality~\cite{wo2}, the intermediate quantum groups $S_N\subset G\subset S_N^+$ therefore correspond to the tensor categories $\mathcal C$ with
\[
\mathrm{span}(NC)\subset\mathcal C\subset\mathrm{span}(P).
\]
Here $\mathcal C$ is allowed to contain arbitrary \emph{linear combinations} of partitions, not just partitions themselves.

\emph{Easiness.} A closed subgroup $S_N\subset G\subset O_N^+$ is called \emph{easy} when its intertwiner spaces are spanned by partitions, in the sense that there are sets $D(k,l)\subset P(k,l)$ such that
\[
\mathrm{Hom}_G(u^{\otimes k},u^{\otimes l})=\mathrm{span}\big\{T_\pi:\pi\in D(k,l)\big\}.
\]
The axioms of tensor categories translate into the fact that $D$ is a \emph{category of partitions}: it is stable under horizontal and vertical concatenation and under upside-down turning, and it contains the identity and the semicircle~\cite{bsp}. The basic examples are $O_N\subset O_N^+$, coming from all pairings and noncrossing pairings, $S_N\subset S_N^+$, coming from all partitions and noncrossing partitions, and the hyperoctahedral and bistochastic versions $H_N\subset H_N^+$ and $B_N\subset B_N^+$, coming from partitions with blocks of even size, and from singletons and pairings. In each case the liberation operation simply removes the crossings. See~\cite{bcs,rw} for the orthogonal classification, and~\cite{tw} for the unitary case. The easy quantum groups $S_N\subset G\subset S_N^+$ correspond to the categories of partitions $NC\subset D\subset P$, and these are easy to classify: a crossing partition, together with $NC$, generates all of $P$, so $D\in\{NC,P\}$. Thus $S_N^+$ is the only easy quantum group strictly containing $S_N$ inside $S_N^+$~\cite{bcs}.

\subsection{The conjecture}

The classification of the quantum subgroups of $S_4^+$ in~\cite{bb} contains no quantum group strictly between $S_4$ and $S_4^+$. This led Banica and Bichon to the following statement.

\begin{conjecture}[maximality]\label{conj:max}
For every $N$, the inclusion $S_N\subset S_N^+$ is maximal: there is no compact quantum group $G$ with $S_N\subsetneq G\subsetneq S_N^+$.
\end{conjecture}

\emph{Status.} The conjecture holds trivially for $N\le3$, and it holds for $N=4$ by~\cite{bb}. The case $N=5$ was settled in~\cite{bah}, using Jones' theory of subfactors~\cite{jo1} and the classification of subfactors of index at most $5$ and exactly $5$~\cite{imps,jms}. For $N\ge6$ it is open.

\emph{Three approaches.} As explained in the introduction, there are three natural approaches to the conjecture. The first one is via subgroups: given a non-classical $G\subset S_N^+$, decide whether $\langle S_N,G\rangle=S_N^+$. This is the natural source of counterexamples, and it is the subject of Sections~\ref{sec:duals}--\ref{sec:hadamard}. The second one is via exotic quantum groups: develop the theory of a hypothetical $S_N\subsetneq G\subsetneq S_N^+$, with the hope of reaching a contradiction, as in Section~\ref{sec:exotic}. The third one is via partitions: work directly with the categories $\spn(NC)\subset\mathcal C\subset\spn(P)$, level by level, as in~\cite{bah,mc2}. This is our main approach, developed in Sections~\ref{sec:method}--\ref{sec:higher}.

\emph{Ramifications.} The same question makes sense in the orthogonal case. Inside Wang's free orthogonal group $O_N^+$~\cite{wa95}, the half-liberated orthogonal group $O_N^*$, defined by the relations $abc=cba$, fits into
\[
O_N\subset O_N^*\subset O_N^+ ,
\]
and $O_N\subset O_N^*$ is known to be maximal~\cite{bbcc}. The conjecture is that $O_N^*$ is the only intermediate object, i.e.\ that this chain is complete. There are unitary and projective versions as well. As in the permutation case, all of this lies beyond easiness. These questions are discussed in Sections~\ref{sec:orthproblem}--\ref{sec:related}.


%======================================================================
% Section: groupduals
%======================================================================

\section{Group duals and Fourier families}\label{sec:duals}

\subsection{Group duals in $S_N^+$, after Bichon}\label{sec:bichon}

\paragraph{Conventions.}
A closed subgroup $G\subset S_N^+$ acts on $C(X)=\C^N$ by $\alpha(e_j)=\sum_ie_i\otimes u_{ij}$, and its fixed vectors at level $k$ are
\[
\Fix_G(u^{\otimes k})=\{\xi:u^{\otimes k}(\xi\otimes1)=\xi\otimes1\},\qquad (u^{\otimes k})_{i_1\dots i_k,j_1\dots j_k}=u_{i_1j_1}\cdots u_{i_kj_k}.
\]
For $\pi\in P(k)$ we use the vector $\xi_\pi=T_\pi$, with $\pi$ viewed as an element of $P(0,k)$, having no upper legs. In practice, this means
\[
\xi_\pi=\sum_{i:\,\pi\le\ker i}e_{i_1}\otimes\dots\otimes e_{i_k},
\]
where $\ker i$ is the partition of $\{1,\dots,k\}$ obtained by joining two legs when they carry the same index, so that the sum is over the multi-indices $i$ which are constant on the blocks of $\pi$. We have $\Fix_{S_N}(u^{\otimes k})=\spn\{\xi_\pi:\pi\in P(k)\}$, and, for $N\ge4$, $\Fix_{S_N^+}(u^{\otimes k})=\spn\{\xi_\pi:\pi\in NC(k)\}$, of dimension $C_k$.

Given two closed subgroups $H,K\subset S_N^+$, we denote by $\langle H,K\rangle$ the smallest closed subgroup of $S_N^+$ containing both. In Tannakian terms, $G=\langle H,K\rangle$ corresponds to the intersection of the categories,
\[
\mathcal C_G=\mathcal C_H\cap\mathcal C_K,
\]
where $\mathcal C_G$ denotes the tensor category of intertwiners of $G$. Since $u$ is self-conjugate, the intertwiner spaces are determined by the fixed vectors, and so the quantity to compute, for a given $G\subset S_N^+$, is
\[
\Fix_{\langle S_N,G\rangle}(u^{\otimes k})=\Fix_G(u^{\otimes k})\cap\spn\{\xi_\pi:\pi\in P(k)\}.
\]
This space contains $\Fix_{S_N^+}(u^{\otimes k})$, and for $N\ge4$ we have $\langle S_N,G\rangle=S_N^+$ if and only if its dimension is $C_k$ for every $k$.

For a group dual, $C(\widehat\Gamma)=C^*(\Gamma)$ with $\Delta(g)=g\otimes g$. An embedding $\widehat\Gamma\subset S_N^+$ is given by a magic unitary $u$ with coefficients $u_{ij}\in C^*(\Gamma)$, satisfying $\Delta(u_{ij})=\sum_ku_{ik}\otimes u_{kj}$. We always assume that the coefficients $u_{ij}$ generate $C^*(\Gamma)$, which means that the morphism $C(S_N^+)\to C^*(\Gamma)$ is surjective, so that $\widehat\Gamma$ is a closed quantum subgroup of $S_N^+$. This is no loss of generality: in general, the $u_{ij}$ generate $C^*(\Gamma_0)$ for a subgroup $\Gamma_0\subset\Gamma$, and we can replace $\Gamma$ by $\Gamma_0$. The quantum subgroups of $\widehat\Gamma$ are the $\widehat\Lambda$ for quotients $\Gamma\to\Lambda$, and $\widehat\Gamma$ is classical iff $\Gamma$ is abelian.

\paragraph{Bichon's theorem.}
A coaction of $\C[\Gamma]$ on $\C^N$ is the same thing as a $\Gamma$-grading $\C^N=\bigoplus_gA_g$, where $A_g=\{a:\alpha(a)=a\otimes g\}$. In this language, Bichon proved the following~\cite[Prop.~5.2, Cor.~5.3]{bi}:
\begin{itemize}
\item an ergodic grading of $\C^m$ comes from a transitive abelian permutation group;
\item consequently, $\widehat\Gamma\subset S_N^+$ iff $\Gamma$ is a quotient of $G_1*\dots*G_k$, where the $G_i\subset S_{m_i}$ are transitive abelian groups and $\sum m_i=N$.
\end{itemize}
Since transitive abelian groups act regularly, $G_i$ can be any finite abelian group, with $m_i=|G_i|$.

We restate Bichon's result in the form used below, with explicit matrices.

\begin{definition}[Fourier families]\label{def:family}
A \emph{Fourier family} for $\Gamma$ is a family $\mathcal H=(H_1,\dots,H_k)$ of finite abelian subgroups of $\Gamma$ that generate $\Gamma$. Repetitions and trivial subgroups are allowed. Its size is $N=\sum_i|H_i|$. The family acts on the points
\[
X_{\mathcal H}=\widehat H_1\sqcup\dots\sqcup\widehat H_k
\]
through the magic unitary $u^{\mathcal H}=\bigoplus_iu^{(i)}$, where
\[
u^{(i)}_{\chi\psi}=P^{(i)}_{\chi\bar\psi},\qquad P^{(i)}_\rho=\frac1{|H_i|}\sum_{h\in H_i}\rho(h)\,h\qquad(\chi,\psi,\rho\in\widehat H_i).
\]
Equivalently, $u^{(i)}=F_i\,\mathrm{diag}(h)_{h\in H_i}\,F_i^*$, where $F_i$ is the normalized Fourier matrix of $H_i$, given by $(F_i)_{\chi h}=\chi(h)/\sqrt{|H_i|}$.
\end{definition}

\begin{theorem}[Bichon, restated]\label{thm:bichon}
For a Fourier family $\mathcal H$, the matrix $u^{\mathcal H}$ defines a closed quantum subgroup $\widehat\Gamma\subset S_N^+$. Conversely, every closed quantum subgroup $\widehat\Gamma\subset S_N^+$ is conjugate by a permutation to $u^{\mathcal H}$, for a unique unordered family $\mathcal H$. Here $H_i$ is the support of the grading restricted to the $i$-th orbit.
\end{theorem}

\begin{proof}[Comments]
The direct statement is a short computation. The $P_\rho$ are the minimal projections of $\C[H_i]$, rows and columns of $u^{(i)}$ are permutations of $\{P_\rho\}_\rho$, and $\Delta(u_{\chi\psi})=\sum_\eta u_{\chi\eta}\otimes u_{\eta\psi}$ follows from orthogonality of characters.

The converse is in~\cite{bi}. The only point to add is the explicit identification. On an orbit $O$ with grading support $H$, each $A_h\cap C(O)$ is spanned by a nowhere-vanishing function $v_h$. Normalizing $v_h(x_0)=1$ at a base point gives $v_hv_k=v_{hk}$. Then $x\mapsto(h\mapsto v_h(x))$ is a bijection $O\simeq\widehat H$ which turns the action into the Fourier block. The $H_i$ are intrinsic, being the supports of the grading, so they are invariants under conjugation by $S_N$. A change of base point translates $\widehat H_i$.
\end{proof}

\begin{remark}
Up to relabeling the points, \emph{actions of $\widehat\Gamma$ on $N$ points are exactly Fourier families of size $N$}. This is the right parametrization for our purposes, since $\langle S_N,G\rangle$ is invariant under conjugation by $S_N$.
\end{remark}

\begin{example}\label{ex:blocks}
\begin{itemize}
\item $H=\Z_2=\langle g\rangle$: $u=M(p)=\begin{pmatrix}p&1-p\\1-p&p\end{pmatrix}$ with $p=\frac{1+g}2$.
\item $H=\Z_n=\langle g\rangle$: $u$ is the circulant matrix built on the spectral projections of $g$. This is the cyclic case.
\item $H=\{e,a,b,c\}\simeq\Z_2^2$, the Klein block. Writing $P_x$ for the projection of the character with kernel $\{e,x\}$ (and $P_e=\frac14(e+a+b+c)$), we get
\[
u=\begin{pmatrix}P_e&P_a&P_b&P_c\\P_a&P_e&P_c&P_b\\P_b&P_c&P_e&P_a\\P_c&P_b&P_a&P_e\end{pmatrix}.
\]
In general, the pattern of $u^{(i)}$ is the division table $(\chi,\psi)\mapsto\chi\psi^{-1}$ of $\widehat H_i$, which is circulant for $\Z_n$, and is not circulant here.
\end{itemize}
\end{example}

\paragraph{Notation used throughout.}
\begin{itemize}
\item \emph{Fourier basis.} $f_{(i,h)}\in C(X_{\mathcal H})$, for $h\in H_i$, is the function $\chi\mapsto\chi(h)$ on $\widehat H_i$, extended by $0$. These are orthogonal, $\|f_{(i,h)}\|^2=|H_i|$, and $u(f_{(i,h)}\otimes1)=f_{(i,h)}\otimes h$.
\item \emph{Labelings.} A labeling of $k$ legs is $a=((i_1,h_1),\dots,(i_k,h_k))$ with $h_l\in H_{i_l}$. Write $f_a=\bigotimes_lf_{(i_l,h_l)}$, and $w(a)=h_1\cdots h_k\in\Gamma$ for its \emph{word} (ordered product).
\item \emph{Standard tori.} For finite abelian groups $\mathcal H=(H_1,\dots,H_k)$, the standard torus is $T_{\mathcal H}=\widehat{H_1*\dots*H_k}$ with its tautological family.
\end{itemize}

\begin{lemma}[quotients]\label{lem:quotient}
Let $q:\Gamma\to\Lambda$ be a quotient. Then $\widehat\Lambda\subset\widehat\Gamma$ acts on the same points, and its family is obtained as follows. Each orbit $\widehat H_i$ splits into $|\ker(q_{|H_i})|$ orbits, namely the cosets of $q(H_i)^\wedge\subset\widehat H_i$. Each of the new orbits carries the subgroup $q(H_i)$.
\end{lemma}

\begin{proof}
The $\Lambda$-grading is $f_{(i,h)}\mapsto q(h)$. Its identity component on $\widehat H_i$ consists of the functions constant on cosets of $\ker(q_{|H_i})^\perp=q(H_i)^\wedge$.
\end{proof}

\begin{example}
Sending $\Z_4\to\Z_2$ in $\Z_4*\Z_2$ turns the $4$-point orbit into two $2$-point orbits with the same subgroup, i.e.\ $u=\mathrm{diag}(M(p),M(p),M(q))$.
\end{example}

%=====================================================================
\subsection{Reduction to just-non-abelian groups}\label{sec:jna}

We first record two simple principles, valid for arbitrary closed subgroups $G\subset S_N^+$. We say that $G$ \emph{moves at most $M$ points} if, up to a permutation of the coordinates, its magic unitary has the form $u=v\oplus1_{N-M}$, with $v$ of size $M$.

\begin{proposition}[monotonicity and locality]\label{prop:monloc}
Let $G\subset S_N^+$ be a closed subgroup, with $N\ge4$.
\begin{enumerate}
\item If $K\subset G$ is a closed subgroup, then $\langle S_N,K\rangle\subset\langle S_N,G\rangle$. In particular, if $\langle S_N,K\rangle=S_N^+$, then $\langle S_N,G\rangle=S_N^+$.
\item If $G$ has a non-classical closed subgroup $K$ which moves at most $5$ points, then $\langle S_N,G\rangle=S_N^+$.
\end{enumerate}
\end{proposition}

\begin{proof}
(1) This is clear from the definition of $\langle\,\cdot\,,\cdot\,\rangle$.

(2) Since $\langle S_N,\sigma G\sigma^{-1}\rangle=\sigma\langle S_N,G\rangle\sigma^{-1}$ for $\sigma\in S_N$, we can assume that $K$ moves only the first $M\le5$ points, so that $K\subset S_M^+$, embedded in $S_N^+$ by fixing the last $N-M$ points. As $K$ is not classical and $S_M^+=S_M$ for $M\le3$, we have $M\in\{4,5\}$. The maximality of $S_M\subset S_M^+$ for $M=4,5$, from~\cite{bah,bb}, gives $\langle S_M,K\rangle=S_M^+$. Since $S_M\subset S_N$, we obtain $S_M^+\subset\langle S_N,G\rangle$, and so $\langle S_N,G\rangle\supset\langle S_N,S_M^+\rangle=S_N^+$, by~\cite[Cor.~3.4]{bcf} for $N\ge5$, and trivially for $N=4$.
\end{proof}

For a group dual $\widehat\Gamma\subset S_N^+$, given by a Fourier family $\mathcal H=(H_1,\dots,H_k)$, the closed quantum subgroups are the duals $\widehat\Lambda$ of the quotients $q:\Gamma\to\Lambda$. By Lemma~\ref{lem:quotient}, $\widehat\Lambda$ moves exactly the points of the orbits $\widehat H_i$ with $q(H_i)\ne1$, so their number is $\sum_{q(H_i)\ne1}|H_i|$. Thus, by monotonicity, we are led to look for small quotients of $\Gamma$ which are still non-abelian. This suggests the following notion.

\begin{definition}
A group $\Gamma$ is called \emph{just-non-abelian} (JNA) if it is non-abelian, and all its proper quotients are abelian. Equivalently, since $\Gamma/K$ is abelian if and only if $K\supset\Gamma'$, where $\Gamma'=[\Gamma,\Gamma]$ is the derived subgroup, $\Gamma$ is JNA when $\Gamma'\ne1$, and every nontrivial normal subgroup of $\Gamma$ contains $\Gamma'$.
\end{definition}

\begin{lemma}\label{lem:zorn}
Every non-abelian finitely generated group has a JNA quotient.
\end{lemma}

\begin{proof}
Let $\Gamma=\langle g_1,\dots,g_n\rangle$. A normal subgroup $K\subset\Gamma$ gives a non-abelian quotient $\Gamma/K$ if and only if $[g_a,g_b]\notin K$ for some pair $(a,b)$. The set of such normal subgroups contains $\{1\}$, and is inductive. Indeed, given a chain $(K_t)$ of such subgroups, since there are finitely many pairs, there is a pair $(a,b)$ with $[g_a,g_b]\notin K_t$ for cofinally many $t$, and hence for all $t$, so that $[g_a,g_b]\notin\bigcup_tK_t$. By Zorn's lemma we obtain a maximal such $K$, and then $\Gamma/K$ is JNA.
\end{proof}

\begin{proposition}\label{prop:reduction}
Let $N\ge4$. If $\langle S_N,\widehat\Lambda\rangle=S_N^+$ for all JNA groups $\Lambda$ and all their Fourier families of size $N$, then $\langle S_N,\widehat\Gamma\rangle=S_N^+$ for all non-classical group duals $\widehat\Gamma\subset S_N^+$. Moreover, when $\Gamma$ is residually finite, for instance a free product of finite abelian groups, as for the standard tori $T_{\mathcal H}$, the JNA quotient $\Lambda$ used for $\widehat\Gamma$ can be chosen finite.
\end{proposition}

\begin{proof}
Let $\widehat\Gamma\subset S_N^+$ be non-classical, given by a Fourier family $\mathcal H=(H_1,\dots,H_k)$. Then $\Gamma$ is non-abelian, and finitely generated, by the finite subgroups $H_i$. By Lemma~\ref{lem:zorn}, it has a JNA quotient $q:\Gamma\to\Lambda$. The quantum group $\widehat\Lambda\subset\widehat\Gamma$ acts on the same $N$ points, through the family $(q(H_1),\dots,q(H_k))$ split as in Lemma~\ref{lem:quotient}, which generates $\Lambda$. By assumption $\langle S_N,\widehat\Lambda\rangle=S_N^+$, and so $\langle S_N,\widehat\Gamma\rangle=S_N^+$, by monotonicity. If $\Gamma$ is residually finite, pick generators $g_a,g_b$ with $[g_a,g_b]\ne1$, and a normal subgroup $K_0$ of finite index not containing $[g_a,g_b]$. Then $\Gamma/K_0$ is finite and non-abelian, and applying Lemma~\ref{lem:zorn} to it gives a finite JNA quotient of $\Gamma$. Finally, free products of finite groups are residually finite.
\end{proof}

The finite JNA groups are well known, and can be described as follows.

\begin{proposition}\label{prop:jna}
A finite group $\Gamma$ is JNA if and only if $\Gamma'\ne1$ is its unique minimal normal subgroup. This happens if and only if $\Gamma$ belongs to one of the following three classes, which are mutually exclusive:
\begin{enumerate}[label=(\alph*)]
\item \emph{Central type:} $\Gamma$ is a $p$-group with $|\Gamma'|=p$ and cyclic center.
\item \emph{Affine type:} $\Gamma=\F_p^m\rtimes A$, with $A\subset\F_{p^m}^\times$ a nontrivial subgroup not contained in a proper subfield, acting by multiplication.
\item \emph{Non-abelian type:} $\Gamma'\simeq T^m$ with $T$ non-abelian simple, $\Gamma'$ is a minimal normal subgroup of $\Gamma$, and its centralizer is trivial, $C_\Gamma(\Gamma')=1$.
\end{enumerate}
\end{proposition}

\begin{proof}
If $\Gamma$ is JNA, every minimal normal subgroup contains $\Gamma'\ne1$, hence equals it. Conversely, if $\Gamma'\ne1$ is the unique minimal normal subgroup, every nontrivial normal subgroup contains it. As for the classification, the solvable case goes back to Newman~\cite{ne1,ne2}, and the general case, including the non-solvable one, is treated by Jain and Shukla~\cite{js}, who also review Newman's results. In any case the argument is short. $V=\Gamma'$ is a minimal normal subgroup, so $V\simeq T^m$ with $T$ simple. If $V$ is central, all subgroups of $Z(\Gamma)$ contain $V$, and $\Gamma$ is nilpotent, which leads to (a). If $V=\F_p^m$ is not central, the abelian group $\Gamma/C_\Gamma(V)$ acts faithfully and irreducibly on $V$, so it embeds in $\F_{p^m}^\times$ by Schur's lemma, and a coprime action argument gives $C_\Gamma(V)=V$, which leads to (b). If $V$ is non-abelian, $C_\Gamma(V)$ is normal and meets $V$ trivially, so it is trivial, which is (c). The converse is a direct check.
\end{proof}

\begin{example}\label{ex:jna}
Here are some examples of finite JNA groups.
\begin{enumerate}[label=(\alph*)]
\item The dihedral group $D_4$ of order $8$, the quaternion group $Q_8$, the extraspecial groups $p^{1+2n}_\pm$, and the Pauli groups $\Z_4\circ2^{1+2n}$.
\item The groups $S_3$, $A_4=\F_2^2\rtimes\Z_3$, the dihedral groups $D_p=\F_p\rtimes\Z_2$ with $p$ odd prime, and $\F_7\rtimes\Z_3$, $\F_5\rtimes\Z_4$.
\item The non-abelian simple groups, the symmetric groups $S_n$ with $n\ge5$, the groups $PGL_2(q)$ with $q\ge4$, and the wreath products $T\wr\Z_m$, with $T$ non-abelian simple.
\end{enumerate}
In case (c), the minimality of $\Gamma'$ cannot be dropped. For instance $\Gamma=A_5\times A_5$ satisfies $\Gamma'\simeq A_5^2$ and $C_\Gamma(\Gamma')=1$, but it is not JNA, having $A_5$ as a quotient.
\end{example}

For group duals, being JNA has the following consequence. Given a Fourier family $\mathcal H=(H_1,\dots,H_k)$, we denote by $n(\mathcal H)=\sum_{H_i\ne1}|H_i|$ the number of points moved by the corresponding action.

\begin{corollary}[locality for JNA groups]\label{cor:jnaloc}
Let $\Gamma$ be JNA, and $\widehat\Gamma\subset S_N^+$ be given by a Fourier family $\mathcal H$. Then $\widehat\Gamma$ is its only non-classical closed quantum subgroup. Thus Proposition~\ref{prop:monloc}~(2) applies to $\widehat\Gamma$ if and only if $n(\mathcal H)\le5$.
\end{corollary}

\begin{proof}
The closed quantum subgroups of $\widehat\Gamma$ are the duals of the quotients of $\Gamma$, and the proper quotients are abelian, with classical duals. The second assertion follows, the number of points moved by $\widehat\Gamma$ being $n(\mathcal H)$.
\end{proof}

Observe that $n(\mathcal H)$ counts all the nontrivial blocks of $\mathcal H$, even when some of them already generate $\Gamma$: the family $(\langle(12)\rangle,\langle(23)\rangle,\langle(13)\rangle)$ of $S_3$ has $n(\mathcal H)=6$, and is not covered by locality. So what remains to be tested are the JNA groups with Fourier families having $n(\mathcal H)\ge6$, and all three classes occur below. In class (a) we will meet $D_4$, $Q_8$, $3^{1+2}_+$, and the magic square group $2^{1+4}_+$ of Proposition~\ref{prop:magic}. In class (b) we will meet $S_3$, $A_4$, $D_5$ and $\F_7\rtimes\Z_3$, see Example~\ref{ex:table}. In class (c) we will meet $A_7$, as the JNA quotient of the solution group $3.A_7$ of~\cite{drs}, acting on $210$ points, see Proposition~\ref{prop:a7}. Class (c) has much smaller actions as well, for instance $A_5$ acting on $6$ points through the family $(\langle(123)\rangle,\langle(345)\rangle)$, see Proposition~\ref{prop:classc}.

%=====================================================================
\subsection{The laws lemma for abelian blocks}\label{sec:laws}

Fix a Fourier family $\mathcal H$ for $\Gamma$, and use the notation of Section~\ref{sec:bichon}.

\begin{proposition}\label{prop:coords}
\begin{enumerate}
\item $\Fix_{\widehat\Gamma}(u^{\otimes k})=\spn\{f_a:w(a)=e\}$.
\item Write $\xi_\pi=\sum_ac_\pi(a)f_a$. Then $c_\pi(a)\ne0$ iff each block $b$ of $\pi$ carries a single orbit index $i_b$ and satisfies $\prod_{l\in b}h_l=e$ in $H_{i_b}$. In that case $c_\pi(a)=\prod_{b\in\pi}|H_{i_b}|^{1-|b|}$.
\end{enumerate}
\end{proposition}

\begin{proof}
(1) $u^{\otimes k}(f_a\otimes1)=f_a\otimes w(a)$, and distinct group elements are linearly independent.

(2) Use $e_{(i,\chi)}=|H_i|^{-1}\sum_h\overline{\chi(h)}f_{(i,h)}$ in $\sum_xe_x^{\otimes b}$, block by block, and then $\sum_\chi\chi(h)=|H_i|\delta_{h,e}$.
\end{proof}

\begin{definition}
Two labelings are \emph{equivalent} if they differ by a bijection $\sigma$ of orbit indices together with isomorphisms $\theta_i:H_i\to H_{\sigma(i)}$ applied to the labels. A \emph{pattern} is an equivalence class. Concretely, a pattern consists of:
\begin{itemize}
\item the partition $\kappa$ of the legs by orbit;
\item for each block of $\kappa$, an abelian isomorphism type with labels in it, up to automorphisms.
\end{itemize}
A pattern is \emph{balanced} if the labels multiply to $e$ in each block of $\kappa$. It is \emph{active} if some realization has $w(a)\ne e$. A balanced inactive pattern is a \emph{law}: an identity holding for all choices of distinct orbit subgroups of the given types, and all identifications. A pattern is \emph{freely trivial} if its words are trivial in the free product $H_1*\dots*H_k$, for all its realizations. For instance, on a single orbit, the labels $(h,h^{-1})$ give a freely trivial pattern. Freely trivial patterns are laws for every $\Gamma$, and we call the other laws \emph{nontrivial}.
\end{definition}

\begin{lemma}[laws lemma]\label{lem:laws}
For $\xi\in\spn\{\xi_\pi\}$:
\begin{enumerate}
\item $c_\xi$ is constant on patterns and vanishes off balanced ones;
\item $\xi\in\Fix_{\langle S_N,\widehat\Gamma\rangle}$ iff $c_\xi$ vanishes on all active patterns.
\end{enumerate}
Hence the joint intertwiner spaces depend only on the multiset of types of the $H_i$ and on the set of active patterns.
\end{lemma}

\begin{proof}
The vectors $\xi_\pi$ are invariant under $S_N$. Three kinds of permutations act diagonally, up to phases, on the Fourier basis:
\begin{itemize}
\item translations of $\widehat H_i$, which multiply $f_{(i,h)}$ by $\overline{\psi(h)}$; these give balancedness;
\item automorphisms of $H_i$;
\item swaps of isomorphic orbits.
\end{itemize}
The last two generate the equivalence of labelings. Then apply Proposition~\ref{prop:coords}(1).
\end{proof}

\begin{corollary}\label{cor:laws}
If $(\Gamma,\mathcal H)$ has no nontrivial laws with $\le k$ legs, then up to level $k$ the joint quantum group $\langle S_N,\widehat\Gamma\rangle$ has the same intertwiners as $\langle S_N,T_{\mathcal H}\rangle$. Indeed, in the torus $T_{\mathcal H}$ the laws are exactly the freely trivial patterns.
\end{corollary}

\begin{remark}\label{rem:laws}
Thus, for group duals, the question becomes that of knowing \emph{which short balanced patterns are nontrivial laws}. If there are no nontrivial laws with $\le k$ legs, then $\widehat\Gamma$ behaves up to level $k$ like the standard torus, which is free in the relevant sense. So a counterexample to the conjecture should be looked for among JNA groups having many short nontrivial laws.
\end{remark}

%=====================================================================
\subsection{First candidates}\label{sec:examples}

\begin{example}\label{ex:table}
The following are JNA groups with Fourier families of size $N\ge6$, so none of them is killed by locality. The list is illustrative, not exhaustive.
\begin{center}
\begin{tabular}{@{}llll@{}}\toprule
$\Gamma$ & class & family & $N$\\\midrule
$S_3$ & (b) & $\langle(12)\rangle,\langle(23)\rangle,\langle(13)\rangle$ & 6\\
$D_4$ & (a) & $\langle s\rangle,\ \langle r\rangle\simeq\Z_4$ & 6\\
$D_4$ & (a) & $\langle sr\rangle,\ \{1,r^2,s,sr^2\}$ (Klein) & 6\\
$A_4$ & (b) & $\langle(123)\rangle,\langle(234)\rangle$ & 6\\
$3^{1+2}_+$ & (a) & $\langle x\rangle,\langle y\rangle\simeq\Z_3$ & 6\\
$\F_7\rtimes\Z_3$ & (b) & two complements $\Z_3$ & 6\\
$D_5$ & (b) & $\langle s\rangle,\langle sr\rangle,\langle sr^2\rangle$ & 6\\
$A_4$ & (b) & $\langle(123)\rangle,\ V_4$ (Klein) & 7\\
$D_5$ & (b) & $\langle s\rangle,\ \langle r\rangle\simeq\Z_5$ & 7\\
$Q_8$ & (a) & $\langle i\rangle,\langle j\rangle\simeq\Z_4$ & 8\\
$D_4$ & (a) & $\{1,r^2,s,sr^2\},\ \{1,r^2,sr,sr^3\}$ (two Klein) & 8\\
$2^{1+4}_+$ & (a) & nine $\Z_2$'s (magic square) & 18\\\bottomrule
\end{tabular}
\end{center}
\end{example}

\subsection{Computations}

\begin{remark}[method]
The program \texttt{duals.py} works as follows.
\begin{itemize}
\item Input: a finite group $\Gamma$, given by permutations or matrices, together with a Fourier family.
\item It enumerates the labelings that are balanced on each orbit and have nontrivial word.
\item For each of them it computes the row $(c_\pi(a))_{\pi\in P(k)}$ given by Proposition~\ref{prop:coords}.
\item It returns
\[
\dim\Fix_{\langle S_N,\widehat\Gamma\rangle}(k)=\dim\spn\{\xi_\pi\}-\mathrm{rank}\,(c_\pi(a))_{a\ \mathrm{bad},\,\pi},
\]
with the rank computed exactly modulo the prime $p=1048573$. Only the crossing columns are kept, the noncrossing ones vanishing identically, and the rank is computed after a random compression $M=RC$, which can only lower it.
\end{itemize}
\end{remark}

\begin{remark}[validation]
This computation was checked in several ways.
\begin{itemize}
\item It agrees with an independent dense computation, using the Haar projection in the regular representation, on four small cases at $k\le4$. These include a Klein block for $D_4$ at $N=6$ and for $A_4$ at $N=7$.
\item It gives Bell numbers $2,5,15$ for the classical Klein group acting regularly on $4$ points, as it should.
\item It reproduces $14,42$ for the magic square at $k=4,5$.
\end{itemize}
\end{remark}

\begin{proposition}[exact]\label{prop:dualcomp}
For the $11$ actions of Example~\ref{ex:table} other than the magic square, and for $k=4,5,6,7,8$,
\[
\dim\Fix_{\langle S_N,\widehat\Gamma\rangle}(u^{\otimes k})=14,\ 42,\ 132,\ 429,\ 1430,
\]
the Catalan numbers. So none of them is a counterexample through level $8$. This covers classes (a) and (b), cyclic and Klein blocks, and families with three blocks on $6$ points.
\end{proposition}

Each run takes about a second for $k\le7$, and between $20$ and $40$ seconds for $k=8$. Although the ranks are computed modulo $p$, after a random compression, the result is exact. Indeed, the coefficients $c_\pi(a)$ are rational, with denominators products of the orders $|H_i|$, which are prime to $p$. Thus the rank modulo $p$, and a fortiori that of the compressed matrix, is at most the rank over $\Q$, so the computed dimension is an upper bound for the true one. On the other hand, the true dimension is at least $C_k$, because $\Fix_{S_N^+}(u^{\otimes k})$ is contained in it. So a Catalan output proves equality. Here $\dim\spn\{\xi_\pi\}$ must be taken exactly; for $k>N$ it is the number of partitions of $\{1,\dots,k\}$ with at most $N$ blocks.

\subsection{Conclusions}\label{sec:dualsconcl}

As a conclusion, in this section we used the group duals $\widehat\Gamma\subset S_N^+$ as seeds, in the search for a counterexample to the maximality conjecture of the form $S_N\subsetneq\langle S_N,\widehat\Gamma\rangle\subsetneq S_N^+$. Here is what we found.

\emph{Structure.} By Bichon's theorem, the group duals in $S_N^+$ are exactly the Fourier families of finite abelian subgroups generating $\Gamma$. Non-cyclic orbits do occur, the simplest being the Klein block.

\emph{Reductions.} By monotonicity and locality, it is enough to consider JNA groups, and their Fourier families moving at least $6$ points. These groups are well known, and come in three classes: central type (a), affine type (b), and non-abelian type (c). Moreover, by the laws lemma, the joint intertwiners only depend on which short balanced patterns are nontrivial laws, and a group dual with no nontrivial laws up to $k$ legs behaves up to level $k$ like the corresponding free torus.

\emph{Computations.} The $11$ actions of Example~\ref{ex:table} from classes (a) and (b), with $N=6,7,8$ and including non-cyclic orbits, all give the Catalan numbers through level $8$. Thus, in all the cases we tested, $\widehat\Gamma$ is absorbed by $S_N$, as far as we can see, and no counterexample appears. Further group duals, coming from graphs, will appear in Section~\ref{sec:graphs}, with the same outcome.

\emph{Outlook.} All this suggests the following statement, which would be of independent interest, independently of the maximality conjecture.

\begin{conjecture}\label{conj:duals}
For every non-classical group dual $\widehat\Gamma\subset S_N^+$ we have $\langle S_N,\widehat\Gamma\rangle=S_N^+$.
\end{conjecture}

By Proposition~\ref{prop:reduction}, it is enough to prove this for the JNA groups $\Gamma$, and their Fourier families with $n(\mathcal H)\ge6$. The evidence consists of the $11$ actions of Example~\ref{ex:table}, of the magic square group (class (a), $N=18$), of the four class (c) actions at $N=6$ of Proposition~\ref{prop:classc} below, and of $A_7$ (class (c), $N=210$, no laws of length $\le8$), all Catalan through level $8$. The natural tool is the laws lemma, which reduces everything to the question of which short balanced patterns are laws. Class by class the laws are known, being central of order $p$ for (a), and metabelian for (b). So a strategy would be to find, for classes (a) and (b) and at each level $k$, enough active patterns to kill all the crossing combinations, using the explicit coefficients $c_\pi(a)=\prod_b|H_{i_b}|^{1-|b|}$ of Proposition~\ref{prop:coords}. A first target is the case where all the blocks are $\Z_2$.

On the computational side, the smallest actions of class (c) occur at $N=6$, and we have tested them as well.

\begin{proposition}[exact]\label{prop:classc}
For the following actions of class (c), with $N=6$, and for $k=4,5,6,7,8$, the dimensions $\dim\Fix_{\langle S_N,\widehat\Gamma\rangle}(u^{\otimes k})$ are the Catalan numbers $14,42,132,429,1430$:
\begin{itemize}
\item $A_5$ with $(\langle(123)\rangle,\langle(345)\rangle)\simeq(\Z_3,\Z_3)$;
\item $A_5$ with three $\Z_2$'s, generated by $(23)(45)$, $(14)(25)$, $(12)(45)$;
\item $A_5$ with $(\Z_2,\Z_2^2)$, given by $\langle(15)(23)\rangle$ and the Klein group on $\{1,2,3,4\}$;
\item $S_5$ with $(\langle(1234)\rangle,\langle(45)\rangle)\simeq(\Z_4,\Z_2)$.
\end{itemize}
\end{proposition}

This was computed with \texttt{duals.py}, as in Proposition~\ref{prop:dualcomp}, in about $20$ seconds per action at $k=8$. Class (c) groups have the fewest short laws, so by the laws lemma they should be the ``most free'', and this is what we find. As a continuation, it would be interesting to test $A_5$ at $N=7,8$, for instance with $(\Z_2,\Z_5)$, then $PSL_2(7)$, and so on. Note that $A_5$, being $(2,3)$-generated, also acts on $2+3=5$ points, but this action is absorbed by maximality at $N=5$, and says nothing about the actions on $6$ points. The computation is cheap, of the same size as those of Proposition~\ref{prop:dualcomp}.

Finally, Proposition~\ref{prop:reduction} reduces the problem to the finite JNA groups only when $\Gamma$ is residually finite. This leads to the following question.

\begin{question}\label{qu:infjna}
Does every non-abelian group generated by finitely many finite abelian subgroups have a finite non-abelian quotient? If not, which infinite JNA groups occur as group duals in $S_N^+$, with $n(\mathcal H)\ge6$?
\end{question}

The standard tori $T_{\mathcal H}$ are free products of finite groups, hence residually finite, but their quotients need not be. An infinite simple group generated by finitely many finite abelian subgroups would be JNA, with no finite JNA quotient, and if such groups exist with only two or three small blocks, then their duals are test cases that no finite computation of the above type sees. Note, however, that an infinite simple quotient of $\Z_2*\Z_3$ would act on $5$ points only, and be absorbed. So the question really concerns the case $n(\mathcal H)\ge6$, and is mostly one of group theory.


%======================================================================
% Section: graphs
%======================================================================

\section{Quantum automorphisms of graphs}\label{sec:graphs}

Recall that we are looking for seeds: non-classical quantum subgroups $K\subset S_N^+$, as small as possible, such that $\langle S_N,K\rangle$ might fail to be $S_N^+$. The group duals were the first natural family. The second one consists of the quantum automorphism groups of graphs, and of their quantum subgroups.

\subsection{Quantum automorphism groups of graphs}

Let $X$ be a finite graph with vertex set $\{1,\dots,N\}$ and adjacency matrix $A=A_X$. Its quantum automorphism group $G^+(X)\subset S_N^+$ is given by the relations $uA=Au$ on the fundamental magic unitary. Equivalently, it is the largest quantum subgroup of $S_N^+$ whose coaction on $C(X)=\C^N$ commutes with $A$, viewed as an operator on $C(X)$. Its classical version is the automorphism group $G(X)=G^+(X)\cap S_N$. The graph is said to have \emph{quantum symmetry} when $G^+(X)\ne G(X)$.

For many examples and computations of $G^+(X)$ see~\cite{sz,sch2} and the references there. For each such graph, $\langle S_N,G^+(X)\rangle$ is a candidate for an exotic intermediate quantum group. Monotonicity and locality (Proposition~\ref{prop:monloc}) reduce the question to the quantum subgroups of $G^+(X)$. The natural ones are group duals.

\subsection{Disjoint automorphisms}

Schmidt's criterion~\cite{sch} says that a graph with two non-trivial automorphisms of disjoint supports has quantum symmetry. In the language of Section~\ref{sec:duals}, this becomes an explicit group dual inside $G^+(X)$.

\begin{proposition}\label{prop:disjoint}
Let $C_1,\dots,C_r\subset G(X)$ be abelian groups of automorphisms whose supports $\supp C_i$ are pairwise disjoint. Let
\[
\Gamma=\widehat{C_1}*\dots*\widehat{C_r}.
\]
Then $\widehat\Gamma\subset G^+(X)$, acting by the Fourier family with the following blocks:
\begin{itemize}
\item one block for each $C_i$-orbit $O\subset\supp C_i$, with subgroup $H_O=\mathrm{Stab}_{C_i}(O)^\perp\subset\widehat{C_i}$;
\item trivial blocks at the points outside $\bigcup_i\supp C_i$.
\end{itemize}
This group dual is non-classical as soon as two of the $C_i$ are nontrivial.
\end{proposition}

\begin{proof}
On $C(\supp C_i)$, consider the $\widehat{C_i}$-grading coming from the classical action of $C_i$: the component of degree $\chi$ consists of the functions $f$ with $f\circ c^{-1}=\chi(c)f$ for all $c\in C_i$. Extend this by the trivial grading on the functions supported outside $\bigcup_i\supp C_i$. Since the supports are disjoint, the result is a $\Gamma$-grading of $C(X)$:
\begin{itemize}
\item products of functions with different supports vanish;
\item the unit is a sum of degree-$e$ indicator functions.
\end{itemize}
So by Theorem~\ref{thm:bichon} it is a coaction of $\widehat\Gamma$, with the stated orbits. Recall that an abelian group acts on each of its orbits regularly, through $C_i/\mathrm{Stab}(O)$.

It remains to check that $A$ preserves every graded component.
\begin{itemize}
\item For $\chi\ne1$ in $\widehat{C_i}$, the component of degree $\chi$ is the $\chi$-eigenspace of $C_i$ acting on all of $C(X)$. Indeed, an eigenfunction for a nontrivial character vanishes at the fixed points of $C_i$. This eigenspace is preserved by $A$, since $A$ commutes with $C_i$.
\item The component of degree $e$ is the space of functions invariant under all the $C_i$, which is also preserved by $A$.
\item All other components are zero.
\end{itemize}
Hence the coaction commutes with $A$, and $\widehat\Gamma\subset G^+(X)$. If two of the $C_i$ are nontrivial, then $\Gamma$ is non-abelian.
\end{proof}

\begin{remark}
Some comments on this result.
\begin{itemize}
\item With $r=2$ and $C_1=\langle\sigma\rangle$, $C_2=\langle\tau\rangle$, this recovers Schmidt's criterion, and identifies the quantum symmetry that the criterion detects.
\item Replacing $\sigma,\tau$ by suitable powers, we may assume that they have prime orders $\ell,\ell'$. Then $\Gamma=\Z_\ell*\Z_{\ell'}$, acting through one $\Z_\ell$-block for each $\ell$-cycle of $\sigma$ and one $\Z_{\ell'}$-block for each $\ell'$-cycle of $\tau$. A smaller support only helps.
\end{itemize}
\end{remark}

\begin{corollary}\label{cor:small}
If $X$ has two non-trivial automorphisms $\sigma,\tau$ with disjoint supports and $|\supp\sigma|+|\supp\tau|\le5$, then $\langle S_N,G^+(X)\rangle=S_N^+$. This applies, for instance, to any graph containing two disjoint pairs of twins, i.e.\ pairs of vertices with the same neighborhood.
\end{corollary}

\begin{proof}
By Proposition~\ref{prop:disjoint} and monotonicity, it is enough to treat the group dual. It moves at most $5$ points and is non-classical. So locality applies (Proposition~\ref{prop:monloc}): maximality for $N\le5$, together with~\cite[Cor.~3.4]{bcf}.
\end{proof}

When $|\supp\sigma|+|\supp\tau|\ge6$, locality does not apply. The group dual of Proposition~\ref{prop:disjoint} then has \emph{repeated} blocks. For instance, $\sigma=(12)(34)$ and $\tau=(56)$ give the family $(\langle x\rangle,\langle x\rangle,\langle y\rangle)$ of $\Z_2*\Z_2$, i.e.\ $u=\mathrm{diag}(M(p),M(p),M(q))$. Such actions have no quantum subgroups with smaller support. By Proposition~\ref{prop:reduction} we pass to JNA quotients of $\Z_\ell*\Z_{\ell'}$, and test them. The smallest cases are those where the two supports have total size $6$, $7$ or $8$, so that $N\le9$, and where the quotients are the smallest JNA groups of Example~\ref{ex:table}, namely $S_3$, $D_4$, $D_5$ and $A_4$. The test itself is the one of Section~\ref{sec:duals}, with the same program, a family with repeated blocks being just a family with repetitions, and the outcome is the same as before.

\begin{proposition}[computational]\label{prop:repeated}
For the following group duals with repeated blocks, the joint dimensions with $S_N$ are Catalan for $4\le k\le7$ (for $4\le k\le6$ when $N=9$):
\begin{center}
\begin{tabular}{@{}llll@{}}\toprule
$\Gamma$ & family (blocks) & from $(\sigma,\tau)$ & $N$\\\midrule
$S_3$ & $\langle(12)\rangle^2,\ \langle(23)\rangle$ & $(12)(34)$, $(56)$ & 6\\
$D_4$ & $\langle s\rangle^2,\ \langle sr\rangle$ & $(12)(34)$, $(56)$ & 6\\
$D_5$ & $\langle s\rangle^2,\ \langle sr\rangle$ & $(12)(34)$, $(56)$ & 6\\
$S_3$ & $\langle(12)\rangle^2,\ \langle(123)\rangle$ & $(12)(34)$, $3$-cycle & 7\\
$S_3$ & $\langle(12)\rangle^2,\ \langle(23)\rangle^2$ & $(12)(34)$, $(56)(78)$ & 8\\
$S_3$ & $\langle(12)\rangle^3,\ \langle(23)\rangle$ & $(12)(34)(56)$, $(78)$ & 8\\
$A_4$ & $\langle(123)\rangle^2,\ \langle(234)\rangle$ & two $3$-cycles, one $3$-cycle & 9\\\bottomrule
\end{tabular}
\end{center}
Here $H^m$ means the subgroup $H$ repeated $m$ times.
\end{proposition}

So no counterexample arises from disjoint automorphisms in these ranges. A general proof that repeated blocks are always absorbed is open. It would follow from the general statement suggested in Section~\ref{sec:dualsconcl}.

\subsection{Global quantum symmetry: solution groups}

The graphs that escape Proposition~\ref{prop:disjoint} are those whose quantum symmetry is \emph{global}. The main known examples come from binary linear systems~\cite{drs,rs}.
\begin{itemize}
\item A system $Mx=b$ over $\Z_2$, with $n$ variables, has a solution group $\Gamma$ generated by involutions $x_1,\dots,x_n$ (and a central involution $J$). Variables occurring in a common equation commute, and each equation $i$ gives the relation $\prod_{j}x_j^{M_{ij}}=J^{b_i}$.
\item In the graphs of~\cite{drs,rs}, $\Gamma$ (or a quotient of it) acts on the $2n$ variable vertices by
\[
u=\mathrm{diag}\big(M(p_1),\dots,M(p_n)\big),\qquad p_j=\tfrac{1+x_j}2 .
\]
In the language of Section~\ref{sec:duals}, this is the Fourier family $(\langle x_1\rangle,\dots,\langle x_n\rangle)$, with all blocks of type $\Z_2$.
\end{itemize}
We test this action of $\widehat\Gamma$ on $2n$ points as a candidate in its own right.

When all the blocks are of type $\Z_2$, the notions of Section~\ref{sec:laws} simplify, as follows. Since $\mathrm{Aut}(\Z_2)$ is trivial, a label is either trivial or not. So a pattern on $k$ legs is simply a pair $(\kappa,M)$, where $\kappa\in P(k)$ records which legs lie on the same orbit, that is, carry the same generator, and $M\subset\{1,\dots,k\}$ is the set of legs carrying the nontrivial element. Such a pattern is balanced when each block of $\kappa$ contains an even number of legs of $M$, and it is active when some assignment of \emph{distinct} generators to the blocks of $\kappa$ makes the ordered word along $M$ nontrivial in $\Gamma$.

\begin{corollary}[laws lemma, $\Z_2$-blocks]\label{cor:z2laws}
Let $\widehat\Gamma\subset S_{2n}^+$ act through a Fourier family of $\Z_2$-blocks $\langle x_1\rangle,\dots,\langle x_n\rangle$.
\begin{enumerate}
\item The spaces $\Fix_{\langle S_{2n},\widehat\Gamma\rangle}(u^{\otimes k})$ depend only on the set of active patterns $(\kappa,M)$ with $k$ legs.
\item If $\Gamma$ has no nontrivial laws with $\le k$ legs, then $\Fix_{\langle S_{2n},\widehat\Gamma\rangle}(u^{\otimes l})=\Fix_{S_{2n}^+}(u^{\otimes l})$ for all $l\le k$.
\end{enumerate}
\end{corollary}

\begin{proof}
(1) This is Lemma~\ref{lem:laws}, with all $H_i\simeq\Z_2$.

(2) By Corollary~\ref{cor:laws}, up to level $k$ the quantum group $\langle S_{2n},\widehat\Gamma\rangle$ has the same intertwiners as $\langle S_{2n},T\rangle$, where $T$ is the dual of $\Z_2^{*n}$, with its tautological family. By killing all but two generators, $T$ contains the dual of $\Z_2*\Z_2$, which is non-classical and moves $4$ points. So, by locality, Proposition~\ref{prop:monloc}, we have
\[
\langle S_{2n},\widehat\Gamma\rangle=\langle S_{2n},T\rangle=S_{2n}^+
\]
at all levels $l\le k$, as claimed.
\end{proof}

\begin{proposition}[Mermin--Peres magic square~\cite{me}, $N=18$]\label{prop:magic}
Here the magic square system has $9$ variables and $6$ equations. Its realization by real $2$-qubit Pauli operators has image $2^{1+4}_+$, of class~(a) in Proposition~\ref{prop:jna}. The dual of this image is a quantum subgroup of $\widehat\Gamma$. For it, the joint dimensions at $k=4,\dots,8$ are
\[
14,\ 42,\ 132,\ 429,\ 1430 ,
\]
which are Catalan, compared with the Bell numbers $15,52,203,877,4140$. This is despite the short law $(ab)^4=1$, which comes from exponent $4$.
\end{proposition}

The computation was done by exact linear algebra modulo a prime, with the patterns $(\kappa,M)$ of Corollary~\ref{cor:z2laws} (\texttt{sol.py}, $41$ seconds at $k=8$). The values at $k=4,5$ were confirmed independently by \texttt{duals.py}. At $k=8$, exactly one of the $17284$ balanced patterns which are not freely trivial is a law, namely $(ab)^4$.

\begin{proposition}[the $A_7$ system, $N=210$]\label{prop:a7}
The system of~\cite{drs} has $105$ variables, namely the double transpositions of $A_7$ (the edges of the Kneser graph $K(7,2)$), and $140$ equations. Its solution group is $3.A_7$. Up to relabeling, there are $291$ words of length $\le8$ in which every letter occurs an even number of times and which are not freely trivial, and each of them admits distinct double transpositions with nontrivial product in $A_7$ (\texttt{sol.py}). So there are no nontrivial laws of length $\le8$, already in the JNA quotient $A_7$. By Corollary~\ref{cor:z2laws}, $\langle S_{210},\widehat\Gamma\rangle$ agrees with $S_{210}^+$ through level $8$.
\end{proposition}

\subsection{Vertex-transitive graphs on at most $13$ vertices}

For small vertex-transitive graphs, the existence of quantum symmetry is completely known:
\begin{itemize}
\item $N\le11$ by Banica and Bichon (see~\cite{sch2}), except the Petersen graph, which was settled by Schmidt~\cite{sch};
\item $N=12$ by Schanz~\cite{sz};
\item $N=13$ by Chassaniol~\cite{cha1,cha2}.
\end{itemize}
In this range, a graph has quantum symmetry iff it has two disjoint automorphisms~\cite[Observation~1]{sz}. We combine this with Proposition~\ref{prop:disjoint}.

\begin{proposition}[computational]\label{prop:vtcount}
We generated all vertex-transitive graphs on $N\le13$ vertices, as the Cayley graphs of all groups of order $N$ together with the Petersen graph and its complement. The numbers obtained agree with the known counts. For each graph we determined the minimal total support $s=|\supp\sigma|+|\supp\tau|$ of two automorphisms of prime order with disjoint supports:
\begin{center}
\begin{tabular}{@{}lcccccccccc@{}}\toprule
$N$ & 4&5&6&7&8&9&10&11&12&13\\\midrule
vertex-transitive graphs & 4&3&8&4&14&9&22&8&74&14\\
no disjoint automorphisms & 0&1&2&2&4&5&10&6&32&12\\
$s=4$ (two pairs of twins) & 4&2&6&2&8&4&8&2&20&2\\
$s=8$, $\sigma,\tau$ of type $(2,2)$ & &&&&2&&4&&14&\\
$s=12$, $\sigma,\tau$ of type $(2,2,2)$ & &&&&&&&&8&\\\bottomrule
\end{tabular}
\end{center}
At $N=12$, the $42$ graphs with disjoint automorphisms form the $21$ complementary pairs with quantum symmetry found in~\cite{sz}. At $N=8$ the $s=8$ graphs are the cube and its complement. At $N=10$ they are $2C_5$, $K_5\times K_2$ and their complements.
\end{proposition}

So every remaining case has $\sigma$ and $\tau$ both products of $m\in\{2,3\}$ disjoint transpositions. By Proposition~\ref{prop:disjoint}, the corresponding group dual is the dual of $\Z_2*\Z_2$ acting with blocks $(\langle x\rangle^m,\langle y\rangle^m)$. This is the $4$-point action $\mathrm{diag}(M(p),M(q))$ with every point repeated $m$ times. We now exploit this structure.

\begin{definition}
The \emph{$m$-fold repetition} is the morphism $\varphi_m:S_n^+\to S_{nm}^+$ given by $u\mapsto u\otimes1_m$. Here the point set is $\{1,\dots,n\}\times\{1,\dots,m\}$, and $(u\otimes1_m)_{(i,c),(j,c')}=u_{ij}\delta_{cc'}$, which is magic.
\end{definition}

\begin{lemma}\label{lem:rep}
For quantum subgroups $A,B\subset S_n^+$, we have $\varphi_m(\langle A,B\rangle)=\langle\varphi_m(A),\varphi_m(B)\rangle$.
\end{lemma}

\begin{proof}
For $L\subset S_n^+$, the fundamental representation of $\varphi_m(L)$ is $u_L\otimes1_m$ on $\C^n\otimes\C^m$, so
\[
\mathrm{Hom}_{\varphi_m(L)}(k,l)=\mathrm{Hom}_L(k,l)\otimes B\big((\C^m)^{\otimes k},(\C^m)^{\otimes l}\big).
\]
The intertwiners of a generated quantum group are the intersections of the intertwiners. Since $(U\otimes W)\cap(V\otimes W)=(U\cap V)\otimes W$ for subspaces $U,V$, we get
\[
\mathrm{Hom}_{\langle\varphi_m(A),\varphi_m(B)\rangle}=\big(\mathrm{Hom}_A\cap\mathrm{Hom}_B\big)\otimes B(\cdot,\cdot)=\mathrm{Hom}_{\varphi_m(\langle A,B\rangle)} .
\]
Thus the two quantum groups in the statement have the same intertwiners, and so they are equal, by Tannakian duality.
\end{proof}

In order to deal with the remaining cases, we would need to know that repeated copies of $S_4^+$, together with the symmetric group, generate the whole quantum permutation group. This is a topological generation question, of the type studied by Chirvasitu~\cite{ch2} and by Brannan, Chirvasitu and Freslon~\cite{bcf}. Their methods work well in large dimensions, but usually lead to difficult generation questions in small dimensions, and here we are precisely in the smallest dimension, $4$, where $S_4^+$ is still coamenable, and has a rather special representation theory. For our purposes, we would need the following statement.

\begin{conjecture}[repetition lemma]\label{conj:rep}
$\langle S_{4m},\varphi_m(S_4^+)\rangle=S_{4m}^+$ for $m\ge1$.
\end{conjecture}

For $m=1$ there is nothing to prove. For $m=2,3$ the joint dimensions are Catalan for $k\le7$ (exact, \texttt{doubling.py}). The computation uses the fact that a vector $\sum_\pi a_\pi\xi_\pi$ is fixed by $\varphi_m(S_4^+)$ if and only if, for every partition $\kappa$ of the legs with at most $m$ blocks, $\sum_{\pi\le\kappa}a_\pi\xi_\pi^{(4)}$ lies in the span of the noncrossing partition vectors at $N=4$. With the Gram matrix of the vectors $\xi_\pi^{(4)}$, this becomes a rank computation, done modulo a prime, so that a Catalan output is a proof, as in Section~\ref{sec:duals}.

\begin{theorem}\label{thm:vt}
Let $X$ be a vertex-transitive graph on $N\le13$ vertices.
\begin{enumerate}
\item If $X$ has no quantum symmetry, then $\langle S_N,G^+(X)\rangle=S_N$.
\item If $X$ has two pairs of twins with disjoint supports, then $\langle S_N,G^+(X)\rangle=S_N^+$. This covers all graphs with quantum symmetry except the $14$ complementary pairs of Proposition~\ref{prop:vtcount} with $s\in\{8,12\}$.
\item If $s=8$ and $N\ge10$ (the $2+7$ pairs at $N=10,12$), then $\langle S_N,G^+(X)\rangle=S_N^+$.
\item For the remaining $5$ pairs (the cube at $N=8$, and the $s=12$ graphs at $N=12$), $\langle S_N,G^+(X)\rangle$ contains a copy of $\varphi_m(S_4^+)$, with $m=2$ or $3$, acting on $4m$ of the points. Hence $\langle S_N,G^+(X)\rangle=S_N^+$, provided that Conjecture~\ref{conj:rep} holds for that $m$.
\end{enumerate}
In particular, no vertex-transitive graph on at most $13$ vertices gives an exotic quantum group $S_N\subsetneq\langle S_N,G^+(X)\rangle\subsetneq S_N^+$, except possibly the cube and the four $s=12$ pairs, and those only if the repetition lemma fails.
\end{theorem}

\begin{proof}
(1) is clear.

(2) is Corollary~\ref{cor:small}.

(3) Here $m=2$, and there are $N-8\ge2$ fixed points. Let $G_5=G_4\oplus1$ be the $4$-point action of the dual of $\Z_2*\Z_2$ with one extra fixed point. Then $\langle S_5,G_5\rangle=S_5^+$, by the maximality of $S_5\subset S_5^+$~\cite{bah}. By Proposition~\ref{prop:disjoint}, $G^+(X)$ contains, up to relabeling, $\varphi_2(G_5)$, and by Lemma~\ref{lem:rep}, $\langle S_{10},\varphi_2(G_5)\rangle\supset\varphi_2(S_5^+)$, the diagonal embedding of $S_5^+$ in $S_{10}^+$. Then~\cite[Prop.~3.10 and Ex.~3.9]{bcf} gives $\langle S_{10},\varphi_2(S_5^+)\rangle=S_{10}^+$, and finally $\langle S_N,S_{10}^+\rangle=S_N^+$ by~\cite[Cor.~3.4]{bcf}.

(4) By Proposition~\ref{prop:disjoint}, $G^+(X)$ contains $\varphi_m(G_4)$, extended by the identity on the remaining $N-4m$ points, with $G_4$ the dual of $\Z_2*\Z_2$ on $4$ points. Since $S_4\subset S_4^+$ is maximal, $\langle S_4,G_4\rangle=S_4^+$, so by Lemma~\ref{lem:rep} the quantum group $\langle S_{4m},\varphi_m(G_4)\rangle$, acting on the support, contains $\varphi_m(S_4^+)$. Given the conjecture, it contains $S_{4m}^+$, and $\langle S_N,S_{4m}^+\rangle=S_N^+$ by~\cite[Cor.~3.4]{bcf}.
\end{proof}

\begin{remark}
For $S_M^+$ with $M\ge5$ in place of $S_4^+$, the repetition statement is known: diagonal embeddings are $(M,N)$-large, so $\langle S_{KM},\mathrm{diag}(S_M^+)\rangle=S_{KM}^+$~\cite[Ex.~3.9, Prop.~3.10]{bcf}. Their proof uses the linear independence of noncrossing partitions on a space of dimension $M-1\ge4$, which fails for $M=4$. So Conjecture~\ref{conj:rep} is exactly the missing case $M=4$. Note also that Lemma~\ref{lem:rep} applies to any repetition of a locally absorbed action. So the conjecture, for general $m$, would settle every group dual whose blocks are those of a $\le5$-point action repeated uniformly.
\end{remark}

\subsection{Lessons}

\begin{remark}
We draw the following lessons from this section.
\begin{itemize}
\item \emph{Locality.} Quantum symmetry living on a few points (disjoint automorphisms with small supports) is absorbed by $S_N$. This follows from~\cite{bcf}.
\item \emph{Repeated blocks.} Disjoint automorphisms with larger supports give group duals with repeated blocks. These are absorbed in every case computed (Proposition~\ref{prop:repeated}). For uniform repetitions they reduce to the repetition lemma (Conjecture~\ref{conj:rep}), and this covers all vertex-transitive graphs on $\le13$ vertices (Theorem~\ref{thm:vt}).
\item \emph{Freeness.} Global quantum symmetry coming from large perfect groups ($A_7$, class~(c)) has no short nontrivial laws. Such groups look free at every computable level, and are absorbed.
\item \emph{Short laws are not enough.} The magic square group ($2^{1+4}_+$, class~(a)) has short nontrivial laws, and is still absorbed through level~$8$.
\end{itemize}
So graphs, as far as they can be tested, reduce to group duals, and add nothing beyond Section~\ref{sec:duals}. Whether a graph with global quantum symmetry could ever give an exotic $\langle S_N,G^+(X)\rangle$ is the same question as for group duals: whether some JNA group has enough short laws.
\end{remark}


%======================================================================
% Section: twists
%======================================================================

\section{Cocycle twists of \texorpdfstring{$S_N$}{SN} and \texorpdfstring{$PU_n$}{PUn}}\label{sec:twists}

We continue our search for seeds $K\subset S_N^+$. The group duals and the quantum automorphism groups of graphs are absorbed by $S_N$, as far as we can compute, so we turn to a third family, the cocycle twists of $S_N$ and of its subgroups, and then to the twists of the projective unitary groups $PU_n$, acting on $N=n^2$ points.

\subsection{Twists of $S_N$}\label{sec:twsn}

Let $A\subset S_N$ be an abelian subgroup, and let $\sigma$ be a $2$-cocycle on the dual group $\widehat A$ with values in roots of unity. The restriction map $C(S_N)\to C(A)=\C[\widehat A]$ allows us to view $\sigma$ as a Hopf $2$-cocycle on $C(S_N)$. Twisting by it, in the sense of Drinfeld~\cite{dr}, Doi~\cite{doi} and Enock--Vainerman~\cite{ev,va}, gives a finite quantum group $S_N^\sigma$:
\begin{itemize}
\item its tensor category is equivalent to that of $S_N$;
\item its coaction is a twisted version of the coaction of $S_N$ on $\C^N$.
\end{itemize}
Twists of this type, with $A\simeq\Z_2^n$ generated by the disjoint transpositions $(2i-1,2i)$ of $S_{2n}$, and $\sigma$ with values $\pm1$, go back to Nikshych~\cite{ni} and Bichon~\cite{bi00}. Bichon constructed in this way a family of noncommutative and noncocommutative Kac algebras $\mathcal O_{\mathbf p}(S_{2n})$, whose representation categories are monoidally equivalent to that of $S_{2n}$, and whose defining relations contain the magic unitary relations, so that they are quantum permutation groups~\cite[Def.~1, Thm.~2]{bi00}. Banica, Bichon and Natale give a sufficient condition for a twist of a quantum permutation algebra to be a quantum permutation algebra, and observe that Bichon's twistings satisfy it~\cite[Prop.~8.1, Ex.~8.2]{bbn}. Proposition~\ref{prop:comm} below is our own formulation of what is needed in the setting of $S_N$.

The twisted coaction lives on the algebra $\C^N_\sigma$. This is $\C^N$ with the product deformed along the $\widehat A$-grading coming from the action of $A$: for homogeneous $f,g$ of degrees $\chi,\psi$,
\[
f*g=\sigma(\chi,\psi)\,fg .
\]
For $S_N^\sigma$ to be a quantum \emph{permutation} group, this algebra must be commutative.

\begin{proposition}\label{prop:comm}
For an orbit $O$ of $A$, write $A_O$ for the stabilizer of a point of $O$, and $A_O^\perp\subset\widehat A$ for the characters trivial on $A_O$. Then $\C^N_\sigma$ is commutative iff, for every orbit $O$, the restriction of $\sigma$ to $A_O^\perp$ is symmetric, i.e.\ cohomologically trivial. In that case $\C^N_\sigma\simeq\C^N$ and $S_N^\sigma\subset S_N^+$.
\end{proposition}

\begin{proof}
The functions supported on different orbits multiply to zero, so it suffices to look at one orbit $O$. The degrees occurring in $C(O)$ are exactly the characters in $A_O^\perp\simeq\widehat{A/A_O}$, since $A$ acts on $O$ regularly through $A/A_O$. So $C(O)$ is commutative iff $\sigma(\chi,\psi)=\sigma(\psi,\chi)$ on $A_O^\perp$. For abelian groups and $\C^\times$ coefficients, this is equivalent to $\sigma|_{A_O^\perp}$ being a coboundary. In that case $\C^N_\sigma$ is a commutative finite-dimensional $C^*$-algebra of dimension $N$, hence isomorphic to $\C^N$, and the twisted coaction becomes a coaction of $S_N^\sigma$ on $\C^N$. Since $S_N^+$ is the universal compact quantum group acting on $\C^N$~\cite{wa}, we obtain $S_N^\sigma\subset S_N^+$. See also~\cite[Prop.~8.1, Ex.~8.2]{bbn}, and~\cite{bi00} for Bichon's twists.
\end{proof}

\begin{remark}\label{rem:twcases}
Two special cases:
\begin{itemize}
\item If every orbit group $A/A_O$ is cyclic, the condition holds automatically, since $2$-cocycles on cyclic groups are trivial. This covers the twists of Proposition~\ref{prop:twres} below.
\item If $A$ acts \emph{regularly} and $\sigma$ is non-degenerate, then $\C^N_\sigma$ is a matrix algebra, and $S_N^\sigma$ is not a quantum permutation group: it acts on $M_n$, with $N=n^2$. In this case the dimension count of Proposition~\ref{prop:J} below can give joint dimensions below Catalan, which is impossible inside $S_N^+$. This case reappears in Section~\ref{sec:twalg}, read backwards.
\end{itemize}
Proposition~\ref{prop:comm} allows more than the cyclic case: non-cyclic orbits, such as Klein orbits, are fine as long as $\sigma$ only couples \emph{different} orbits. This mirrors the role of non-cyclic orbits in Section~\ref{sec:duals}.
\end{remark}

\subsection{Twisted intertwiners in the Fourier basis}

Take $A=\Z_m^n$ acting on $N=mn$ points through $n$ disjoint $m$-cycles. Take the bicharacter
\[
\sigma(e_i,e_j)=\omega^{b_{ij}}\quad(i<j),\qquad \omega=e^{2\pi i/m},
\]
with $b$ strictly upper triangular. On the orbit $q$ we use the Fourier basis $f_{(q,a)}$, $a\in\Z_m$, as in Section~\ref{sec:duals}. These vectors are weight vectors for $A$.

\begin{proposition}\label{prop:J}
\begin{enumerate}
\item In the coordinates $(q_s,a_s)_{s\le k}$, a partition vector $\xi_\pi$ is supported on the tuples that are \emph{balanced} for $\pi$: on each block, $q$ is constant and $\sum a_s\equiv0\pmod m$. This is Proposition~\ref{prop:coords} with $H=\Z_m$.
\item The twist acts diagonally in these coordinates:
\[
\Fix_{S_N^\sigma}(k)=\Theta\big(\Fix_{S_N}(k)\big),\qquad \Theta=\prod_{s<t}\omega^{\,b_{q_sq_t}a_sa_t}.
\]
\item Consequently, with $V=\spn\{\xi_\pi\}$,
\[
\dim\Fix_{\langle S_N,S_N^\sigma\rangle}(k)=\dim(V\cap\Theta V)=2\dim V-\mathrm{rank}\,[V\,|\,\Theta V].
\]
\end{enumerate}
\end{proposition}

We compute these ranks exactly modulo the prime $p=1048261\equiv1\pmod{60}$, which contains the roots of unity of orders $2,3,4,5$, after a random compression, which can only lower the rank, so that a Catalan output is a proof (\texttt{twists.py}). In practice $\omega$ is replaced by a primitive root of unity of the same order in $\mathbb F_p$, the matrix $V$ is the one of Section~\ref{sec:duals}, with columns indexed by $P(k)$ and rows by the balanced tuples, and $\Theta V$ is obtained from it by rescaling the rows, by the formula in Proposition~\ref{prop:J}(2). The whole computation therefore takes place over $\mathbb F_p$, with no floating point at all.

\emph{The ground field.} All dimensions above are dimensions of intersections of subspaces defined over $\mathbb Q(\omega)$. So they do not change over any algebraically closed field of characteristic $0$, which is the setting of~\cite{bbn}.

\subsection{Results}

\begin{proposition}[computational]\label{prop:twres}
For the following twists, the joint dimensions are the Catalan numbers at all the levels listed.
\begin{center}
\begin{tabular}{@{}llll@{}}\toprule
$N$ & $A$ & cocycles tested, up to symmetry & levels $k$\\\midrule
4 & $\Z_2^2$, on two transpositions & the non-trivial one ($S_4^\tau$) & $4$--$6$\\
6 & $\Z_2^3$ & all $3$ sign graphs & $4$--$7$\\
8 & $\Z_2^4$ & all $10$ sign graphs & $4$--$7$\\
6 & $\Z_3^2$ & $b=1$ & $4$--$7$\\
8 & $\Z_4^2$ & $b=1,2$ & $4$--$7$\\
9 & $\Z_3^3$ & all $6$ classes up to orbit permutations & $4$--$6$\\
10 & $\Z_5^2$ & $b=1$ & $4$--$6$\\\bottomrule
\end{tabular}
\end{center}
For $m=2$, the cocycle is encoded by the \emph{sign graph} on the $n$ transpositions, with an edge $i$--$j$ iff $\sigma(e_i,e_j)\sigma(e_j,e_i)^{-1}=-1$.
\end{proposition}

The $N=4$ line is a consistency check. There $S_4^\tau\subset S_4^+$ is non-classical, so $\langle S_4,S_4^\tau\rangle=S_4^+$ by maximality at $N=4$~\cite{bb}. The check also confirms that, with our conventions, $NC$ sits inside the twisted category.

\begin{proposition}[computational, non-cyclic and mixed orbits]\label{prop:twklein}
Let $A=H_1\times\dots\times H_r$, where each $H_q$ acts regularly on its own orbit, and let $\sigma$ be a bicharacter coupling different orbits only. Such twists are quantum permutation groups by Proposition~\ref{prop:comm}. For the following cases, the joint dimensions are Catalan for $k=4,5,6,7$.
\begin{center}
\begin{tabular}{@{}lll@{}}\toprule
$N$ & orbit groups $H_q$ & coupling\\\midrule
6 & $\Z_2^2,\ \Z_2$ & non-trivial\\
6 & $\Z_4,\ \Z_2$ & non-trivial\\
8 & $\Z_2^2,\ \Z_2^2$ & rank $1$ and rank $2$\\
8 & $\Z_2^2,\ \Z_4$ & non-trivial\\
8 & $\Z_2^2,\ \Z_2,\ \Z_2$ & chain ($1$--$2$--$3$) and star (center $\Z_2^2$)\\\bottomrule
\end{tabular}
\end{center}
\end{proposition}

\begin{remark}[locality, heuristically]
A bicharacter is a product of contributions from pairs of orbits. Restricted to the points of two coupled orbits, a twist looks like a twist of a smaller symmetric group, which is non-classical on few points. For two coupled transpositions this is $S_4^\tau$, which is absorbed by locality (Proposition~\ref{prop:monloc}). We have not proved that $S_N^\sigma$ actually contains such a piece as a quantum subgroup. When the sign graph has more edges, the cross terms get in the way. So this is a heuristic explanation of the tables, not a theorem.
\end{remark}

\subsection{Twists of smaller subgroups}

By monotonicity, smaller seeds are better. So instead of twisting all of $S_N$, we twist a subgroup $L$ with $A\subset L\subset S_N$. The same construction gives $L^\sigma\subset S_N^\sigma$, and the twisted algebra is the same $\C^N_\sigma$. So Proposition~\ref{prop:comm} applies unchanged, and so does the formula
\[
\Fix_{L^\sigma}(k)=\Theta\big(\Fix_L(k)\big).
\]
Hence $\xi\in\spn\{\xi_\pi\}$ lies in $\Fix_{L^\sigma}$ iff $\Theta^*\xi$ is $L$-invariant. Invariance under $A$ is automatic, so it suffices to check the other generators of $L$.

When $L$ normalizes $A$, its elements act on the Fourier basis by monomial matrices: they permute the orbits and act by automorphisms on the labels, up to phases. The test then stays inside the Fourier coordinates (\texttt{twists.py}). For $A=\Z_m^n$ acting on $n$ disjoint $m$-cycles, the full normalizer is
\[
N_{S_N}(A)=(\Z_m\rtimes\Z_m^\times)\wr S_n ,
\]
the wreath product of the affine group of $\Z_m$ with $S_n$, and all its elements are of this type.

\begin{remark}[which $L^\sigma$ are non-classical]\label{rem:classical}
If every element of $L$ preserves the alternating form $\sigma(\chi,\psi)\sigma(\psi,\chi)^{-1}$, the twist has no effect: $L^\sigma$ is classical, and the joint dimensions are the Bell numbers. We observed this in three cases:
\begin{itemize}
\item $L=D_4\subset S_4$, which normalizes $A=\Z_2^2$, since every automorphism of $\Z_2^2$ preserves its non-degenerate form;
\item the triangle sign graph at $N=6$, with $L$ swapping two orbits;
\item two Klein orbits with a rank-$1$ coupling, with $L$ swapping the orbits.
\end{itemize}
The useful seeds are therefore $L=A\rtimes\langle g\rangle$ with $g$ changing the form: a permutation of the orbits that is not an automorphism of the sign graph, or the inversion $a\mapsto-a$ on a $\Z_3$-orbit. These are essentially the smallest non-classical seeds of twist type. A joint dimension strictly below Bell proves that $L^\sigma$ is non-classical.
\end{remark}

\begin{proposition}[computational]\label{prop:twsub}
For the following seeds $L^\sigma$, with $L=A\rtimes\langle g\rangle$, the joint dimensions are Catalan for $k=4,5,6,7$. In particular all these $L^\sigma$ are non-classical, and they are absorbed by $S_N$ through level $7$.
\begin{center}
\begin{tabular}{@{}llll@{}}\toprule
$N$ & $A$ & $\sigma$ & $g$\\\midrule
6 & $\Z_2^3$ & edge $\{1,2\}$ & swap of orbits $2,3$\\
6 & $\Z_2^3$ & path $1$--$2$--$3$ & swap of orbits $1,2$\\
6 & $\Z_3^2$ & $b=1$ & swap of the two orbits\\
6 & $\Z_3^2$ & $b=1$ & $a\mapsto-a$ on one orbit\\
8 & $\Z_2^4$ & edge $\{1,2\}$ & swap of orbits $2,3$\\
8 & $\Z_2^4$ & edges $\{1,2\},\{3,4\}$ & swap of orbits $2,3$\\
8 & $\Z_2^4$ & path $1$--$2$--$3$--$4$ & swap of orbits $1,2$\\
8 & $\Z_4^2$ & $b=1$ & swap of the two orbits\\\bottomrule
\end{tabular}
\end{center}
\end{proposition}

As a continuation of this, it would be interesting to twist subgroups which do not normalize $A$.

\begin{question}\label{qu:nonnormal}
Compute $\langle S_N,L^\sigma\rangle$ for subgroups $A\subset L\subset S_N$ which do not normalize $A$, for instance primitive groups.
\end{question}

When $L$ normalizes $A$, the test stays in the Fourier coordinates, where the generators of $L$ act by monomial matrices. Otherwise these generators are dense in the Fourier basis, and the invariance test is more expensive, but levels $k\le6$ for $N\le9$ should be within reach.

So, at the levels we can reach, twists of $S_N$ and of its subgroups behave like group duals and graphs: everything is absorbed, including the smallest twisted seeds. All the twists above, however, live on bicharacters, which only couple pairs of orbits. We now turn to twists of a different kind, where the group being twisted is not a permutation group at all.

%---------------------------------------------------------------------
\subsection{Twisting an algebra}\label{sec:twalg}

The mechanism of Proposition~\ref{prop:comm} needs neither $S_N$ nor $\C^N$. Let $G$ be a compact group, acting by automorphisms on a finite-dimensional $C^*$-algebra $B$ and preserving a faithful trace $\mathrm{tr}$, and let $A\subset G$ be a finite abelian subgroup. The restricted action grades $B=\bigoplus_{\chi\in\widehat A}B_\chi$, and a $2$-cocycle $\sigma$ on $\widehat A$ with values in $\T$ defines, as in Section~\ref{sec:twsn}, the twisted algebra $B_\sigma$, with product
\[
f*g=\sigma(\chi,\psi)fg\qquad(f\in B_\chi,\ g\in B_\psi),
\]
and the twisted compact quantum group $G^\sigma$, whose algebra $C(G^\sigma)$ is $C(G)$ with the product twisted by $\sigma$ on the left and by $\sigma^{-1}$ on the right, with respect to the gradings coming from the left and right translations by $A$.

\begin{proposition}\label{prop:twalg}
$G^\sigma$ coacts on $B_\sigma$, preserving $\mathrm{tr}$. If $B_\sigma$ is commutative and $\mathrm{tr}$ becomes the uniform trace on $B_\sigma\simeq\C^N$, where $N=\dim B$, then $G^\sigma\subset S_N^+$. If moreover $G$ acts faithfully on $B$, this is a closed quantum subgroup.
\end{proposition}

\begin{proof}
The coaction $\beta:B\to B\otimes C(G)$ is unchanged as a linear map. It is multiplicative for the twisted products on both sides, as one checks on homogeneous elements, which is the standard twisting of comodule algebras~\cite{doi}, see also~\cite[Sec.~8]{bbn}. The invariance of $\mathrm{tr}$ is a linear identity, so it survives. When $B_\sigma\simeq\C^N$ with the uniform trace, Wang's universality~\cite{wa} gives a morphism $C(S_N^+)\to C(G^\sigma)$. Its image contains the coefficients of $\beta$, and when $G$ acts faithfully these generate $C(G)$, hence also $C(G^\sigma)$, since the twisted product of two coefficients is a coefficient of the tensor product of the corresponding comodules.
\end{proof}

Proposition~\ref{prop:comm} is the case $B=\C^N$, $G=S_N$. The second item of Remark~\ref{rem:twcases} goes from $\C^N$ to a matrix algebra. The example below goes the other way, from a matrix algebra to $\C^N$, and produces quantum permutation groups which are twists of the projective unitary groups.

Let $H$ be a finite abelian group, with $|H|=n$, and $\Lambda=H\times\widehat H$, so that $|\Lambda|=n^2=N$. The Weyl matrices $W_{(i,a)}=\sum_ka(k)E_{k,k+i}$, acting on $\ell^2(H)=\C^n$, satisfy
\[
W_\lambda W_\mu=c(\lambda,\mu)W_{\lambda+\mu},\qquad c\big((i,a),(j,b)\big)=b(i),
\]
so that $M_n=C^*_c(\Lambda)$, the twisted group algebra~\cite[Prop.~2.7]{bn}. The commutator of two Weyl matrices is a phase,
\[
W_\lambda W_\mu W_\lambda^*=\varepsilon(\lambda,\mu)W_\mu,\qquad \varepsilon(\lambda,\mu)=c(\lambda,\mu)\overline{c(\mu,\lambda)},
\]
with $\varepsilon$ a non-degenerate alternating bicharacter on $\Lambda$, i.e.\ $\varepsilon((i,a),(j,b))=b(i)\overline{a(j)}$. So conjugation by the Weyl matrices gives an abelian subgroup $A\simeq\Lambda$ of $PU_n=\Aut(M_n)$, the grading of $M_n$ has the one-dimensional components $\C W_\mu$, and we identify $\widehat A$ with $\Lambda$ through $\varepsilon$. Twisting by $\sigma=\bar c$ gives $W_\lambda*W_\mu=W_{\lambda+\mu}$, so that
\[
(M_n)_\sigma=\C[\Lambda]\simeq C(\widehat\Lambda)=\C^N ,
\]
with the trace $\mathrm{tr}(W_\lambda)=\delta_{\lambda,0}$ becoming the uniform trace. Since $PU_n$ acts faithfully on $M_n$, Proposition~\ref{prop:twalg} gives a closed quantum subgroup
\[
PU_n^\sigma\subset S_N^+ ,\qquad N=n^2 .
\]
For $H=\Z_2$ this is $PU_2^\sigma=SO_3^{-1}=S_4^+$~\cite{bb}. The classical subgroup $A$ survives in $PU_n^\sigma$, and acts on the $N$ points $\widehat\Lambda$ by translations, regularly.

%---------------------------------------------------------------------
\subsection{Flat matrix models}\label{sec:flat}

The quantum groups $PU_n^\sigma$ appear in the work of Banica and Nechita on flat matrix models~\cite{bn}. Consider the representation
\[
\pi_H:C(S_N^+)\to C(U_n,M_N),\qquad w_{\lambda\mu}\mapsto\big[x\mapsto\Proj(W_\lambda xW_\mu^*)\big],
\]
where $\Proj$ is taken inside $L(M_n)\simeq M_N$, for the scalar product $\langle a,b\rangle=\mathrm{tr}(ab^*)$. This is a matrix model in the sense that, for each $x\in U_n$, the matrix $(\Proj(W_\lambda xW_\mu^*))_{\lambda\mu}$ is magic: its entries are rank one projections, and on each row and column they project on the orthonormal basis $\{W_\lambda xW_\mu^*/\sqrt n\}_\lambda$, resp.\ $\{W_\lambda xW_\mu^*/\sqrt n\}_\mu$, of $M_n$. For $H=\Z_2$ the Weyl matrices are the Pauli matrices, up to phases, and this is the Pauli representation of $S_4^+$ mentioned in Section~\ref{sec:maximality}, whose magic unitary is
\[
\big(\Proj(c_ixc_j)\big)_{ij},\qquad x\in SU_2 .
\] It is shown in~\cite[Thm.~4.3]{bn} that the main character of the Hopf image of $\pi_H$ has the same law as $x\mapsto|\mathrm{Tr}(x)|^2$ on $U_n$, which is the main character of $PU_n$, and it is suggested there that the Hopf image should be a twist of $PU_n$. This is indeed the case.

\begin{theorem}\label{thm:bn}
The representation $\pi_H$ factors as $C(S_N^+)\to C(PU_n^\sigma)\xrightarrow{\theta}C(U_n,M_N)$, with $\theta$ injective on the dense Hopf $*$-algebra $\mathcal O(PU_n^\sigma)$. Hence the Hopf image of $\pi_H$ is $PU_n^\sigma\subset S_N^+$.
\end{theorem}

\begin{proof}
\emph{Conventions.} The subgroup $A$ acts on $PU_n$ by left and right translations, and this bigrades $\mathcal O(PU_n)=\bigoplus_{\alpha,\beta\in\Lambda}\mathcal O_{\alpha\beta}$, where $f\in\mathcal O_{\alpha\beta}$ means $f(W_\gamma xW_\delta)=\varepsilon(\gamma,\alpha)\varepsilon(\delta,\beta)f(x)$. The coaction $\beta(a)(x)=xax^*$ of $PU_n$ on $M_n$ is $\beta(W_\mu)=\sum_\nu W_\nu\otimes c_{\nu\mu}$, with $c_{\nu\mu}(x)=\mathrm{tr}(W_\nu^*xW_\mu x^*)$, and the relations $W_\gamma^*W_\nu^*W_\gamma=\varepsilon(\gamma,\nu)W_\nu^*$, $W_\delta W_\mu W_\delta^*=\varepsilon(\delta,\mu)W_\mu$ show that $c_{\nu\mu}\in\mathcal O_{\nu\mu}$. Comparing the coefficients of $xW_\mu W_{\mu'}x^*=(xW_\mu x^*)(xW_{\mu'}x^*)$ gives
\[
c(\mu,\mu')\,c_{\nu,\mu+\mu'}=\sum_{\nu_1+\nu_2=\nu}c(\nu_1,\nu_2)\,c_{\nu_1\mu}c_{\nu_2\mu'} .
\]
So the twisted product on $\mathcal O(PU_n)$ for which $\beta:(M_n)_\sigma\to(M_n)_\sigma\otimes\mathcal O(PU_n^\sigma)$ is multiplicative, as in Proposition~\ref{prop:twalg}, is $a\cdot b=c(\alpha,\alpha')\overline{c(\beta,\beta')}\,ab$, for $a\in\mathcal O_{\alpha\beta}$ and $b\in\mathcal O_{\alpha'\beta'}$, see~\cite{doi,ev}. As for the points, the minimal projections of $(M_n)_\sigma=\C[\Lambda]$ are $e_\lambda=N^{-1}\sum_\beta\overline{\varepsilon(\lambda,\beta)}W_\beta$, with $W_\nu=\sum_\lambda\varepsilon(\lambda,\nu)e_\lambda$, and the magic unitary of $PU_n^\sigma$, given by $\beta(e_\mu)=\sum_\lambda e_\lambda\otimes v_{\lambda\mu}$, is therefore
\[
v_{\lambda\mu}=\frac1N\sum_{\nu,\beta}\varepsilon(\lambda,\nu)\overline{\varepsilon(\mu,\beta)}\,c_{\nu\beta} .
\]

\emph{The embedding.} Let $L_a,R_a\in L(M_n)$ be the left and right multiplications by $a\in M_n$, and define $\Phi:\mathcal O(PU_n^\sigma)\to C(U_n,L(M_n))=C(U_n,M_N)$ by
\[
\Phi(f)(x)=f(x)\,L_{W_\alpha}R_{W_\beta^*}\qquad(f\in\mathcal O_{\alpha\beta}).
\]
Since left and right multiplications commute, with $L_aL_b=L_{ab}$ and $R_aR_b=R_{ba}$, we get $L_{W_\alpha}L_{W_{\alpha'}}=c(\alpha,\alpha')L_{W_{\alpha+\alpha'}}$ and $R_{W_\beta^*}R_{W_{\beta'}^*}=R_{(W_\beta W_{\beta'})^*}=\overline{c(\beta,\beta')}R_{W_{\beta+\beta'}^*}$, so $\Phi(a\cdot b)=\Phi(a)\Phi(b)$. Moreover $\Phi$ is injective, the $N^2$ operators $L_{W_\alpha}R_{W_\beta^*}$ being a basis of $L(M_n)$.

\emph{The comparison.} Write $P_x=\Proj(x)=\sum_{\alpha\beta}p_{\alpha\beta}(x)L_{W_\alpha}R_{W_\beta}$. Since $\mathrm{Tr}(L_aR_b)=\mathrm{Tr}(a)\mathrm{Tr}(b)$, the operators $L_{W_\alpha}R_{W_\beta}$ are orthogonal for the trace of $L(M_n)$, with norms $n^2$, and since $P_x$ is the rank one projection onto the unit vector $x$, we get
\[
p_{\alpha\beta}(x)=n^{-2}\,\mathrm{tr}(W_\alpha^*xW_\beta^*x^*)=N^{-1}\,\overline{c(-\beta,\beta)}\,c_{\alpha,-\beta}(x),
\]
by using $W_\beta^*=\overline{c(-\beta,\beta)}W_{-\beta}$. Now $\Proj(W_\lambda xW_\mu^*)=UP_xU^*$, with $U=L_{W_\lambda}R_{W_\mu^*}$, and the Weyl relations give $UL_{W_\alpha}R_{W_\beta}U^*=\varepsilon(\lambda,\alpha)\varepsilon(\mu,\beta)L_{W_\alpha}R_{W_\beta}$. Thus, by substituting $\beta\to-\beta$ and using $R_{W_{-\beta}}=c(\beta,-\beta)R_{W_\beta^*}$, $c(\beta,-\beta)=c(-\beta,\beta)$ and $\varepsilon(\mu,-\beta)=\overline{\varepsilon(\mu,\beta)}$, we obtain
\[
\pi_H(w_{\lambda\mu})(x)=\frac1N\sum_{\alpha,\beta}\varepsilon(\lambda,\alpha)\overline{\varepsilon(\mu,\beta)}\,c_{\alpha\beta}(x)\,L_{W_\alpha}R_{W_\beta^*}=\Phi(v_{\lambda\mu})(x),
\]
with all the phases cancelling exactly. Hence $\pi_H=\Phi\circ q$ on $\mathcal O(S_N^+)$, where $q:w\mapsto v$ is the quotient map onto $\mathcal O(PU_n^\sigma)$. Since $\Phi$ is injective, $\ker\pi_H=\ker q$, which is a Hopf ideal, and so the Hopf image of $\pi_H$ is $PU_n^\sigma$. Equivalently, the Tannakian category of the Hopf image, which consists of the maps $T$ with $T\pi_H(w^{\otimes k})=\pi_H(w^{\otimes l})T$, see~\cite{bqp}, is that of $PU_n^\sigma$.
\end{proof}

In particular $\pi_H$ is inner faithful onto $PU_n^\sigma$, and~\cite[Thm.~4.3]{bn} becomes the statement that twisting does not change the dimensions of the fixed point spaces. We have also checked Theorem~\ref{thm:bn} numerically, by verifying that the vectors of Proposition~\ref{prop:pufix} below are fixed by $\pi_H$ at random points $x\in U_n$, up to $10^{-14}$, for $H=\Z_3$ at levels $k\le4$, and for $H=\Z_4,\Z_2^2$ at level $3$.

The fixed vectors have the same shape as in Proposition~\ref{prop:J}. We use the Fourier basis of $\C^N=C(\widehat\Lambda)$, which is the basis $\{W_\alpha\}$ of $\C[\Lambda]$, and the operator $\Theta$ on $(\C^N)^{\otimes k}$, diagonal in this basis,
\[
\Theta\,(W_{\alpha_1}\otimes\dots\otimes W_{\alpha_k})=\prod_{s<t}\sigma(\alpha_s,\alpha_t)\,W_{\alpha_1}\otimes\dots\otimes W_{\alpha_k},
\]
which is the analogue of the operator $\Theta$ of Proposition~\ref{prop:J}, with the cocycle $\sigma=\bar c$ on $\Lambda$ in place of the bicharacter $\omega^{b}$ there. The fixed vectors of $PU_n$ itself are the permutation operators, and the twist moves them by $\Theta$, as follows.

\begin{proposition}\label{prop:pufix}
For $\tau\in S_k$, let $P_\tau\in\mathrm{End}((\C^n)^{\otimes k})=M_n^{\otimes k}$ be the permutation of the tensor legs, and expand
\[
P_\tau=\sum_\alpha c_\tau(\alpha)\,W_{\alpha_1}\otimes\dots\otimes W_{\alpha_k}.
\]
Then $c_\tau(\alpha)\ne0$ iff $\alpha$ is balanced on the cycles of $\tau$, i.e.\ $\sum_{s\in\gamma}\alpha_s=0$ for each cycle $\gamma$, in which case $|c_\tau(\alpha)|=n^{\#\mathrm{cycles}(\tau)-k}$. Moreover
\[
\Fix_{PU_n^\sigma}(k)=\Theta\big(\spn\{c_\tau:\tau\in S_k\}\big),
\]
and its dimension is the number of permutations in $S_k$ having no increasing subsequence of length $n+1$.
\end{proposition}

\begin{proof}
By Schur--Weyl duality, $\Fix_{PU_n}(M_n^{\otimes k})$ is spanned by the $P_\tau$. The twist acts on fixed vectors by $\Theta$, exactly as in Proposition~\ref{prop:J}. The coefficients are $c_\tau(\alpha)=\langle P_\tau,W_\alpha\rangle=n^{-k}\mathrm{Tr}(P_\tau\,W_{\alpha_1}^*\otimes\dots\otimes W_{\alpha_k}^*)$, where $W_\alpha=W_{\alpha_1}\otimes\dots\otimes W_{\alpha_k}$, the $W_\alpha$ being orthonormal for $\langle a,b\rangle=\mathrm{tr}(ab^*)$. This is a product over the cycles of $\tau$ of normalized traces of products of Weyl matrices, which vanish unless the labels sum to zero. The dimension count is the one of~\cite[Sec.~4]{bn}.
\end{proof}

So both families of vectors live on balanced tuples, the partition vectors on the blocks of $\pi$ (Proposition~\ref{prop:J}), and the twisted vectors on the cycles of $\tau$. The joint dimensions are therefore again $\dim(V\cap\Theta W_0)$, with $V=\spn\{\xi_\pi\}$ and $W_0=\spn\{c_\tau\}$, computed as in Proposition~\ref{prop:J}. One more symmetry helps. The symplectic group $Sp(\Lambda)$ of $\varepsilon$ acts on $\widehat\Lambda$, hence on the $N$ points, and we have checked that it fixes all the vectors $\Theta c_\tau$, at the levels below. Since it also fixes the partition vectors, all ranks can be computed on the orbit representatives of $Sp(\Lambda)$ acting diagonally on $\Lambda^k$, which cuts the size by a factor close to $|Sp(\Lambda)|$, which is $24,48,720,120$ for $H=\Z_3,\Z_4,\Z_2^2,\Z_5$.

%---------------------------------------------------------------------
\subsection{Results for \texorpdfstring{$PU_n^\sigma$}{PUn sigma}}\label{sec:pures}

Before giving the numbers, observe that $PU_n^\sigma$ sits at the crossing of two maximality problems. The point is that twisting by $\sigma$, along the common regular subgroup $A$, identifies the closed quantum subgroups of $S_N^+$ containing $A$ with those of its twist containing $A$, preserving inclusions and hence generation. The twist of $S_N^+$ along a regular $A$ with non-degenerate $\bar\sigma$ is $\Aut^+(M_n,\mathrm{tr})$~\cite{bbcu}, the twist of $PU_n^\sigma$ back is $PU_n$, and the twist of $S_N$ is the quantum group $S_N^{\bar\sigma}\subset\Aut^+(M_n)$ of Remark~\ref{rem:twcases}. This correspondence is standard, but we include a proof.

\begin{lemma}\label{lem:twcorr}
Let $G$ be a compact quantum group with a closed abelian subgroup $A$, and $\sigma$ a $2$-cocycle on $\widehat A$. Then $K\mapsto K^\sigma$ is an inclusion-preserving bijection from the closed quantum subgroups $A\subset K\subset G$ onto the closed quantum subgroups $A\subset K'\subset G^\sigma$, with inverse $K'\mapsto K'^{\bar\sigma}$. In particular $\langle K,L\rangle^\sigma=\langle K^\sigma,L^\sigma\rangle$.
\end{lemma}

\begin{proof}
Let $r:\mathcal O(G)\to\mathcal O(A)$ be the restriction. The algebra $\mathcal O(G^\sigma)$ is $\mathcal O(G)$ with the same coalgebra structure, and with the product $a\cdot b=\sum\sigma(r(a_1),r(b_1))\,a_2b_2\,\sigma^{-1}(r(a_3),r(b_3))$~\cite{doi}. If $A\subset K$, then $r$ vanishes on the Hopf ideal $I=\ker(\mathcal O(G)\to\mathcal O(K))$, so for $a\in I$ only the terms with $a_2\in I$ survive, and $I$ is an ideal for the twisted product as well, and a Hopf ideal of $\mathcal O(G^\sigma)$. The quotient is $\mathcal O(K^\sigma)$, and since $\mathcal O(A)$ is commutative and cocommutative its twist is itself, so $A\subset K^\sigma$. Twisting back by $\bar\sigma$ gives the inverse map, and both maps preserve inclusions, which are inclusions of the ideals $I$. Since $\langle K,L\rangle$ is the smallest closed subgroup containing $K$ and $L$, which both contain $A$, generation is preserved too.
\end{proof}

This gives:

\begin{proposition}\label{prop:crossing}
$\langle S_N,PU_n^\sigma\rangle=S_N^+$ if and only if $\langle S_N^{\bar\sigma},PU_n\rangle=\Aut^+(M_n,\mathrm{tr})$. In particular, this equality follows both from the maximality of $S_N\subset S_N^+$, and from the maximality of $PU_n\subset PO_n^+=\Aut^+(M_n)$ in Conjecture~\ref{conj:mn}.
\end{proposition}

\begin{proof}
The equivalence follows from Lemma~\ref{lem:twcorr}, applied to $G=S_N^+$ and to the subgroups $S_N$ and $PU_n^\sigma$. For the second assertion we need $PU_n^\sigma\not\subset S_N$ and $S_N^{\bar\sigma}\not\subset PU_n$. The first holds because $PU_n^\sigma$, having the fusion rules of $PU_n$, is infinite, while $S_N$ is finite. For the second, at level $3$ the fixed point spaces have dimensions $6$ for $PU_n$ and $5$ for $S_N^{\bar\sigma}$, for any $n\ge3$, and a quantum subgroup has at least as many fixed vectors as the ambient quantum group.
\end{proof}

\begin{proposition}[computational]\label{prop:pures}
For the following $H$, the joint dimensions $\dim\Fix_{\langle S_N,PU_n^\sigma\rangle}(k)$ are the Catalan numbers $14,42,132,429$ at all the levels listed, and the noncrossing partition vectors lie in $\Fix_{PU_n^\sigma}(k)$.
\begin{center}
\begin{tabular}{@{}lllll@{}}\toprule
$N$ & $H$ & levels $k$ & $\dim\Fix_{PU_n^\sigma}(k)$ & $|Sp(\Lambda)|$\\\midrule
$4$ & $\Z_2$ & $4$--$6$ & $14,42,132$ & $6$\\
$9$ & $\Z_3$ & $4$--$7$ & $23,103,513,2761$ & $24$\\
$16$ & $\Z_4$ & $4$--$6$ & $24,119,694$ & $48$\\
$16$ & $\Z_2^2$ & $4$--$6$ & $24,119,694$ & $720$\\
$25$ & $\Z_5$ & $4$--$5$ & $24,120$ & $120$\\\bottomrule
\end{tabular}
\end{center}
\end{proposition}

All the ranks were computed exactly modulo the prime $p=1048261\equiv1\pmod{60}$, after a random compression, which can only lower the rank, so that the Catalan outputs are proofs (\texttt{twistpu.py}). The rows were reduced to the orbit representatives of $Sp(\Lambda)$, after checking on random samples that $Sp(\Lambda)$ fixes the vectors $\Theta c_\tau$. For $H=\Z_3$, levels $4$ and $5$ were also computed without this reduction, with the same results. The convention for $\Theta$ matters here. With the coefficients $c_\tau(\alpha)=\langle P_\tau,W_\alpha\rangle$, the twist acts by $\Theta(\alpha)=\prod_{s<t}c(\alpha_s,\alpha_t)=\prod_{s<t}\overline{\sigma(\alpha_s,\alpha_t)}$. With the conjugate convention, the vectors $\Theta c_\tau$ are no longer $Sp(\Lambda)$-invariant, and the joint dimensions drop below Catalan, to $1$ and $22$ at $k=4,5$ for $H=\Z_3$, which is impossible inside $S_N^+$: this is the matrix algebra case of Remark~\ref{rem:twcases}. The line $N=4$ is a consistency check: there $PU_2^\sigma=S_4^+$, and the twisted vectors span exactly the noncrossing partitions.

\begin{remark}
These seeds are far from classical, and quite different from the previous ones. $PU_n^\sigma$ is infinite, with the fusion rules of $PU_n$, its fixed point spaces are not spanned by partitions, and already at level $3$ they are larger than those of $S_N$. Its classical part contains the regular subgroup $A$. Nevertheless, everything is absorbed, through level $7$ at $N=9$, level $6$ at $N=16$ and level $5$ at $N=25$. Note also that, by Proposition~\ref{prop:crossing}, a counterexample of this type would disprove two conjectures at once.
\end{remark}

We end with some questions. The main one is whether everything is absorbed at all levels, that is, whether $\langle S_N,PU_n^\sigma\rangle=S_N^+$ for all $n\ge3$, with $N=n^2$, and it is our feeling that this is the case. By Proposition~\ref{prop:crossing} this is equivalent to $\langle S_N^{\bar\sigma},PU_n\rangle=\Aut^+(M_n,\mathrm{tr})$, so it would follow from Conjecture~\ref{conj:mn}, in the liftable case. A possible route is to find inside the generated quantum group a non-classical piece moving few points, as in Proposition~\ref{prop:monloc}, for instance a copy of $S_4^+$ coming from a subgroup $\Z_2\times\Z_2$ of $\Lambda$, when $n$ is even. Computationally, the next steps are level $8$ at $N=9$, level $7$ at $N=16$ and level $6$ at $N=25$, by using the Clifford group $\widehat\Lambda\rtimes Sp(\Lambda)$, which fixes all the twisted fixed vectors that we computed, and which should be the classical part of $PU_n^\sigma$. Also, all the twists above come from bicharacters, and the next step is to use a non-abelian group of central type $G$, with $C^*_\sigma(G)\simeq M_n$, and to identify the Hopf image of the flat model $x\mapsto\Proj(g_ixg_j^*)$ coming from the standard basis of $C^*_\sigma(G)$, as in~\cite[Thm.~2.8]{bn}, which should be a twist of $PU_n$ by a Drinfeld twist supported on $G$. For the smallest such groups see~\cite{bbn,eg,kr}. Finally, Bichon's pseudo-twists, which give the quantum groups $DC_n^\tau$ at $N=4$, see~\cite{bb}, are not twists of the whole tensor category, so they may avoid the locality mechanism, and would be worth testing as seeds.


%======================================================================
% Section: hadamard
%======================================================================

\section{Hadamard models and deformations}\label{sec:hadamard}

\subsection{Hadamard models}\label{sec:hadmodels}

Let $H\in M_N(\T)$ be a complex Hadamard matrix, with rows $H_1,\dots,H_N$. The vectors $\xi_{ij}=H_i/H_j$ form a magic basis of $\C^N$, so the rank-one projections $u_{ij}=\Proj(\xi_{ij})$ form a magic unitary in $M_N(\C)$. We denote by $G_H\subset S_N^+$ the Hopf image of the corresponding representation $C(S_N^+)\to M_N(\C)$, i.e.\ the smallest quantum subgroup through which it factorizes~\cite{bqp,bbs}. For complex Hadamard matrices in general, see~\cite{tz}. Philosophically, $H$ is ``the Fourier matrix of $G_H$''.

We will use these quantum groups $G_H$ as seeds: the question is whether $\langle S_N,G_H\rangle=S_N^+$, as soon as $G_H$ is not classical. The test is computable for two reasons:
\begin{itemize}
\item the intertwiners of the Hopf image are the intertwiners of the model, i.e.\ the vectors with $u^{\otimes k}(\xi\otimes1)=\xi\otimes1$;
\item these are the vectors fixed by the contraction $T_k=(\mathrm{id}\otimes\mathrm{tr})(u^{\otimes k})$.
\end{itemize}
So $\dim\Fix_{\langle S_N,G_H\rangle}(k)$ is the multiplicity of the eigenvalue $1$ of the compression of $T_k$ to $\spn\{\xi_\pi:\pi\in P(k)\}$. The matrix $T_k$ is computed from the Gram matrix of the vectors $\xi_{ij}$, through cyclic traces~\cite[Thm.~6.5]{bbs}. With our convention for scalar products, and with $\xi_r=\xi_{i_rj_r}$, the formula reads
\[
\mathrm{tr}(u_{i_1j_1}\cdots u_{i_kj_k})=\frac{1}{N^{k+1}}\,\langle\xi_2,\xi_1\rangle\langle\xi_3,\xi_2\rangle\cdots\langle\xi_1,\xi_k\rangle .
\]

\begin{theorem}[{\cite[Thm.~6.4]{bbs}}]\label{thm:classical}
For a complex Hadamard matrix $H$, the following are equivalent:
\begin{enumerate}
\item $G_H$ is classical;
\item $G_H$ is a group dual;
\item $H$ is equivalent to a generalized Fourier matrix $F_G=F_{N_1}\otimes\dots\otimes F_{N_k}$, with $G=\Z_{N_1}\times\dots\times\Z_{N_k}$.
\end{enumerate}
In that case $G_H=G$, acting regularly.
\end{theorem}

So Hadamard models never produce non-classical group duals. Their non-classical Hopf images lie outside Section~\ref{sec:duals}.

\begin{remark}[tame and wild matrices]
The correspondence $H\mapsto G_H$ suggests that the wilder the matrix, the bigger the quantum group. For our purposes we need \emph{tame} matrices, whose $G_H$ is non-classical but small. The natural tame matrices are the deformations of Fourier matrices. Deciding in general whether $G_H=S_N^+$ is a hard open problem, and we do not attempt it. We only run the joint test on concrete matrices.
\end{remark}

\subsection{Regular Hadamard matrices at $N\le7$}

\begin{definition}[{\cite{bbs}}]
A complex Hadamard matrix is \emph{regular} if the scalar product of any two distinct rows decomposes as a sum of cycles, i.e.\ of rotated sums of all $p$-th roots of unity, with $p$ prime.
\end{definition}

All the examples below are regular.
\begin{itemize}
\item \emph{$N\le5$.} The only complex Hadamard matrices, up to equivalence, are $F_2$, $F_3$, $F_4^q$ (the Di{\c t}{\u a} deformations of $F_2\otimes F_2$) and $F_5$~\cite[Thm.~5.4]{bbs}.
\item \emph{$N=6$.} The regular complex Hadamard matrices are~\cite[Thm.~11.2]{bbs}:
\begin{itemize}
\item the Tao matrix $T_6$~\cite{tao};
\item the Haagerup family $H_6^q$~\cite{ha};
\item the Di{\c t}{\u a} deformations~\cite{di} $F_{23}^{rs}=F_2\otimes_lF_3$ and $F_{32}^{rs}=F_3\otimes_lF_2$.
\end{itemize}
Here $(h\otimes_lk)_{ia,jb}=h_{ij}l_{aj}k_{ab}$, with a dephased parameter matrix $l$ whose two free entries are $r,s\in\T$.
\item \emph{$N=7$.} The regular complex Hadamard matrices are $F_7$, the Petrescu family $P_7^q$~\cite{pe,szo}, and the family $R_7^q$ of Lampio, {\"O}sterg{\aa}rd and Sz{\"o}ll{\H o}si~\cite{los}. This is a 2026 theorem of GPT and Teo, the paper being in preparation.
\end{itemize}

\begin{remark}[normalizations]
We use the following forms of the matrices above.
\begin{itemize}
\item \emph{The matrices $T_6$ and $H_6^q$.} We use them as printed in~\cite[\S5]{bbs}.
\item \emph{The families at $N=7$.} With $w=e^{2\pi i/6}$, we use the following forms of $P_7^q$ and $R_7^q$:
\[
{\setlength{\arraycolsep}{4pt}
P_7^q=\begin{pmatrix}-q&q&w^2&1&w^2&1&w^2\\q&-q&w^2&1&1&w^2&w^2\\w^2&w^2&-w^2&1&w^2&w^2&1\\1&1&1&-1&w^2&w^2&w^2\\w^2&1&w^2&w^2&-\bar qw^2&\bar qw^2&1\\1&w^2&w^2&w^2&\bar qw^2&-\bar qw^2&1\\w^2&w^2&1&w^2&1&1&-1\end{pmatrix}
\qquad
R_7^q=\begin{pmatrix}1&1&1&1&1&1&1\\1&1&1&w^2&w^3&w^3&w^4\\1&w&w^4&w^4&q&qw^3&w^2\\1&w^2&w^3&1&qw^4&qw&w^4\\1&w^3&w^5&w^3&qw^3&q&w\\1&w^4&w^2&qw^2&q&qw^5&qw^3\\1&w^4&w^2&qw^5&qw^2&qw^3&q\end{pmatrix}}
\]
\end{itemize}
\end{remark}

\subsection{Deformed Fourier matrices reduce to group duals}

Let $A,B$ be finite abelian groups and $Q\in M_{A\times B}(\T)$. Write $G_Q$ for the Hopf image of the Di{\c t}{\u a} deformation $F_A\otimes_QF_B$, and set
\begin{align*}
\Gamma_{A,B}&=B^{*A}\Big/\Big\langle\big[c_1^{(i_1)}\cdots c_s^{(i_s)},\,d_1^{(j_1)}\cdots d_s^{(j_s)}\big]=1\ \Big|\ \textstyle\sum_rc_r=\sum_rd_r=0\Big\rangle\\
&\simeq\Z^{(|A|-1)(|B|-1)}\rtimes B .
\end{align*}
Here the superscripts index the $|A|$ copies of $B$.

\begin{theorem}[Banica--Bichon~\cite{bbrw}]\label{thm:bbrw}
$C(G_Q)$ is a Hopf quotient of the crossed coproduct $C^*(\Gamma_{A,B})\rtimes C(A)$. For generic $Q$ it is isomorphic to it.
\end{theorem}

In other words, $G_Q$ is built from the group dual $\widehat{\Gamma_{A,B}}$ and the classical group $A$, and the deformation parameter $Q$ only decides how much of this survives. What we need is the group-dual part, which survives in all cases, as follows.

\begin{proposition}\label{prop:dual}
For every $Q$, the quantum group $G_Q$ contains a group dual $\widehat{\Gamma_Q}$, where $\Gamma_Q$ is a quotient of $\Gamma_{A,B}$ and $\Gamma_Q=\Gamma_{A,B}$ for generic $Q$. It acts on the $N=|A||B|$ points through the Fourier family of the $|A|$ copies of $B$, in the sense of Section~\ref{sec:duals}.
\end{proposition}

\begin{proof}[Sketch]
In $C^*(\Gamma_{A,B})\rtimes C(A)$ the subalgebra $C(A)$ is central. Dividing the Hopf image $L=C(G_Q)$ by the ideal generated by $C(A)^+$ gives a Hopf quotient, which is a quotient of $C^*(\Gamma_{A,B})$, hence a group algebra $C^*(\Gamma_Q)$ (compare the proof of~\cite[Lemma~4.5]{bbrw}). For generic $Q$, Theorem~\ref{thm:bbrw} gives $\Gamma_Q=\Gamma_{A,B}$. Under the counit of $C(A)$, the magic unitary $w_{ia,jb}=u^{(i)}_{ab}v_{ij}$ of~\cite[Lemma~3.8]{bbrw} becomes $\delta_{ij}u^{(i)}_{ab}$. Each $u^{(i)}$ is the Fourier block of the $i$-th copy of $B$.
\end{proof}

\begin{corollary}\label{cor:fx2}
If $Q$ is generic and $B=\Z_2$, then $\langle S_N,G_Q\rangle=S_N^+$. In particular this holds for $F_{32}^{rs}$ with generic $r,s$.
\end{corollary}

\begin{proof}
Kill all the copies of $B$ except two in $\Gamma_{A,\Z_2}$. The quotient is $\Gamma_{\Z_2,\Z_2}=\Z\rtimes\Z_2=D_\infty$, acting through two $\Z_2$-blocks, i.e.\ on $4$ points. It is non-abelian, so locality (Proposition~\ref{prop:monloc}, together with~\cite[Cor.~3.4]{bcf}) and monotonicity give the result.
\end{proof}

\begin{remark}
Two comments on the cases not covered by Corollary~\ref{cor:fx2}.
\begin{itemize}
\item For $F_{23}^{rs}$ with generic parameters, $\Gamma_{\Z_2,\Z_3}=\Z^2\rtimes\Z_3$ acts through two $\Z_3$-blocks on $6$ points, so locality does not apply. It has $A_4=\Z_2^2\rtimes\Z_3$ as a quotient. So $\langle S_6,G_Q\rangle$ contains the candidate of $\widehat{A_4}$ with two $\Z_3$-blocks, which is Catalan through level $8$ (Proposition~\ref{prop:dualcomp}).
\item For special $Q$, the same reduction works as long as $\Gamma_Q$ remains non-abelian.
\end{itemize}
\end{remark}

\subsection{Computations}

\begin{proposition}[computational]\label{prop:hadres}
The joint dimensions $\dim\Fix_{\langle S_N,G_H\rangle}(k)$ are as follows.
\begin{center}\small
\setlength{\tabcolsep}{4pt}
\begin{tabular}{@{}llll@{}}\toprule
$H$ & parameters & $\dim\Fix(G_H)(2),(3)$ & joint, levels\\\midrule
$T_6$ & -- & $2,5$ & Catalan, $k\le5$\\
$H_6^q$ & $q=1,i,e^{2\pi i/3},e^{i\pi/3}$, generic & $2,5$ & Catalan, $k\le5$\\
$F_{23}^{rs},F_{32}^{rs}$ & $r,s$ all $12$th roots of unity & & $k=4$: see below\\
$F_{23}^{rs},F_{32}^{rs}$ & generic, and $5$ root-of-unity pairs & $4,18$ & Catalan, $k\le5$\\
$P_7^q$ & $9$ values, incl.\ roots of unity & $2,5$ & Catalan, $k\le5$\\
$R_7^q$ & $q=\pm1,i,e^{2\pi i/3},e^{i\pi/3}$, generic & $2,5$ & Catalan, $k\le5$\\\bottomrule
\end{tabular}
\end{center}
For the Di{\c t}{\u a} deformations at $k=4$, $132$ of the $144$ root-of-unity pairs give Catalan ($14$), for both families. The remaining $12$ give Bell ($15$). At those points $\dim\Fix(G_H)(2),(3)=6,36$, which is the regular action of $\Z_6$: the deformation is Fourier-equivalent there, and $G_H$ is classical, in accordance with Theorem~\ref{thm:classical}. All these numbers were computed in floating point, from the closed-walk formula of Proposition~\ref{prop:walks} below (\texttt{hadamard.py}).
\end{proposition}

\begin{remark}
The Hopf images split into two kinds:
\begin{itemize}
\item \emph{Wild.} $T_6$, $H_6^q$, $P_7^q$ and $R_7^q$ have $\dim\Fix(G_H)$ equal to Catalan already at levels $2$ and $3$, and in fact at all the levels we could reach (Section~\ref{sec:paley}). They look like $S_N^+$ themselves, as the tame/wild heuristic predicts, and for them the joint test is automatic.
\item \emph{Tame.} The non-Fourier Di{\c t}{\u a} points have small Hopf images, with $\dim\Fix(G_H)(2),(3)=4,18$. These are genuinely useful seeds, and they are still absorbed.
\end{itemize}
\end{remark}

\subsection{Wildness and the Paley matrices}\label{sec:paley}

The joint test has content only at the levels where $G_H$ itself has more fixed vectors than $S_N^+$. Indeed, $\mathcal C_{\langle S_N,G_H\rangle}=\mathcal C_{S_N}\cap\mathcal C_{G_H}$, and the noncrossing partition vectors are fixed by both. So if $\dim\Fix(G_H)(k)=C_k$, then $\dim\Fix_{\langle S_N,G_H\rangle}(k)=C_k$ as well, whatever $S_N$ does.

\begin{definition}
A complex Hadamard matrix $H\in M_N(\T)$, with $N\ge4$, is \emph{wild at level $k$} if $\dim\Fix(G_H)(k)=C_k$, and \emph{wild} if it is wild at all levels.
\end{definition}

Since $u=\bar u$, the fixed spaces at all levels determine the category of $G_H$, so $H$ is wild if and only if $G_H=S_N^+$. A wild matrix is useless as a seed: it is a matrix model of $S_N^+$ itself. Wildness at level $k$ can be tested on $G_H$ alone, without $S_N$, by computing the fixed space of the operator $T_k$ of Subsection~\ref{sec:hadmodels}, namely
\[
T_k=(\mathrm{id}\otimes\mathrm{tr})\pi_H(u^{\otimes k})=\Big(\mathrm{tr}(u_{i_1j_1}\cdots u_{i_kj_k})\Big)_{i,j}\in M_{N^k}(\C),
\]
whose fixed vectors are those of $G_H$ at level $k$. This matrix has $N^{2k}$ entries, so it cannot be stored beyond small $k$, but there is no need to, thanks to the following formula.

\begin{proposition}\label{prop:walks}
Let $H\in M_N(\T)$ be complex Hadamard, and let $a_{yx}\in\C^N$ be the vectors $a_{yx}(i)=H_{iy}\bar H_{ix}$. Then
\[
T_k=\frac{1}{N^{k+1}}\sum_{x_1,\dots,x_k=1}^N a_{x_kx_1}a_{x_kx_1}^*\otimes a_{x_1x_2}a_{x_1x_2}^*\otimes\dots\otimes a_{x_{k-1}x_k}a_{x_{k-1}x_k}^*,
\]
a sum over the closed walks $x_1\to x_2\to\dots\to x_k\to x_1$ on the columns of $H$. Consequently:
\begin{enumerate}
\item $T_k$ is a positive contraction;
\item the orthogonal complement $W_k$ of $\spn\{\xi_\pi:\pi\in NC(k)\}$ is invariant under $T_k$, and $H$ is wild at level $k$ if and only if $\lambda_k(H)=\max\sigma(T_k|_{W_k})<1$;
\item $\lambda_k(H)\le\lambda_{k+1}(H)$ for all $k$.
\end{enumerate}
\end{proposition}

\begin{proof}
Since $u_{ij}=\Proj(\xi_{ij})$ with $\|\xi_{ij}\|^2=N$, the trace of a product of rank-one projections gives
\[
\mathrm{tr}(u_{i_1j_1}\cdots u_{i_kj_k})=\frac{1}{N^{k+1}}\prod_{r=1}^k\sum_{x_r=1}^N\overline{\xi_{i_rj_r}(x_r)}\,\xi_{i_{r+1}j_{r+1}}(x_r),
\]
with indices $r$ taken modulo $k$. Expand $\xi_{ij}(x)=H_{ix}\bar H_{jx}$. The slot $r$ appears in the factors $r-1$ and $r$, and contributes $H_{i_rx_{r-1}}\bar H_{i_rx_r}\cdot\overline{H_{j_rx_{r-1}}\bar H_{j_rx_r}}=a_{x_{r-1}x_r}(i_r)\overline{a_{x_{r-1}x_r}(j_r)}$, which is the formula. Each term is a tensor product of positive rank-one operators, so $T_k\ge0$, and $\|T_k\|\le1$ because $T_k=(\mathrm{id}\otimes\mathrm{tr})(U)$ with $U$ the image of the unitary $u^{\otimes k}$. The noncrossing vectors are fixed, because $G_H\subset S_N^+$, so their span is invariant, and so is $W_k$, since $T_k$ is self-adjoint. For a positive contraction, the fixed vectors are the eigenvectors at the top eigenvalue $1$, which gives (2). For (3), let $v\in W_k$ and $v'=v\otimes\sum_ie_i$. Since $\sum_ju_{ij}=1$, we have $T_{k+1}v'=(T_kv)\otimes\sum_ie_i$. Also $v'\in W_{k+1}$, because summing $\xi_\pi$, $\pi\in NC(k+1)$, over its last index gives a multiple of $\xi_{\pi'}$, with $\pi'\in NC(k)$ the restriction of $\pi$. So the Rayleigh quotients of $v$ and $v'$ agree.
\end{proof}

Computationally, $T_kv$ is obtained by a transfer along the walk, at a cost of order $kN^{k+3}$, and $\lambda_k(H)$ by the Lanczos method on $W_k$. The number $1-\lambda_k(H)$ is a spectral gap, measuring how wild $H$ is at level $k$. For the Fourier-equivalent and the tame Di{\c t}{\u a} points of Proposition~\ref{prop:hadres} we find $\lambda_2=1$, as we should.

Real Hadamard matrices are a natural place to look, and they tend to be forgotten in the complex theory. The basic constructions are those of Paley~\cite{pa}, see also~\cite{bha}. Let $q$ be an odd prime power, $\chi$ the quadratic character of $\F_q$, with $\chi(0)=0$, and $Q=(\chi(b-a))_{a,b\in\F_q}$ the Jacobsthal matrix, which satisfies
\[
QQ^t=q1_q-J,\qquad QJ=JQ=0,\qquad Q^t=\chi(-1)Q,
\]
with $J$ the all-one matrix, these identities being what makes the matrices below Hadamard. With $\mathbf 1$ the all-one column vector, the Paley I matrix, for $q\equiv3\pmod4$, and the Paley II matrix, for $q\equiv1\pmod4$, are
\[
H=1_{q+1}+\begin{pmatrix}0&\mathbf 1^t\\-\mathbf 1&Q\end{pmatrix},\qquad
H=C\otimes\begin{pmatrix}1&1\\1&-1\end{pmatrix}+1_{q+1}\otimes\begin{pmatrix}1&-1\\-1&-1\end{pmatrix},\quad C=\begin{pmatrix}0&\mathbf 1^t\\\mathbf 1&Q\end{pmatrix},
\]
of orders $N=q+1$ and $N=2(q+1)$ respectively.
For $q\equiv3\pmod4$ the matrix $Q$ is antisymmetric and Paley I is skew-Hadamard, while for $q\equiv1\pmod4$ the matrix $C$ is a symmetric conference matrix. The small cases are as follows:
\begin{itemize}
\item $N=4,8$: Paley I with $q=3,7$ gives the Walsh matrices $F_2^{\otimes2},F_2^{\otimes3}$, the only real Hadamard matrices of these orders, and $G_H$ is classical.
\item $N=12$: Paley I with $q=11$ and Paley II with $q=5$ are equivalent, the real Hadamard matrix of order $12$ being unique.
\item $N=20$: Paley I with $q=19$ and Paley II with $q=9$ are inequivalent, two of the three equivalence classes found by Hall~\cite{hal}.
\item $N=24$: Paley I with $q=23$. $N=28$: Paley I with $q=27$ and Paley II with $q=13$, inequivalent (their $4$-row profiles differ). The next orders, $N=32$ (Paley I, $q=31$) and $N=36$ (Paley II, $q=17$), are beyond our reach.
\end{itemize}
For a real $H$ the vectors $\xi_{ij}=H_iH_j$ are real and symmetric in $i,j$, so the model satisfies $u_{ij}=u_{ji}$, which means that the antipode fixes the generators, $S(u_{ij})=u_{ij}$, a feature without analogue for genuinely complex matrices. Whether it says anything about $G_H$ is not clear to us. We believe that the Paley matrices of order $N\ge12$ are wild, i.e.\ $G_H=S_N^+$, and the evidence is as follows. Note that at every level where $H$ is wild the joint test with $S_N$ is empty, so for these matrices the numbers $\lambda_k(H)$ are the whole story, and the question is how far from $1$ they stay when $k$ grows. The cost of the walk transfer, of order $kN^{k+3}$, limits us to $k\le7$ at $N=6,7$, to $k\le6$ at $N=12$, and to $k\le4$ at $N=20,24,28$.

\begin{proposition}[numerical]\label{prop:wildnum}
The values of $\lambda_k(H)$, rounded to three decimals, are as follows. All of them are $<1$, so these matrices are wild at all the levels shown.
\begin{center}\small
\setlength{\tabcolsep}{5pt}
\begin{tabular}{@{}llcccccc@{}}\toprule
$H$ & parameters & $k=2$ & $3$ & $4$ & $5$ & $6$ & $7$\\\midrule
$T_6$ & -- & $0.250$ & $0.428$ & $0.619$ & $0.619$ & $0.651$ & $0.651$\\
$H_6^q$ & $q=1,i$ & $0.370$ & $0.370$ & $0.667$ & $0.667$ & $0.667$ & $0.667$\\
 & $q=e^{2\pi i/3},e^{i\pi/3}$ & $0.499$ & $0.499$ & $0.665$ & $0.665$ & $0.694$ & \\
 & generic & $0.518$ & $0.518$ & $0.641$ & $0.641$ & $0.666$ & $0.666$\\
$P_7^q$ & $q=1,e^{2\pi i/3}$ & $0.353$ & $0.570$ & $0.570$ & $0.570$ & $0.571$ & \\
 & $q=i$ & $0.353$ & $0.503$ & $0.562$ & $0.593$ & $0.593$ & \\
 & generic & $0.353$ & $0.551$ & $0.551$ & $0.570$ & $0.591$ & $0.591$\\
$R_7^q$ & $q=\pm1,e^{2\pi i/3}$ & $0.353$ & $0.570$ & $0.570$ & $0.570$ & $0.571$ & \\
 & $q=i$ & $0.353$ & $0.503$ & $0.562$ & $0.593$ & $0.593$ & \\
 & generic & $0.353$ & $0.555$ & $0.555$ & $0.565$ & $0.588$ & $0.588$\\\bottomrule
\end{tabular}

\medskip
\begin{tabular}{@{}lccccc@{}}\toprule
Paley matrix & $k=2$ & $3$ & $4$ & $5$ & $6$\\\midrule
$N=12$, I ($q=11$) and II ($q=5$) & $0.185$ & $0.309$ & $0.399$ & $0.399$ & $0.433$\\
$N=20$, I ($q=19$) & $0.292$ & $0.292$ & $0.292$ & & \\
$N=20$, II ($q=9$) & $0.392$ & $0.462$ & $0.508$ & & \\
$N=24$, I ($q=23$) & $0.207$ & $0.207$ & $0.240$ & & \\
$N=28$, I ($q=27$) & $0.200$ & $0.200$ & $0.200$ & & \\
$N=28$, II ($q=13$) & $0.528$ & $0.570$ & $0.570$ & & \\\bottomrule
\end{tabular}
\end{center}
Here ``generic'' means $q=e^{2\pi it}$ with $t=0.1234$ for $H_6^q$ and $P_7^q$, and $t=0.3711$ for $R_7^q$. The eigenvalues were computed by the Lanczos method on $W_k$, with relative tolerance $10^{-6}$, applying $T_k$ as a cyclic chain of $4$-index tensors, $(T_k)_{ij}=N^{-k-1}\prod_sQ(i_s,j_s,i_{s+1},j_{s+1})$ with $Q(a,b,c,d)=\sum_x\bar H_{ax}H_{bx}H_{cx}\bar H_{dx}$ (\texttt{hadamard.py}), and the noncrossing vectors were checked to be fixed by $T_k$ up to $2\cdot10^{-15}$. At $N=12$ the two Paley matrices give the same values, as they should.
\end{proposition}

Several features are visible. The values are nondecreasing in $k$, as they must be by Proposition~\ref{prop:walks}, and they often stay constant for several levels. The Paley I matrices are much wilder than the complex matrices at $N=6,7$, and get wilder as $N$ grows. The Paley II matrices are less wild than Paley I at the same $N$, and comparable to the matrices at $N=6,7$. None of these matrices is a useful seed: at all the levels shown, $\langle S_N,G_H\rangle$ is Catalan for free. As explained above, it is our feeling that the Paley matrices are wild, and this leads to the following question.

\begin{question}\label{qu:paley}
Prove that the Paley matrices of order $N\ge12$ are wild, i.e.\ that $G_H=S_N^+$, or at least that $\lambda_k(H)<1$ for all $k$.
\end{question}

The values above are often constant in $k$ over several levels, for instance $\lambda_k=0.292$ for $k=2,3,4$ at $N=20$, Paley I, and at $N=12$ they seem to be $\lambda_2=5/27$, $\lambda_3=25/81$, $\lambda_4=\lambda_5=97/243$, as rational guesses from the decimals. So an exact formula for $\lambda_k$ might be within reach, using the large automorphism groups of these matrices.

\subsection{Conclusion}

\begin{theorem}\label{thm:hadmain}
Let $H$ be a regular complex Hadamard matrix with $N\le6$, or one of $F_7$, $P_7^q$, $R_7^q$. Then exactly one of the following holds:
\begin{enumerate}
\item $H$ is Fourier-equivalent, $G_H$ is classical, and $\langle S_N,G_H\rangle=S_N$;
\item $\langle S_N,G_H\rangle=S_N^+$.
\end{enumerate}
The statement is proved for $N\le5$, by maximality at $N=4,5$, and for $F_{32}^{rs}$ with generic parameters (Corollary~\ref{cor:fx2}). In the remaining cases it is verified through level $5$, on the parameter values of Proposition~\ref{prop:hadres}, and through level $6$ for $T_6$, $H_6^q$, $P_7^q$, $R_7^q$ on the parameter values of Proposition~\ref{prop:wildnum}, where these matrices are wild.
\end{theorem}

We end with a question about the joint test itself. Let $M_k(H)$ be the matrix with entries $\langle T_k\xi_\sigma,\xi_\pi\rangle$, for $\pi,\sigma\in P(k)$, that is, the compression of $T_k$ to the span of the partition vectors. Its entries are the values of the model state on the $S_N$-bi-invariant sums $\sum u_{i_1j_1}\cdots u_{i_kj_k}$, with $i,j$ constant on the blocks of $\pi,\sigma$. Since $T_k$ is a positive contraction, $\dim\Fix_{\langle S_N,G_H\rangle}(k)$ is the corank of $G_k-M_k$, with $G_k$ the Gram matrix of the partition vectors. So the joint test only sees $M_k$, and if $M_k(H^q)$ does not depend on $q$, then a whole family is a single test case.

\begin{question}\label{qu:invis}
For which families $H^q$ of complex Hadamard matrices is $M_k(H^q)$ independent of $q$, and why?
\end{question}

Here is what we know (computational). At $k\le3$, the matrix $M_k$ is the same for all complex Hadamard matrices of a given size, as checked at $N=6$ for $T_6$, $H_6^q$, Di{\c t}{\u a} and $F_6$. At $k=4$, the matrices $M_4(P_7^q)$ and $M_4(R_7^q)$ are independent of $q$, and equal to each other, recall that $P_7^q$ is equivalent to $\bar R_7^q$, but they differ from $M_4(F_7)$. On the other hand, $M_4(H_6^q)$ depends on $q$, and so does $M_4$ for the Di{\c t}{\u a} families, where the joint dimension at $k=4$ jumps from $14$ to $15$ at the Fourier points. Also, the numbers $\lambda_k(P_7^q)$ do depend on $q$, so it is only the compression which is invisible. Is $M_k(P_7^q)$ independent of $q$ at all levels $k$? What distinguishes $P_7^q$ from $H_6^q$ here, perhaps the absence of a Fourier point on the family? Is $M_k$ determined by some finite set of invariants of $H$, such as the multisets $\{\sum_iH_{iy_1}\bar H_{ix_1}\cdots H_{iy_s}\bar H_{ix_s}\}$ along closed walks, as the closed-walk formula of Proposition~\ref{prop:walks} suggests?



%======================================================================
% Section: exotic
%======================================================================

\section{Exotic quantum permutation groups}\label{sec:exotic}

Throughout this section, $G$ denotes a hypothetical \emph{exotic} quantum permutation group, $S_N\subsetneq G\subsetneq S_N^+$. By maximality at $N\le5$, such a $G$ can only exist for $N\ge6$. We collect what can be said about it, in soft terms: representation theory, orbitals, subgroups, diagonal algebras, and analysis.

\subsection{Representation theory and moments}

By Tannakian duality, $G$ corresponds to a tensor category $\mathcal C_G$ with
\[
\spn(NC)\subsetneq\mathcal C_G\subsetneq\spn(P).
\]
In particular $C_k\le\dim\Fix_G(k)\le\dim\Fix_{S_N}(k)$ for all $k$, where $C_k$ is the Catalan number. Moreover $\mathcal C_G$ is not spanned by partitions: there is no easy quantum group strictly between $S_N$ and $S_N^+$~\cite{bsp}. So $G$ is non-easy, and $\mathcal C_G$ must contain a linear combination of crossing partitions that is not itself a partition~\cite{mc2}.

\begin{proposition}[McCarthy~\cite{mc1}]\label{prop:mc}
The Haar state of $G$ is tracial and invariant under the antipode and under relabeling. All its integrals of degree $\le4$ in the coordinates $u_{ij}$ coincide with those of $S_N^+$. In particular the fourth moment of the character is $14$, i.e.\ $\dim\Fix_G(k)=C_k$ for $k\le4$.
\end{proposition}

So an exotic $G$ is invisible to $S_N^+$ through level $4$. The level computations of Sections~\ref{sec:method}--\ref{sec:higher} push this further.
McCarthy also shows that no exotic category is generated by a linear combination of three crossing partitions~\cite{mc2}.

\subsection{Orbitals}

Following Lupini, Man{\v c}inska and Roberson~\cite{lmr}, write $(i_1,\dots,i_m)\sim_m(j_1,\dots,j_m)$ when $u_{i_1j_1}\cdots u_{i_mj_m}\ne0$. Say that $G$ has \emph{free $m$-orbitals} if such a monomial vanishes only for the trivial reason, namely two consecutive factors in a common row or column.

\begin{theorem}[McCarthy~\cite{mc1}]\label{thm:orbitals}
An exotic $G$ has free orbits, free $2$-orbitals and free $3$-orbitals.
\end{theorem}

\begin{remark}
Two comments on this result.
\begin{itemize}
\item Freeness of orbits and $2$-orbitals is inherited from $S_N$, through abelianization. Freeness of $3$-orbitals is the main result of~\cite{mc1}.
\item McCarthy asks whether \emph{any} proper quantum subgroup of $S_N^+$ can have free $3$-orbitals. A negative answer would prove maximality.
\end{itemize}
\end{remark}

\subsection{What $G$ cannot contain}

The subgroup approach of Sections~\ref{sec:duals}--\ref{sec:hadamard} turns into a list of properties of $G$, through the following trivial observation.

\begin{lemma}\label{lem:contain}
If $K\subset G$, then $\langle S_N,K\rangle\subset G$. In particular:
\begin{itemize}
\item if $\langle S_N,K\rangle=S_N^+$, then $K\not\subset G$;
\item if $\langle S_N,K\rangle$ agrees with $S_N^+$ through level $k$, and $K\subset G$, then so does $G$.
\end{itemize}
\end{lemma}

For $M\le N$, view $S_M^+\subset S_N^+$ as acting on $M$ chosen points and fixing the others, i.e.\ $v\mapsto\mathrm{diag}(v,1_{N-M})$. Define $G_M=G\cap S_M^+$, the intersection of quantum subgroups of $S_N^+$. It is the largest quantum subgroup of $G$ that only moves these $M$ points.

\begin{proposition}[restriction]\label{prop:restriction}
For $4\le M\le N$, either $G_M=S_M$, or $G_M$ is exotic at $M$, i.e.\ $S_M\subsetneq G_M\subsetneq S_M^+$. In particular:
\begin{enumerate}
\item if $S_M\subset S_M^+$ is maximal, then $G_M=S_M$;
\item every quantum subgroup of $G$ that moves at most $5$ points is classical;
\item more generally, if maximality holds at all $M\le m$, then every quantum subgroup of $G$ moving at most $m$ points is classical.
\end{enumerate}
\end{proposition}

\begin{proof}
We have $S_M\subset S_N\subset G$ and $S_M\subset S_M^+$, so that we have inclusions as follows:
\[
S_M\subset G_M\subset S_M^+ .
\]
If $G_M=S_M^+$, then $S_M^+\subset G$. Since $4\le M\le N$, the result in~\cite[Cor.~3.4]{bcf} gives
\[
S_N^+=\langle S_N,S_M^+\rangle\subset G ,
\]
contradicting $G\ne S_N^+$. This proves the dichotomy, and (1) follows.

For (2) and (3), let $K\subset G$ move at most $m$ points. Conjugating by an element of $S_N\subset G$, we may assume these are the first $m$ points. Then $K\subset G_m$, and $G_m=S_m$ by (1), with the cases $m\le3$ being trivial. Maximality at $4$ and $5$ is known~\cite{bah,bb}.
\end{proof}

So exotic quantum permutations are \emph{global}: they are invisible on every set of at most $5$ points, and on every set of size $M$ at which maximality holds. Inductively, the first exotic $G$, if any, restricts to $S_M$ on every proper subset.

\begin{corollary}\label{cor:nocontain}
An exotic $G$ contains none of the following:
\begin{itemize}
\item a non-classical quantum subgroup moving $\le5$ points, for instance $S_4^\tau$, or the dual of $\Z_2*\Z_2$ or of $\Z_2*\Z_3$ on $4$ or $5$ points;
\item a diagonal copy of $S_M^+$ with $M\ge5$~\cite[Ex.~3.9, Prop.~3.10]{bcf};
\item $G^+(X)$ for a graph $X$ with two disjoint pairs of twins, or more generally with two disjoint automorphisms of total support $\le5$ (Corollary~\ref{cor:small});
\item $G_H$ for a generic Di{\c t}{\u a} deformation $F_A\otimes_QF_2$ (Corollary~\ref{cor:fx2});
\item a group-theoretical quantum group $H_M^{[\Gamma]}\subset H_M^+\subset S_{2M}^+$ of Raum and Weber~\cite{rw2,rw}, with $N=2M$, as soon as the quotient of $\Gamma$ by the normal subgroup generated by $g_3,\dots,g_M$ is non-abelian, for instance $H_M^{[\infty]}$, where $\Gamma=\Z_2^{*M}$. Indeed, $H_M^{[\Gamma]}$ contains the dual of $\Gamma$, acting diagonally through $u_{ij}\mapsto\delta_{ij}g_i$, and hence the dual of this quotient, a non-abelian quotient of $\Z_2*\Z_2$, which acts on $4$ of the $2M$ points.
\end{itemize}
If $G$ contains any of the other seeds of Sections~\ref{sec:duals}--\ref{sec:hadamard}, then $G$ agrees with $S_N^+$ through the level at which that seed was checked (Lemma~\ref{lem:contain}).
\end{corollary}

\subsection{Relabeling and diagonal algebras}\label{sec:diagonal}

Since $S_N\subset G$, every permutation $\sigma\in S_N$ gives a character of $C(G)$, namely $ev_\sigma(u_{ij})=\delta_{i\sigma(j)}$, the evaluation at the permutation matrix of $\sigma$. These characters can be used to move the indices around, inside $C(G)$ itself.

\begin{proposition}[relabeling]\label{prop:relabel}
Assume $S_N\subset G\subset S_N^+$. For any $\sigma,\tau\in S_N$ there is a $*$-automorphism $\alpha_{\sigma\tau}$ of $C(G)$ such that
\[
\alpha_{\sigma\tau}(u_{ij})=u_{\sigma(i)\tau(j)}.
\]
\end{proposition}

\begin{proof}
Set $\alpha_{\sigma\tau}=(ev_{\sigma^{-1}}\otimes\mathrm{id}\otimes ev_\tau)\Delta^{(2)}$, where $\Delta^{(2)}=(\Delta\otimes\mathrm{id})\Delta$. This is a composition of $*$-homomorphisms, and
\[
\alpha_{\sigma\tau}(u_{ij})=\sum_{k,l}ev_{\sigma^{-1}}(u_{ik})\,u_{kl}\,ev_\tau(u_{lj})=u_{\sigma(i)\tau(j)}.
\]
Since $\alpha_{\sigma^{-1}\tau^{-1}}$ is an inverse, $\alpha_{\sigma\tau}$ is an automorphism.
\end{proof}

This is the mechanism behind McCarthy's invariance of the Haar state under relabeling (Proposition~\ref{prop:mc}), and it has the following consequences, from~\cite[Props.~10.38, 10.39]{bqp}, where they are obtained through the Haar state and Hahn--Banach.

\begin{corollary}\label{cor:relabel}
Assume $S_N\subset G\subset S_N^+$.
\begin{enumerate}
\item The algebra $\langle u_{i_1j_1},\dots,u_{i_pj_p}\rangle\subset C(G)$ depends, up to isomorphism, only on $\ker i$ and $\ker j$, that is, on the relative position of the pairs $(i_r,j_r)$ inside the square $\{1,\dots,N\}^2$.
\item If $G$ is not classical, then $[u_{ij},u_{kl}]\ne0$ for all $i\ne k$ and $j\ne l$.
\item The diagonal algebra $D(G)=\langle u_{11},\dots,u_{NN}\rangle$, regarded as a quotient of $C^*(\Z_2^{*N})$ through $g_i\mapsto2u_{ii}-1$, has an $S_N$-invariant kernel.
\end{enumerate}
\end{corollary}

\begin{proof}
(1) The automorphisms $\alpha_{\sigma\tau}$ act transitively on the pairs of multi-indices with given kernels. (2) If $i=k$ or $j=l$, then $u_{ij}$ and $u_{kl}$ commute, being equal or orthogonal projections. So, $G$ being non-classical, some commutator with $i\ne k$, $j\ne l$ is nonzero, and then all of them are, by (1). (3) The automorphisms $\alpha_{\sigma\sigma}$ map $u_{ii}\mapsto u_{\sigma(i)\sigma(i)}$.
\end{proof}

At the level of diagonal algebras we have quotient maps $D(S_N^+)\to D(G)\to D(S_N)$, all of them $S_N$-equivariant. The last one is commutative and explicit: $u_{ii}$ is the indicator of the permutations fixing $i$, so $D(S_N)=C(X_N)$, where $X_N$ is the set of fixed-point sets of permutations, namely the subsets of $\{1,\dots,N\}$ of size $\ne N-1$. So $D(G)$ is an $S_N$-invariant quotient of $D(S_N^+)$ which still projects onto $C(X_N)$. Very little is known about $D(S_N^+)$ itself, and for this we refer to McCarthy~\cite{mc0,mc1,mc2}.

As explained in~\cite[Chapter~8]{bqp}, a promising approach to the diagonal algebras of the various quantum permutation groups $S_N\subset G\subset S_N^+$ is expected to come from looking more generally at the diagonal algebras of quantum semigroups of partial permutations, in the spirit of the constructions of So{\l}tan~\cite{so}. Whether this idea can really lead to concrete results remains to be seen.

\begin{question}\label{qu:diagonal}
Understand the diagonal algebras $D(S_N^+)$ and $D(G)$, for $S_N\subset G\subset S_N^+$, and find out whether they can detect an exotic $G$.
\end{question}

Some comments on this question.
\begin{itemize}
\item Conjecture~10.43 in~\cite{bqp} states that the $u_{ii}$ are free in $C(S_N^+)$, so that $D(S_N^+)\simeq C^*(\Z_2^{*N})$. This seems to fail at $N=4$: in the Pauli model one finds relations of type $p_ap_bp_c=p_ap_dp_c$, and polynomial growth, which is forced by subhomogeneity. This should be rechecked, and for $N\ge5$ the question is open.
\item The free wreath products $\Z_2\wr_*G$ do not see $G$, their diagonal subgroup being always $\widehat{\Z_2^{*N}}$. The $G$-dependent wreath product $\langle G,\Z_2^N\rangle\subset H_N^+$ interpolates between $H_N$ and $H_N^+$, but easy inputs give nothing here, by~\cite{bcs}.
\item As suggested in~\cite{bqp}, $D(G)$ should come from an intermediate quotient $\Z_2^{*N}\to\Gamma\to\Z_2^N$. However, $D(S_N)=C(X_N)$ is a quotient of $C(\Z_2^N)$ which is not a group algebra, so the right intermediate objects are rather the $S_N$-invariant quotients of $C^*(\Z_2^{*N})$, not necessarily group algebras.
\end{itemize}

\subsection{Analysis, growth and Lie theory}

Beyond representation theory, one can ask about the analytic and geometric properties of an exotic $G$: coamenability, growth, and Lie-theoretic aspects. Regarding growth, the natural questions are those of polynomial growth of the dimensions of the irreducible representations, and of the related finiteness conditions, as in~\cite{cdpr,dpr}, polynomial growth implying coamenability. Regarding Lie theory, this is only available, to some extent, in the half-classical case, see~\cite{bv,bdv}, and its adaptation to exotic quantum permutations is an open question.

We focus in what follows on coamenability. For a compact quantum group of Kac type with self-adjoint fundamental representation, coamenability is equivalent to $N$ lying in the support of the law $\mu_G$ of the character $\chi=\sum_iu_{ii}$ (a Kesten-type criterion). Since $0\le\chi\le N$ and $\int_G\chi^k=\dim\Fix_G(k)$, this means
\[
\limsup_{k\to\infty}\,(\dim\Fix_G(k))^{1/k}=N.
\]

\begin{lemma}\label{lem:coamsub}
Let $H\subset G\subset S_N^+$ be closed quantum subgroups. If $G$ is coamenable, then so is $H$.
\end{lemma}

\begin{proof}
We have $\Fix_G(k)\subset\Fix_H(k)$, so the moments of $\mu_H$ dominate those of $\mu_G$. Both growth rates are at most $N$, so if the rate of $G$ is $N$, so is the rate of $H$.
\end{proof}

So coamenability goes down to quantum subgroups, and a non-coamenable quantum subgroup is an obstruction. Let us compare with the two extreme cases.
\begin{itemize}
\item $S_N$ is coamenable, being a finite group.
\item $S_N^+$ is not, for $N\ge5$. The usual witness is the non-amenable group dual $\widehat{\Z_2*\Z_3}$ acting on $5$ points, together with Lemma~\ref{lem:coamsub}.
\item By Proposition~\ref{prop:restriction}, that witness cannot sit inside an exotic $G$. Nor can any non-classical subgroup on $\le5$ points.
\end{itemize}

There is one case where coamenability is automatic, namely when $G$ is finite, i.e.\ when $C(G)$ is finite dimensional. This case is detected by the law of the character.

\begin{proposition}\label{prop:atom}
Let $G\subset S_N^+$. If $G$ is finite, then $\mu_G(\{N\})=1/\dim C(G)$. If $G$ is infinite, then $\mu_G(\{N\})=0$.
\end{proposition}

\begin{proof}
Let $P$ be the spectral projection of $\chi$ at $N$, in the von Neumann algebra $L^\infty(G)$ of the Haar state $h$, so that $\mu_G(\{N\})=h(P)$. From $\chi P=NP$ we get $\sum_iP(1-u_{ii})P=0$, a sum of positive elements, so $(1-u_{ii})P=0$, that is, $u_{ii}P=P=Pu_{ii}$ for all $i$. For $i\ne j$ we get $u_{ij}P=u_{ij}u_{jj}P=0$, and similarly $Pu_{ij}=0$. Thus $xP=Px=\varepsilon(x)P$ for every $x$ in the dense $*$-algebra $\mathcal O(G)$, where $\varepsilon$ is the counit.

Now expand $P$ in $L^2(G)$, in the orthonormal basis $\sqrt{d_\alpha}\,v^\alpha_{ij}$ of coefficients of irreducible representations, $G$ being of Kac type. The coefficient of $P$ on $\sqrt{d_\alpha}\,v^\alpha_{ij}$ is $\langle P,\sqrt{d_\alpha}\,v^\alpha_{ij}\rangle=\sqrt{d_\alpha}\,h((v^\alpha_{ij})^*P)=\sqrt{d_\alpha}\,\delta_{ij}\,h(P)$, and Parseval gives
\[
h(P)=h(P^*P)=\sum_\alpha d_\alpha^2\,h(P)^2.
\]
So either $h(P)=0$, or $\sum_\alpha d_\alpha^2=1/h(P)$ is finite, in which case $G$ is finite, with $\dim C(G)=\sum_\alpha d_\alpha^2=1/h(P)$. Conversely, if $G$ is finite, the counit is a character of the finite dimensional algebra $C(G)$, so it is supported by a nonzero central projection $p$, with $xp=\varepsilon(x)p$. Then $\chi p=Np$, so $p\le P$, and $h(P)\ge h(p)>0$, the Haar state being faithful on $C(G)$. So $h(P)=1/\dim C(G)$.
\end{proof}

So an exotic $G$ is either finite, and then coamenable, with $\dim C(G)>N!$ since $C(G)$ maps onto $C(S_N)$, or infinite, with no atom of $\mu_G$ at $N$. In the infinite case, coamenability means that $N$ is an accumulation point of the support of $\mu_G$.

\begin{question}\label{qu:coam}
Is an exotic $G$ coamenable? Can an exotic $G$ be finite?
\end{question}

Our feeling is that an exotic $G$, if any, should rather not be coamenable. Indeed, the results of McCarthy on the Haar state and the orbitals (Proposition~\ref{prop:mc} and Theorem~\ref{thm:orbitals}) suggest that such a $G$ is ``closer'' to $S_N^+$ than to $S_N$, and the level computations of Sections~\ref{sec:method}--\ref{sec:higher} point in the same direction: for $8\le N\le37$, the first $8$ moments of $\mu_G$ are the Catalan numbers, as for $S_N^+$. This does not decide the question, since the moment bounds $C_k\le\dim\Fix_G(k)\le\dim\Fix_{S_N}(k)$ say nothing about the growth rate, and a non-amenable witness inside $G$ would have to be global, for instance a non-amenable group dual with all its non-trivial pieces on $\ge6$ points, such as $\widehat{\Z_3*\Z_3}$ acting through two $\Z_3$-blocks.

\subsection{Finite exotic quantum permutations}\label{sec:finite}

Let us try to exclude an exotic $G$ which is finite, $\dim C(G)<\infty$, as in Question~\ref{qu:coam}. Finiteness produces a Lagrange theorem, and lower bounds on the size of $G$ growing with the number of settled levels, but no contradiction. We write $D=\dim C(G)$, $m_k=\dim\Fix_G(k)$ and $s_k=\dim\Fix_{S_N}(k)$.

\begin{proposition}\label{prop:finstructure}
Let $G$ be a finite exotic quantum permutation group.
\begin{enumerate}
\item The characters of $C(G)$ are the evaluations $ev_\sigma$, $\sigma\in S_N$. So $C(G)=\C^{N!}\oplus A_{nc}$, where $A_{nc}$ is a sum of matrix blocks $M_{n_b}$ with $n_b\ge2$.
\item \emph{(Lagrange)} $D=q\cdot N!$, for an integer $q\ge2$, the index of $S_N$ in $G$.
\item \emph{(Cosets)} The algebra $C(G/S_N)$ of right $S_N$-invariant elements has dimension $q$, and $C(G/S_N)\simeq\bigoplus_\alpha f_\alpha V_\alpha$ as a $G$-module, where $f_\alpha=\dim V_\alpha^{S_N}$. So
\[
q=\sum_\alpha d_\alpha f_\alpha.
\]
\item The law of $\chi$ splits as $\mu_G=\frac1q\,\mu_{S_N}+\left(1-\frac1q\right)\nu$, where $\nu$ is the normalized law of $\chi$ on $A_{nc}$, a probability measure on $[0,N)$.
\end{enumerate}
\end{proposition}

\begin{proof}
(1) A character maps $u$ to a permutation matrix, so the classical version of $G$ is $S_N$, and the characters of a finite dimensional C$^*$-algebra are its $1$-dimensional blocks. (2) Dually, $\C[S_N]$ is a Hopf subalgebra of $C(G)^*$, and $N!\mid D$ by the Nichols--Zoeller theorem~\cite{nz}. (3) $C(G)$ is free over its quotient $C(S_N)$, which gives the dimension, and by Peter--Weyl the right $S_N$-invariant part of $C(G)=\bigoplus_\alpha V_\alpha\otimes\bar V_\alpha$ is $\bigoplus_\alpha V_\alpha\otimes\bar V_\alpha^{S_N}$. (4) The Haar state is the normalized regular trace, $h=\frac1D\sum_bn_b\mathrm{Tr}_b$, so the classical part has weight $N!/D=1/q$, and on it $\chi$ counts the fixed points. On a block of size $\ge2$, an eigenvector of $\chi$ at $N$ would satisfy $u_{ij}\xi=\delta_{ij}\xi$, and give a $1$-dimensional subrepresentation of an irreducible block.
\end{proof}

By (4), with $t=1/(q-1)$, the moments of $\nu$ are $\nu_k=C_k+t\,(C_k-s_k)$ for all $k$ at which $G$ agrees with $S_N^+$. For $t\to0$ these are the moments of the free Poisson law, but positivity fails quickly as $t$ grows.

\begin{theorem}[exact]\label{thm:finiteindex}
A finite exotic $G$ has index $q=\dim C(G)/N!\ge31$, and $q\ge32$ for $N\ge7$.
\end{theorem}

\begin{proof}
By Theorem~\ref{thm:levelseven}, $m_k=C_k$ for $k\le7$. Since $\nu$ lives on $[0,\infty)$, the Hankel matrix $H=(\nu_{i+j+1})_{0\le i,j\le3}$ is positive semidefinite. For $N\ge7$ we have $s_k=B_k$, the Bell numbers, for $k\le7$, and a direct computation gives $\det H=1-30t-14t^2+6t^3+t^4$. With $x=q-1=1/t$, this is $x^{-4}(x^4-30x^3-14x^2+6x+1)$, which is negative for the integers $1\le x\le30$. So $x\ge31$. For $N=6$ we have $s_7=B_7-1$, and the determinant is $1-29t-16t^2+5t^3+t^4$, which gives $x\ge30$.
\end{proof}

With the level-$8$ result of Corollary~\ref{cor:closeeight}, the Hankel matrix $(\nu_{i+j})_{0\le i,j\le4}$ gives $q\ge123$ for $8\le N\le37$. The thresholds for $t$ obtained from the levels $4,\dots,8$ are $1$, $0.414$, $0.116$, $0.0328$, $0.0082$, decreasing by a factor of about $4$ per level. Representation theory does better.

\begin{proposition}[computational, $8\le N\le37$]\label{prop:finfrob}
A finite exotic $G$ has index $q\ge1+d_4$, where $d_0=1$, $d_1=N-1$ and $d_{k+1}=(N-2)d_k-d_{k-1}$. For instance $q\ge1394$ at $N=8$.
\end{proposition}

\begin{proof}
By Corollary~\ref{cor:closeeight}, $G$ has the intertwiners of $S_N^+$ through level $8$. So $u^{\otimes k}$, $k\le4$, decomposes into irreducible representations $r_0,\dots,r_4$ of dimensions $d_0,\dots,d_4$, with the multiplicities $m_{kj}$ of $S_N^+$, and $m_{44}=1$. Counting $S_N$-fixed vectors gives $s_k=\sum_jm_{kj}f_j$. For $k\le3$ we have $s_k=C_k=m_{k0}$, so $f_1=f_2=f_3=0$, and for $k=4$ we get $15=14+f_4$, so $f_4=1$. Now Proposition~\ref{prop:finstructure}(3) gives $q\ge d_0f_0+d_4f_4$.
\end{proof}

The same argument shows that the more levels are settled, the larger a finite exotic $G$ must be. If $G$ agrees with $S_N^+$ through level $2m$, then $q\ge\sum_{j\le m}d_jf_j$, where the multiplicities $f_j$ are those of $S_N^+$, namely $1,0,0,0,1,1,6,20,86,393,\dots$ for $N$ large. So $q\ge d_m$ for $m\ge4$, and $d_m$ grows roughly like $(N-2)^m$.

All these are lower bounds. To exclude a finite exotic $G$, one needs an upper bound on $q$, or a contradiction between finiteness and the long initial agreement with $S_N^+$, and there are several possible sources for this: the coset algebra $C(G/S_N)$, a $q$-dimensional C$^*$-algebra with an ergodic action of $G$, whose ``double coset'' algebra has dimension $\sum_\alpha f_\alpha^2$; the fusion graph of $r_1$, a finite connected graph with Perron--Frobenius eigenvalue $N-1$ and integer Perron--Frobenius vector $(d_\alpha)$, starting with the chain $r_0-r_1-r_2-r_3$, and satisfying $\sum_\alpha d_\alpha^2=N!\sum_\alpha d_\alpha f_\alpha$; the fixed-point subfactor of the action of $G$ on $\C^N$, of index $N$ and finite depth, whose planar algebra lies between Temperley--Lieb and the spin model planar algebra of $S_N$; and integrality, the spectrum of $\chi$ consisting of totally real algebraic integers in $[0,N]$, with $\sum_km_kz^k$ a rational function with integer coefficients, by Proposition~\ref{prop:atom}. None of these gives an upper bound, as far as we can see: each of them is satisfied by a ``fake'' finite $G$ which looks like $S_N^+$ up to a high level, and then closes up. So the finite case does not seem to be easier than the maximality conjecture itself, or at least we do not know how to properly exploit the finiteness assumption.


%======================================================================
% Section: method
%======================================================================

\section{Partitions, levels and the method}\label{sec:method}

This section and the next two form the core of the paper. They attack the maximality conjecture head-on, through the Tannakian category of a hypothetical intermediate quantum group $S_N\subsetneq G\subsetneq S_N^+$. The idea is simple. The category of $G$ sits between the noncrossing partitions and all partitions, and it must be closed under the usual categorical operations. So we look for the first ``level'' $k$ at which it could contain a vector outside the noncrossing span, and we write down enough necessary conditions to show that no such vector exists.

At each level the conditions come in three layers: linear, quadratic and cubic. What makes the approach work in practice is that every layer can be reduced to finite, exact linear algebra, uniformly in $N$ wherever possible. Such computations cannot be checked by hand, so we put some care into making them \emph{certificates}: objects that can be verified independently, rather than numerical evidence.

This section sets up the framework and the tools. Section~\ref{sec:six} settles level $6$ for all $N\ge6$. Section~\ref{sec:higher} settles level $7$ for all $N\ge6$, and proves a structural theorem at level $8$ for $7\le N\le37$.

\subsection{Partition vectors and the two categories}

Let $N\ge1$, and let $e_1,\dots,e_N$ be the standard basis of $\C^N$. For a partition $\pi$ of $\{1,\dots,k\}$ we use the vector
\[
\xi_\pi=\sum_{i_1,\dots,i_k}\delta_\pi(i)\,e_{i_1}\otimes\dots\otimes e_{i_k}\ \in(\C^N)^{\otimes k},
\]
where $\delta_\pi(i)=1$ if the multi-index $i$ is constant on each block of $\pi$, and $\delta_\pi(i)=0$ otherwise. We write $\Pp(k)$ for the set of all partitions of $k$ points, and $\NC(k)\subset\Pp(k)$ for the noncrossing ones. Recall that
\[
\Fix_{S_N}(u^{\otimes k})=\spn\{\xi_\pi:\pi\in\Pp(k)\},\quad
\Fix_{S_N^+}(u^{\otimes k})=\spn\{\xi_\pi:\pi\in\NC(k)\}.
\]
More generally, partitions with $k$ upper and $l$ lower points act as linear maps $T_\pi:(\C^N)^{\otimes k}\to(\C^N)^{\otimes l}$, and composing or tensoring such maps corresponds to composing or juxtaposing partitions. Each closed loop created by a composition contributes a factor $N$, see~\cite{bsp}.

\begin{lemma}\label{lem:indep}
The vectors $\xi_\pi$, $\pi\in\Pp(k)$, are linearly independent if and only if $k\le N$. For $k=N+1$ there is exactly one linear relation, namely
\[
\sum_{\pi\in\Pp(k)}\mu(0_k,\pi)\,\xi_\pi=0,
\]
where $\mu$ is the M{\"o}bius function of the partition lattice, so that $\mu(0_k,\pi)=\prod_{B\in\pi}(-1)^{|B|-1}(|B|-1)!$.
\end{lemma}

\begin{proof}
By M{\"o}bius inversion, the vector $\sum_\pi\mu(0_k,\pi)\xi_\pi$ is the sum of all $e_{i_1}\otimes\dots\otimes e_{i_k}$ with pairwise distinct indices. The vectors $\delta_{\ker=\pi}$, sums over the multi-indices whose kernel is exactly $\pi$, are related to the $\xi_\pi$ by a unitriangular change of basis. They are nonzero, with disjoint supports, exactly when $\pi$ has at most $N$ blocks. For $k=N+1$, only the partition into singletons has more than $N$ blocks.
\end{proof}

The noncrossing vectors behave better: for $N\ge4$ the $\xi_\pi$, $\pi\in\NC(k)$, are linearly independent for every $k$. This is the nondegeneracy of the Temperley--Lieb Gram matrices at parameter $N\ge4$.

\subsection{Levels}

Let $S_N\subset G\subset S_N^+$, and write $\mathcal C(k)=\Fix_G(u^{\otimes k})$. Then
\[
\spn\{\xi_\pi:\pi\in\NC(k)\}\ \subset\ \mathcal C(k)\ \subset\ \spn\{\xi_\pi:\pi\in\Pp(k)\},
\]
and the family $\mathcal C=(\mathcal C(k))_k$, together with the corresponding spaces of intertwiners, is a tensor category containing all the noncrossing partitions. We abbreviate $\NC(k)$ for the span of the noncrossing partition vectors, when no confusion can arise.

\begin{definition}
The category $\mathcal C$ \emph{first differs at level $k$} if $\mathcal C(l)=\NC(l)$ for all $l<k$, and $\mathcal C(k)\ne\NC(k)$.
\end{definition}

If $G\ne S_N^+$ then $\mathcal C\ne\NC$, so $\mathcal C$ first differs at some level. The maximality conjecture is therefore equivalent to the following: for $S_N\subsetneq G$, the category of $G$ never first differs at any level. The work of this section and the next two is to exclude, one level at a time, that $\mathcal C$ first differs at level $k$.

When $k\le N$, Lemma~\ref{lem:indep} shows that every element of $\spn\{\xi_\pi:\pi\in\Pp(k)\}$ has a well-defined \emph{crossing part}, namely its coordinates on the crossing partitions. Since $\NC(k)\subset\mathcal C(k)$, the crossing part of an element of $\mathcal C(k)$ is itself an element of $\mathcal C(k)$. So $\mathcal C(k)\ne\NC(k)$ iff $\mathcal C(k)$ contains a nonzero purely crossing combination.

\begin{remark}\label{rem:known}
The known results at low levels are as follows.
\begin{itemize}
\item There is no easy exotic $G$ (Banica--Curran--Speicher~\cite{bcs}).
\item Easiness holds at level $2$ (Banica~\cite{bah}) and at level $3$ (McCarthy~\cite{mc1}).
\item The Haar integrals of degree $\le4$ agree with those of $S_N^+$ (McCarthy~\cite{mc1}).
\item $\mathcal C(k)=\NC(k)$ for $k\le5$ was announced by Freslon and Speicher. We reprove it in Proposition~\ref{prop:lowlevels}.
\item At order $6$, no exotic category contains a combination of at most $31$ crossing partitions (McCarthy~\cite{mc2}).
\end{itemize}
\end{remark}

\subsection{The operations}

We use four kinds of operations, all of them legitimate in any category containing $\NC$. For a vector at level $k$, the legs are indexed cyclically by $\Z_k$.

\begin{itemize}
\item \emph{Leg-sums:} $d_i\,\xi_\pi=N^{[\{i\}\in\pi]}\,\xi_{\pi\setminus i}$, which sums out the leg $i$. The factor $N$ appears exactly when $i$ is a singleton of $\pi$.
\item \emph{Merges:} $m_i\,\xi_\pi=\xi_{\pi/(i\sim i+1)}$, which identifies the legs $i$ and $i+1$.
\item \emph{Ribbon contractions:} given $\xi,\eta$ at level $k$, join an arc of $j$ consecutive legs of $\xi$ to an arc of $j$ consecutive legs of $\eta$ by $j$ parallel strings, in the planar way. Each of the two outermost strings may carry one output leg. A ribbon is described by its \emph{configuration} $(s,t,j,o_1,o_2)$: the strings join the legs $s,s+1,\dots,s+j-1$ of $\xi$ to the legs $t+j-1,\dots,t$ of $\eta$, and $o_1,o_2\in\{0,1\}$ say whether the first and the last strings carry an output leg. The output has level $2k-2j+o_1+o_2$, and every closed loop contributes a factor $N$.
\item \emph{Triple contractions:} the composite $x\mapsto\rho_2\big(\rho_1(x\otimes x)\otimes x\big)$ of a ribbon $\rho_1$ landing at level $k$ and a ribbon $\rho_2$ landing at level $<k$.
\end{itemize}

\begin{lemma}\label{lem:ops}
Let $\mathcal C$ be the category of $S_N\subset G\subset S_N^+$.
\begin{enumerate}
\item If $\xi\in\mathcal C(k)$, then $d_i\xi,m_i\xi\in\mathcal C(k-1)$.
\item If $\xi,\eta\in\mathcal C(k)$, then every ribbon contraction of $\xi\otimes\eta$ lies in $\mathcal C$.
\item If $x\in\mathcal C(k)$, then every triple contraction of $x$ lies in $\mathcal C$.
\end{enumerate}
\end{lemma}

\begin{proof}
Each operation is obtained from $\xi$, from $\xi\otimes\eta\in\mathcal C(2k)$, or from $\rho_1(x\otimes x)\otimes x\in\mathcal C(2k)$, by applying a map $T_\sigma$ with $\sigma$ noncrossing. These are a singleton, a merging block, or a family of nested caps, possibly with a three-leg block at each end of the nest. Since $\NC\subset\mathcal C$ and $\mathcal C$ is a tensor category, the result stays in $\mathcal C$. The factors $N$ are the usual loop factors of $T_\sigma$.
\end{proof}

We only ever use these operations to produce \emph{necessary} conditions. Whether they generate all the relevant operations is therefore irrelevant for the non-existence statements below. It is reassuring, however, that a natural class of further linear conditions turns out to be redundant (Remark~\ref{rem:partners}).

\subsection{The three-step method}

Fix $k$, and assume that $\mathcal C$ first differs at level $k$, with $k\le N$. Let $x\in\mathcal C(k)$ be purely crossing.

\begin{enumerate}
\item \emph{Linear step.} The leg-sums and merges send $x$ to $\mathcal C(k-1)=\NC(k-1)$, and they preserve $\NC$. So the crossing coordinates of $d_ix$ and $m_ix$ vanish. This is a linear system $M_k(N)x=0$, with
\[
M_k(N)=A_k+NB_k
\]
for integer matrices $A_k,B_k$, with $2k\,|\Pp(k-1)\setminus\NC(k-1)|$ rows and $|\Pp(k)\setminus\NC(k)|$ columns. We denote its kernel by $K_k(N)$, and we choose a basis $w_1,\dots,w_d$, writing $x=\sum_bc_bw_b$.
\item \emph{Quadratic step.} The ribbon contractions of $x\otimes x$ landing at level $<k$ lie in $\NC$, so their crossing coordinates vanish. These are quadratic forms in $c$, and they must all vanish at $c$.
\item \emph{Cubic step.} If the quadrics do not suffice, we add either \emph{closure conditions} or \emph{triple contractions}:
\begin{itemize}
\item closure: the ribbon contractions of $x\otimes x$ landing at level $k$ must lie in $\NC(k)\oplus(\mathcal C(k)/\NC(k))$;
\item triple contractions: $\rho_2(\rho_1(x\otimes x)\otimes x)$, landing at level $<k$, must have zero crossing part.
\end{itemize}
Both give cubic equations in $c$. Triple contractions have the advantage of requiring no knowledge of the dimension of $\mathcal C(k)/\NC(k)$.
\end{enumerate}

If the resulting system of polynomial equations has no nonzero solution, then $\mathcal C$ cannot first differ at level $k$.

\begin{table}[ht]
\centering
\begin{tabular}{@{}cccccc@{}}
\toprule
level $k$ & $|\Pp(k)|$ & $|\NC(k)|$ & crossing & $\dim K_k(N)$ & quadrics\\
\midrule
$4$ & $15$ & $14$ & $1$ & $1$ & --\\
$5$ & $52$ & $42$ & $10$ & $0$ & --\\
$6$ & $203$ & $132$ & $71$ & $5$ & $9$ of $15$\\
$7$ & $877$ & $429$ & $448$ & $14$ & $105$ of $105$\\
$8$ & $4140$ & $1430$ & $2710$ & $66$ & $1154$ of $2211$\\
\bottomrule
\end{tabular}
\caption{Sizes at each level. ``Quadrics'' is the dimension of the space of quadratic forms produced by the quadratic step, out of the dimension of all quadratic forms in $\dim K_k$ variables. The dimension of $K_k(N)$ is independent of $N$ in the range where it was computed.}
\label{tab:sizes}
\end{table}

Table~\ref{tab:sizes} shows how the problem grows. Each level brings its own phenomenon:
\begin{itemize}
\item at level $5$ the linear step alone suffices;
\item at level $6$ the quadrics leave $4$ lines, and a cubic step is needed;
\item at level $7$ the quadrics already exclude everything;
\item at level $8$ neither the quadrics nor the simplest cubic conditions suffice, but they force a strong symmetry.
\end{itemize}

\subsection{Levels up to five}

\begin{proposition}[exact, all $N\ge5$]\label{prop:lowlevels}
If $S_N\subsetneq G\subset S_N^+$ with $N\ge5$, then $\mathcal C(k)=\NC(k)$ for all $k\le5$.
\end{proposition}

\begin{proof}
We go level by level.
\begin{itemize}
\item \emph{Levels $k\le3$.} Here $\Pp(k)=\NC(k)$, so there is nothing to prove.
\item \emph{Level $4$.} The only crossing partition in $\Pp(4)$ is the basic crossing $\{13\}\{24\}$. If $\mathcal C(4)\ne\NC(4)$, its crossing part gives $\xi_{\{13\}\{24\}}\in\mathcal C(4)$. But $\NC$ and the basic crossing generate the category $\Pp$ of all partitions, which corresponds to $S_N$. So $G\subset S_N$, a contradiction.
\item \emph{Level $5$.} By the level-$4$ case, the linear step applies. The matrix $M_5(N)$ is square, of size $10\times10$: $10$ moves, times the single crossing partition of $\Pp(4)$, against the $10$ crossing partitions of $\Pp(5)$. Its determinant, computed exactly as a polynomial in $N$, is
\[
\det M_5(N)=-(N-4)(N^2-3N+1)^2,
\]
which does not vanish at any integer $N\ge5$. So $K_5(N)=0$. \qedhere
\end{itemize}
\end{proof}

\begin{remark}
Some comments on these factors.
\begin{itemize}
\item The roots of $N^2-3N+1$ are $4\cos^2(\pi/5)$ and $4\cos^2(2\pi/5)$. Similarly, the level-$6$ minor of Proposition~\ref{prop:lin} has the quadratic factor $N^2-4N+2$, with roots $4\cos^2(\pi/8)$ and $4\cos^2(3\pi/8)$. These are Beraha numbers, the values where Temperley--Lieb Gram determinants degenerate.
\item Since $\det M_5(N)$ comes from a square matrix, it is intrinsic up to sign. The level-$6$ quantity is a particular minor, and whether its quadratic factor is intrinsic has not been checked.
\end{itemize}
\end{remark}

\begin{remark}\label{rem:partners}
A natural source of further linear conditions is to contract $x\in\mathcal C(k)$ with noncrossing partitions $\eta$, landing at a level $<k$. We computed these at $N=11$ for $\eta$ of level $\le4$ (for $k=6,7$) and of level $\le3$ (for $k=8$). In every case they impose no condition beyond $M_k(N)x=0$. This supports the expectation that leg-sums and merges generate all the relevant linear operations.
\end{remark}

\subsection{How the computations are certified}\label{sec:certify}

The computations of Sections~\ref{sec:six} and~\ref{sec:higher} are far too large to be checked by hand. They were organized so that each claimed fact rests on one of the following elementary principles, rather than on floating-point evidence.

\begin{enumerate}
\item \emph{Polynomials in $N$ from finitely many values.} Every quantity we compute at level $k$ is a polynomial in $N$, once denominators are cleared, and we bound its degree \emph{a priori}. For instance:
\begin{itemize}
\item a coefficient of a ribbon contraction has degree $\le2\deg w+k$, because every closed loop of a ribbon contains a string;
\item a determinant has degree at most the sum of its row degrees.
\end{itemize}
A polynomial of degree $\le D$ is determined by $D+1$ values, which we compute exactly, and we always check a few further values.
\item \emph{Identities in $\Q[N]$.} When a vector depending on $N$ is claimed to lie in a kernel, the identity $M_k(N)w(N)=0$ is verified as an identity of polynomials, not at sample points.
\item \emph{Multimodular reconstruction.} Large determinants are computed modulo many primes, as polynomials over $\F_p$, and reconstructed by the Chinese remainder theorem. The number of primes is dictated by the Hadamard-type bound
\[
\|\det M\|_\infty\le\prod_i\Big(\sum_j\|M_{ij}\|_1\Big),
\]
where $\|\cdot\|_1$ is the sum of the absolute values of the coefficients. Spare primes are used as a consistency check.
\item \emph{Rational reconstruction.} At a fixed $N$, rational kernel bases are obtained modulo several primes and lifted by rational reconstruction. They are then \emph{verified} by an exact integer computation of $M_k(N)v$.
\item \emph{Ranks modulo $p$.} The rank of an integer matrix modulo a prime is a lower bound for its rank over $\Q$. So any ``full rank'' statement verified modulo $p$ holds in characteristic zero, provided the matrix modulo $p$ is the reduction of the true one; points 2 and 4 guarantee this. To use the rotation symmetry at level $8$ we work modulo $p\equiv1\bmod8$, so that $\F_p$ contains the $8$-th roots of unity.
\item \emph{Algebraic geometry.} Finiteness and degrees of solution sets come from Hilbert functions, via Gotzmann's persistence theorem and Macaulay matrices. Real roots of the resulting polynomials in $N$ are located by exact sign evaluation and Descartes' rule of signs.
\end{enumerate}

\begin{remark}
All code is in Python, using exact rational and modular arithmetic. Floating point is used only for exploration, never in a certificate. The computations ran on a two-core server and on two laptops, which handled the parallel parts: the Step~1 certificate at level $7$, and most values of $N$ at level $8$.
\end{remark}


%======================================================================
% Section: levelsix
%======================================================================

\section{Level six, settled for all \texorpdfstring{$N$}{N}}\label{sec:six}

Level $6$ is the first level not covered by the results recalled in Remark~\ref{rem:known} and Proposition~\ref{prop:lowlevels}. It is also the first level where the three-step method of Section~\ref{sec:method} runs in full. The linear step leaves $5$ dimensions, the quadrics leave $4$ lines, and a cubic step is needed to exclude them. The main result of this section is the following.

\begin{theorem}\label{thm:levelsix}
Let $N\ge6$, and let $S_N\subsetneq G\subset S_N^+$. Then
\[
\Fix(u^{\otimes6})=\spn\{\xi_\pi:\pi\in\NC(6)\}.
\]
Together with Proposition~\ref{prop:lowlevels}, $\dim\Fix(u^{\otimes k})$ is the Catalan number $C_k$ for every $k\le6$. So the first six moments of the main character of $G$ coincide with those of $S_N^+$.
\end{theorem}

The proof treats $N$ as a formal parameter throughout. Every step is reduced to the non-vanishing of an explicit polynomial in $N$, and the finitely many integer roots of these polynomials are all smaller than $6$. After the proof we describe the objects involved in more detail. At $N=6,7,8,11$ we give the exact solution set of the quadrics, and for $N$ formal we determine the parameters at which a candidate could survive. The answer to the latter is $N\in\{0,2,5\}$, and at $N=5$ the survivor is just the linear relation of Lemma~\ref{lem:indep}.

Throughout, $|\Pp(6)|=203$, $|\NC(6)|=132$, and the crossing part has dimension $71$.

\subsection{The linear step, uniformly in $N$}

The matrix $M_6(N)=A_6+NB_6$ has size $120\times71$: $12$ moves, times the $10$ crossing partitions of $\Pp(5)$.

\begin{proposition}[exact, all $N$]\label{prop:lin}
For every integer $N\ge3$ we have $\dim K_6(N)=5$. More precisely:
\begin{enumerate}
\item There are vectors $w_1(N),\dots,w_5(N)$ with entries in $\Z[N]$, of degree $\le4$, such that $M_6(N)w_b(N)=0$ identically in $N$. On the coordinates of the last five crossing partitions, in our ordering, they equal $N(N-2)$ times the standard basis.
\item A $66\times66$ minor of $M_6(N)$ equals
\[
\Delta_6(N)=-N^{29}(N-2)^3(N^2-4N+2).
\]
\end{enumerate}
Since $2\pm\sqrt2\notin\Z$, we have $\Delta_6(N)\ne0$ for $N\ge3$. So $\mathrm{rank}\,M_6(N)=66$, and the $w_b(N)$ form a basis of $K_6(N)$.
\end{proposition}

\begin{proof}
The proof is computational, as follows.
\begin{itemize}
\item The canonical kernel basis (reduced row echelon form) was computed exactly at $N=6,\dots,45$. Its entries were reconstructed as polynomials of degree $\le4$ divided by $N(N-2)$, by interpolation on $25$ values and verification on the other $15$.
\item The identity $M_6(N)w_b(N)=0$ was then checked in $\Q[N]$.
\item The entries of $M_6(N)$ have degree $\le1$, so $\Delta_6$ has degree $\le66$. It was interpolated from $67$ exact values and checked at two more. \qedhere
\end{itemize}
\end{proof}

\begin{remark}
Some comments on this basis.
\begin{itemize}
\item The five basis vectors involve $32$, $40$, $41$, $45$ and $50$ crossing partitions. This is consistent with McCarthy's bound $(31,6)$ stopping exactly there. Whether some element of $K_6(N)$ has support $\le31$ has not been checked.
\item The rotation of the six legs acts on $K_6(N)$ with eigenvalues $1,1,1,\omega,\bar\omega$, where $\omega=e^{2\pi i/3}$.
\end{itemize}
\end{remark}

\subsection{The quadrics, and finiteness}

Write $\xi=\sum_bc_bw_b(N)$ with $c\in\C^5$. A ribbon contraction of $\xi\otimes\xi$ landing at level $\le5$ must lie in $\NC$, so its crossing coordinates vanish. Such levels arise for $j=4,5,6$ strings, and they give $792$ quadratic forms in $c$, whose coefficients are polynomials in $N$.

\begin{lemma}\label{lem:deg}
The coefficients of these quadrics, and of the level-$6$ outputs used below, are polynomials in $N$ of degree $\le14$.
\end{lemma}

\begin{proof}
The entries of $w_b$ have degree $\le4$. Each closed loop of a ribbon contains a string, because the legs not on strings are output legs. So there are at most $6$ loops, and the degree is at most $4+4+6$.
\end{proof}

All these coefficients were computed exactly at $N=1,\dots,15$ and interpolated, then checked at $N=16,\dots,20$ and at $N=50$. The largest degree that actually occurs is $8$.

\begin{proposition}[exact, all $N$: finiteness]\label{prop:finite}
There are $9$ of the quadrics, $q_1,\dots,q_9$, with the following property. Restrict them to the hyperplane $c_3=0$, and form the degree-$3$ Macaulay matrix in the remaining variables $c_0,c_1,c_2,c_4$. It has $36$ rows ($q_a$ times a variable) and $20$ columns (the cubic monomials). One of its $20\times20$ minors equals
\[
D_2(N)=\lambda\,N^{48}(N-1)(N-2)^{31}(N-3)^{9}(N-4)^{8}R_{41}(N),
\]
where $\lambda\in\Q^\times$ and $R_{41}\in\Q[N]$ has degree $41$. Moreover $R_{41}(6)\ne0$, and all the coefficients of $R_{41}(6+x)$ have the same sign. Hence $D_2(N)\ne0$ for every real $N\ge6$.
\end{proposition}

\begin{proof}
The degree of $D_2$ is at most the sum of the row degrees, namely $152$; its actual degree is $138$. It was interpolated from $153$ exact values and checked at two more. The sign statement is Descartes' rule of signs on $(6,\infty)$.
\end{proof}

\begin{corollary}\label{cor:line}
For $N\ge6$, the common zero set of $q_1,\dots,q_9$ in $\mathbb P^4$ is finite. Hence $W=\mathcal C(6)/\NC(6)$ has dimension at most $1$.
\end{corollary}

\begin{proof}
By Proposition~\ref{prop:finite}, the cubics $c_iq_a|_{c_3=0}$ span all cubics in $c_0,c_1,c_2,c_4$, so the zero set misses the hyperplane $c_3=0$. A projective variety of positive dimension meets every hyperplane, so the zero set is finite. The projectivization of $W$ is a linear space contained in it, hence a point at most.
\end{proof}

\subsection{Closure}

Let $W=\C\xi$ be a line, so that $\mathcal C(6)=\NC(6)\oplus\C\xi$. Every ribbon contraction of $\xi\otimes\xi$ landing at level $6$ must then lie in $\NC(6)\oplus\C\xi$. These are the types $(j,o)=(3,0)$ and $(4,2)$, where $o=o_1+o_2$, which makes $72$ contractions. Equivalently, the crossing part $o(c)$ of the output must be proportional to the crossing part $x(c)$ of $\xi$. So all the minors $o_P(c)x_Q(c)-o_Q(c)x_P(c)$ vanish, and these are cubic forms in $c$.

\begin{proposition}[exact, all $N$: closure]\label{prop:closure}
There are $35$ of these cubics, coming from the contractions of types $(3,0)$ and $(4,2)$ at a fixed position, whose coefficient matrix ($35\times35$, in the monomial basis of cubics) has determinant
\[
D_3(N)=\mu\,N^{127}(N-2)^{90}(N-3)^{2}(N-4)^{3}(N-5)\,R_{63}(N),
\]
where $\mu\in\Q^\times$ and $R_{63}\in\Q[N]$ has degree $63$. The polynomial $R_{63}$ has no integer root $N\ge6$:
\begin{itemize}
\item Descartes' rule on $R_{63}(16+x)$ shows that it has no root beyond $16$;
\item exact evaluation at $N=6,\dots,16$ gives the signs $+$, $-$, $-$, $+$, $+$, $+$, $+$, $+$, $+$, $+$, $+$;
\item Descartes' rule on $(6,\infty)$ allows at most $2$ roots.
\end{itemize}
So its real roots in $(6,\infty)$ lie in $(6,7)$ and in $(8,9)$, and $D_3(N)\ne0$ for every integer $N\ge6$.
\end{proposition}

\begin{proof}
The degree of $D_3$ is at most $329$, and its actual degree is $286$. It was interpolated from $330$ exact values and checked at two more.
\end{proof}

\begin{proof}[Proof of Theorem~\ref{thm:levelsix}]
By Proposition~\ref{prop:lowlevels}, $\mathcal C(l)=\NC(l)$ for $l\le5$. Assume $\mathcal C(6)\ne\NC(6)$, so that $\mathcal C$ first differs at level $6$.
\begin{itemize}
\item By Lemma~\ref{lem:ops} and Proposition~\ref{prop:lin}, the crossing part of $\mathcal C(6)$ lies in $K_6(N)$.
\item By Corollary~\ref{cor:line}, $W=\mathcal C(6)/\NC(6)$ is a line $\C\xi$, with $\xi=\sum c_bw_b(N)$ and $c\ne0$.
\item By Lemma~\ref{lem:ops}, the $35$ cubics of Proposition~\ref{prop:closure} vanish at $c$.
\item Since $D_3(N)\ne0$, these cubics span all cubic forms. So every cubic monomial vanishes at $c$, hence $c=0$, a contradiction. \qedhere
\end{itemize}
\end{proof}

\begin{remark}
Some comments on the roles of the two steps.
\begin{itemize}
\item Proposition~\ref{prop:closure} shows that the closure conditions alone, without the quadrics, admit no nonzero solution. The roles of the two steps are therefore clean: the quadrics give $\dim W\le1$, and closure rules out $\dim W=1$.
\item The two real roots of $R_{63}$ depend on the choice of minor, and we attach no meaning to them.
\end{itemize}
\end{remark}

\subsection{The exact solution set at $N=6,7,8,11$}

At individual values of $N$, the full set of $792$ quadrics can be analyzed completely, by Gr\"obner bases over $\Q$, the unknowns being the $5$ coordinates $c_0,\dots,c_4$ of a candidate in a basis of $K_6(N)$. This shows what the method is actually excluding, namely a finite set of complex candidates, all of which are then killed by closure, and it also shows that the quadrics alone are not enough, the closure step being needed in all the cases below.

\begin{proposition}[exact over $\Q$]\label{prop:quad}
For $N=6,7,8,11$, let $I\subset R=\Q[c_0,\dots,c_4]$ be the ideal generated by all the quadrics, in the basis of $K_6(N)$ used in the computation.
\begin{itemize}
\item The quadrics span a $9$-dimensional space.
\item The Hilbert function of $R/I$ is $H(2)=6$ and $H(d)=4$ for $3\le d\le6$.
\item The zero set of $I$ in $\mathbb P^4(\C)$ consists of exactly $4$ simple points, which form a single Galois orbit over $\Q$.
\item At $N=6$ all four points are real. At $N=7,8,11$, two are real and two are complex conjugate.
\item Every point satisfies $c_0=c_2$ and $c_1=c_4$, which reflects the dihedral symmetry.
\item At each of the $4$ points, every one of the $72$ level-$6$ contractions violates closure.
\end{itemize}
\end{proposition}

\begin{proof}
The proof is computational, as follows.
\begin{itemize}
\item \emph{Finiteness and degree.} Since $H(4)=H(5)=4=4^{\langle4\rangle}$, Gotzmann's persistence theorem gives $H(d)=4$ for all $d\ge4$, so the zero set is finite of degree $4$.
\item \emph{Distinct points.} Choose a linear form $\ell$ whose multiplication map $(R/I)_4\to(R/I)_5$ is invertible; this is an exact $4\times4$ determinant. The operators $T_i=\ell^{-1}c_i$ on $(R/I)_4$ commute, and their joint eigenvalues are the values of $c_i/\ell$ at the points. For a generic combination $T_g=\sum g_iT_i$, the characteristic polynomial $\chi_g\in\Q[t]$ is squarefree. This gives $4$ distinct simple points.
\item \emph{Galois orbit.} $\chi_g$ is irreducible, since its factorization pattern modulo some prime is $(4)$.
\item \emph{Closure.} Writing $T_i=p_i(T_g)$ gives a generic point $c=(p_i(t))$ over the field $F=\Q[t]/(\chi_g)$. The closure minors are elements of $F$, and they are nonzero, hence nonzero at all four embeddings $F\hookrightarrow\C$. \qedhere
\end{itemize}
\end{proof}

\begin{remark}
Some comments on these solution sets.
\begin{itemize}
\item At $N=6$, no factorization pattern $(1,3)$ appears modulo primes, while it does at $N=7,8,11$. This suggests a Galois group $D_4$ or smaller at $N=6$, against $S_4$ for the others.
\item Numerically, the relative closure violations at the $4$ points are at least $0.018$ at $N=6$, and at least $0.05$ at $N=7$. So none of the $72$ contractions comes close to being compatible.
\end{itemize}
\end{remark}

\subsection{The survival locus}

Everything at level $6$ is polynomial in $N$. So we may treat $N$ as a formal complex parameter, and ask for which $N\in\C$ some candidate survives all the level-$6$ constraints.

To answer, we work on the coordinate ring of the quadric scheme. For generic $N$ the quadrics cut out $4$ points, and on $(R/I)_4$ we have the commuting operators $T_i=\ell^{-1}c_i$. For a closure cubic $m$, the norm
\[
\mathrm{Nm}_\ell(m)(N)=\det m(T_0,\dots,T_4)=\prod_{\text{points }x}\frac{m(x)}{\ell(x)^3}
\]
vanishes iff $m$ vanishes at one of the points. It is a rational function of $N$.

\begin{proposition}[exact]\label{prop:locus}
Take two random integer combinations $m_1,m_2$ of the $134$ closure cubics coming from the contractions of types $(3,0)$ and $(4,2)$ at a fixed position, and two auxiliary forms $\ell,\ell'$.
\begin{itemize}
\item The four norms $\mathrm{Nm}_\ell(m_i)$, $\mathrm{Nm}_{\ell'}(m_i)$ are rational functions of type $(55,24)$.
\item The gcd of their numerators is
\[
N^{12}(N-2)^{12}(N-5)^{2}.
\]
\item The denominators of $\mathrm{Nm}_\ell$ and $\mathrm{Nm}_{\ell'}$ have no common root.
\end{itemize}
Hence, outside the finitely many $N$ where the Hilbert function of the quadric ideal jumps, a level-$6$ candidate can survive closure only if $N\in\{0,2,5\}$.
\end{proposition}

\begin{proof}
Each norm was computed exactly at $N=6,\dots,199$ (for $\ell$) or $N=6,\dots,95$ (for $\ell'$). It was reconstructed by Pad{\'e} approximation from $80$ values, and checked on all the others. The roots of the denominators are the values of $N$ at which $\ell$ vanishes at one of the points. At those values, the norms with respect to the other form apply.
\end{proof}

\begin{proposition}[the parameter $N=5$]\label{prop:nfive}
At $N=5$ the Hilbert function of the quadric ideal stabilizes at $6$, and the quadric scheme is a single rational point of multiplicity $6$. This point is the crossing part of the relation of Lemma~\ref{lem:indep},
\[
\sum_{\pi\in\Pp(6)}\mu(0_6,\pi)\,\xi_\pi=0\qquad(N=5),
\]
with coordinates $(6,6,6,-1,6)$ in the basis used.
\end{proposition}

\begin{remark}
Some comments on the survival locus.
\begin{itemize}
\item The only parameter $N\ge3$ at which a level-$6$ candidate survives is therefore $N=5$, and the survivor is the linear dependence among the partition vectors. As an actual vector it is noncrossing, so it lies in every category, trivially.
\item In this sense the level-$6$ constraints are \emph{sharp}: away from degenerate parameters, they exclude everything except genuine linear relations.
\item No Jones-type values $4\cos^2(\pi/n)$ appear in the survival locus, only in the linear step.
\item A complete version of Proposition~\ref{prop:locus} would use all $72$ closure contractions, and would determine the finite set where the quadric scheme jumps.
\end{itemize}
\end{remark}

\subsection{Triple contractions at level six}

The closure step uses the fact that $W$ is a line. The triple contractions of Section~\ref{sec:method} avoid this assumption, and level $6$ is a good place to test their strength.

\begin{proposition}[computational, $N=11$]\label{prop:triplesix}
At level $6$ and $N=11$, the quadrics times linear forms span $31$ of the $35$ cubic monomials in $c$. Adding the cubics of the triple contractions with $\rho_1$ of types $(j,o)=(3,0)$ and $(4,2)$ at position $0$, and $\rho_2$ ranging over the ribbons at position $0$ landing at level $\le5$, gives all $35$.
\end{proposition}

So at level $6$, quadrics and triple contractions together exclude every candidate, without passing through Corollary~\ref{cor:line} and closure. This is the version of the cubic step that we use at level $8$, where nothing like Corollary~\ref{cor:line} is available.


%======================================================================
% Section: higher
%======================================================================

\section{Levels seven, eight and nine}\label{sec:higher}

In this section we push the method of Section~\ref{sec:method} further. Level $7$ is settled completely, for all $N\ge6$, the quadrics alone excluding everything, although the candidate space is larger than at level $6$ ($14$ dimensions against $5$). Level $8$ is the first level where the method meets real resistance: the quadrics and the simplest cubic conditions no longer suffice, but they force a strong symmetry, the new vectors being invariant under the half-turn of the eight legs, for $7\le N\le37$. Using the rotation symmetry systematically, we then close level $8$ computationally, $N$ by $N$, for $8\le N\le37$, and level $9$ numerically, by positivity, at eight values of $N$.

\begin{theorem}\label{thm:main}
Let $S_N\subsetneq G\subset S_N^+$.
\begin{enumerate}
\item If $N\ge6$, then $\dim\Fix(u^{\otimes k})=C_k$ for all $k\le7$.
\item If $7\le N\le37$ and $G$ first differs from $S_N^+$ at level $8$, then $\Fix(u^{\otimes8})/\NC(8)$ consists of half-turn-invariant vectors.
\item \emph{(computational)} If $8\le N\le37$, then $\dim\Fix(u^{\otimes k})=C_k$ for all $k\le8$.
\item \emph{(numerical)} If $N\in\{9,\dots,13,17,19,23\}$, then $\dim\Fix(u^{\otimes k})=C_k$ for all $k\le9$.
\end{enumerate}
\end{theorem}

Part 1 is Theorem~\ref{thm:levelseven} below, part 2 is Theorem~\ref{thm:halfturn}, part 3 is Corollary~\ref{cor:closeeight}, and part 4 is Theorem~\ref{thm:nine}.

\subsection{Level seven}\label{sec:levelseven}

Here $|\Pp(7)|=877$ and $|\NC(7)|=429$, so the crossing part has dimension $448$. By Theorem~\ref{thm:levelsix} and Proposition~\ref{prop:lowlevels}, for $S_N\subsetneq G$ and $N\ge6$ we already know $\mathcal C(l)=\NC(l)$ for $l\le6$.

\begin{proposition}[exact, $N\ge7$: linear step]\label{prop:linseven}
Let $M_7(N)=A_7+NB_7$, of size $994\times448$, and $K_7(N)=\ker M_7(N)$.
\begin{enumerate}
\item There are $14$ vectors $w_1(N),\dots,w_{14}(N)$, with entries in $\Q[N]$ of degree $\le9$, such that $M_7(N)w_b(N)=0$ identically in $N$. On $14$ fixed coordinates they equal $\delta(N)$ times the standard basis, where
\[
\delta(N)=\tfrac12(N-1)(N-5)(2N^4-17N^3+50N^2-57N+16),
\]
whose only integer roots are $1$ and $5$.
\item A $434\times434$ minor of $M_7(N)$ equals
\[
\Delta_7(N)=\lambda\,N^{47}(N-1)(N-2)(N-3)(N-4)\,R_{21}(N),
\]
where $R_{21}$ has degree $21$, $R_{21}(7)\ne0$, and $R_{21}(7+x)$ has coefficients of a single sign.
\end{enumerate}
Hence $\dim K_7(N)=14$ for all $N\ge7$, with basis $w_1(N),\dots,w_{14}(N)$.
\end{proposition}

\begin{proof}
The proof is computational, as follows.
\begin{itemize}
\item The kernel was computed exactly, by sparse elimination, at $N=7,\dots,76$. The common denominator was found by Pad{\'e} reconstruction, and the numerators by interpolation, with checks on held-out values.
\item The identity $M_7(N)w_b(N)=0$ was verified in $\Q[N]$.
\item $\Delta_7$ has degree $\le434$; its actual degree is $72$. It was computed modulo $37$ primes of $30$ bits, on two laptops, as polynomials over $\F_p$. It was then reconstructed by Chinese remaindering: the Hadamard bound $2^{824}$ requires $28$ primes, and the other $9$ confirm. \qedhere
\end{itemize}
\end{proof}

\begin{proposition}[exact, $N\ge7$: quadrics]\label{prop:quadseven}
Write $\xi=\sum_bc_bw_b(N)$ with $c\in\C^{14}$. There are $105$ ribbon contractions landing at level $\le6$, from $35$ configurations, whose $105\times105$ coefficient matrix, in the monomial basis of quadratic forms, has determinant
\[
D_7(N)=\mu\,N^{54}(N-1)^{203}(N-2)^{25}(N-3)^{4}(N-4)^{30}(N-5)^{263}(N-6)^{14}\,R_{1058}(N).
\]
Here $R_{1058}$ has degree $1058$ and exactly one real root in $(7,\infty)$, which lies in $(8,9)$. So $D_7(N)\ne0$ for every integer $N\ge7$, and for these $N$ the quadrics span all $105$ quadratic forms on $\C^{14}$.
\end{proposition}

\begin{proof}
The proof is computational, as follows.
\begin{itemize}
\item Each coefficient is a polynomial in $N$ of degree $\le9+9+7=25$, by the argument of Lemma~\ref{lem:deg}. The coefficients were computed exactly at $N=1,\dots,32$ and interpolated, with $6$ check values. The actual degrees are $17$ to $19$.
\item $D_7$ has degree at most $1890$, the sum of the row degrees; its actual degree is $1651$. Its coefficients are bounded by $2^{3762}$. It was computed modulo $140$ primes and reconstructed from $126$ of them; the other $14$ confirm.
\item Descartes' rule gives exactly one sign change for $R_{1058}(7+x)$, hence exactly one real root beyond $7$. Exact evaluation gives the signs $-,-,+$ at $N=7,8,9$. \qedhere
\end{itemize}
\end{proof}

\begin{proposition}[exact: the case $N=6$]\label{prop:sixseven}
At $N=6$ the vectors $\xi_\pi$, $\pi\in\Pp(7)$, satisfy exactly one linear relation (Lemma~\ref{lem:indep}). Let $r$ be its crossing part.
\begin{itemize}
\item $r$ lies in $\ker M_7(6)$, which has dimension $14$.
\item Complete $r$ to a basis $r,e_1,\dots,e_{13}$ of $\ker M_7(6)$, and write $x=c_0r+\sum_{i\ge1}c_ie_i$. Then no quadric involves $c_0$; this was checked in exact arithmetic for all $1147$ rows used.
\item The quadrics span all $91$ quadratic forms in $c_1,\dots,c_{13}$.
\end{itemize}
\end{proposition}

\begin{theorem}\label{thm:levelseven}
Let $N\ge6$ and $S_N\subsetneq G\subset S_N^+$. Then $\Fix(u^{\otimes7})=\spn\{\xi_\pi:\pi\in\NC(7)\}$, and hence $\dim\Fix(u^{\otimes k})=C_k$ for all $k\le7$.
\end{theorem}

\begin{proof}
Let $\xi\in\mathcal C(7)$ be purely crossing.
\begin{itemize}
\item \emph{Case $N\ge7$.} By Proposition~\ref{prop:linseven}, $\xi=\sum c_bw_b(N)$. All the quadrics of Proposition~\ref{prop:quadseven} vanish at $c$. Since they span all quadratic forms, $c_b^2=0$ for all $b$, so $c=0$.
\item \emph{Case $N=6$.} The same argument, with Proposition~\ref{prop:sixseven}, gives $c_1=\dots=c_{13}=0$. So $\xi=c_0r$, which as a vector equals $-c_0$ times the noncrossing part of the relation, and lies in $\NC(7)$. \qedhere
\end{itemize}
\end{proof}

\subsection{Level eight: what the first two steps give}

At level $8$ we have $|\Pp(8)|=4140$, $|\NC(8)|=1430$, and a crossing part of dimension $2710$. A uniform-in-$N$ treatment as at levels $6$ and $7$ is out of reach: a determinant of size $2211\times2211$ would have degree in $N$ in the tens of thousands. So we work at individual values of $N$. The exploratory computations below were done at $N=11$.

\begin{proposition}[computational, $N=11$]\label{prop:eight}
\begin{enumerate}
\item $\dim K_8=66$. This was computed modulo a prime, and confirmed in floating point by a singular value gap of $13$ orders of magnitude.
\item Over all $960$ ribbon configurations, the quadrics span a space of dimension $1154$, out of $2211$.
\item Let $I$ be the ideal of the quadrics. Then $\dim(R/I)_3=2350$. Splitting by the eigenvalues $\omega^j$ of the rotation of the eight legs, where $\omega=e^{2\pi i/8}$, the contributions are
\[
(602,\ 2,\ 596,\ 2,\ 548,\ 2,\ 596,\ 2)\qquad(j=0,\dots,7).
\]
\end{enumerate}
\end{proposition}

\begin{remark}
Two comments on these numbers.
\begin{itemize}
\item Since $2211-1154=1057$ is not of the form $m(m+1)/2$, the common zero set $Z\subset\mathbb P^{65}$ of the quadrics is not a linear space.
\item A Newton search from $40$ random complex starting points never converges to a point of $Z$: the residual always stalls near $2\cdot10^{-3}$. The same holds inside the subspaces fixed by the dihedral group, by the rotations, by $R^2$, by $R^4$ and by the reflection. Numerically, $Z$ appears to be empty, but the degree-$3$ test of item 3 cannot decide this.
\end{itemize}
\end{remark}

\subsection{Triple contractions at level eight}

We now add the cubics of triple contractions (Section~\ref{sec:method}). Since the rotation of the legs is a symmetry of the whole construction, the space of constraints is rotation-invariant. Working modulo the prime $p_0=1048433\equiv1\bmod8$, we use coordinates $y$ on $K_8$ adapted to the rotation, with eigenvalues $\omega^j$. In these coordinates every constraint splits into its homogeneous components of each rotation character, and each component is itself a valid constraint.

\begin{proposition}[computational, $N=11$]\label{prop:tripleeight}
\begin{enumerate}
\item Consider the two triple contractions with
\[
\rho_1=(0,0,4,0,0),\qquad\rho_2\in\{(0,0,5,0,0),(0,0,5,0,1)\}.
\]
They give $519$ cubics. Together with the quadrics times linear forms, they span all cubic monomials of odd rotation character.
\item In the even sectors, which are the parts invariant under the half-turn $R^4$, the deficiencies decrease under further random triple contractions and then stall:
\begin{itemize}
\item $602\to298$ in sector $0$;
\item $596\to37$ in sectors $2$ and $6$;
\item sector $4$ was not run to its plateau ($548\to461$ after $8$ steps).
\end{itemize}
These values stayed constant over more than $25$ further steps.
\end{enumerate}
\end{proposition}

Item 1 has a direct consequence. If all cubic monomials of odd character lie in the ideal of constraints, then so does $y_a^3$ for every coordinate $y_a$ of odd character $j_a$, since $3j_a$ is then odd. So every candidate has vanishing odd coordinates, which means it is invariant under $R^4$.

\subsection{The half-turn theorem}

\begin{theorem}[exact, $7\le N\le37$]\label{thm:halfturn}
Let $7\le N\le37$, and let $S_N\subsetneq G\subset S_N^+$ first differ from $S_N^+$ at level $8$. Then every element of $\Fix(u^{\otimes8})$ is, modulo $\NC(8)$, invariant under the rotation of the eight legs by four positions.
\end{theorem}

\begin{proof}
Fix $N$. The following certificate was computed for each $N$ in the range.
\begin{itemize}
\item \emph{Exact kernel.} A rational basis of $K_8(N)$ is obtained modulo $5$ to $10$ primes close to $2^{31}$ and lifted by rational reconstruction. It is certified by checking $M_8(N)v=0$ in exact integer arithmetic. Together with $\mathrm{rank}\,M_8(N)\ge2644$, which holds modulo a prime, this gives $\dim K_8(N)=66$. The basis reduces modulo $p_0$ to the modular kernel used below, which makes the modular computations reductions of the true ones.
\item \emph{Constraints.} We use the quadrics of the ribbon configurations anchored at position $0$, with all their rotation-character components, together with the $519$ cubics of item 1 of Proposition~\ref{prop:tripleeight}. In every case the quadrics span $1154$ dimensions, with character decomposition $(134,152,135,152,142,152,135,152)$.
\item \emph{Conclusion.} The Macaulay matrix in degree $3$ has full rank, modulo $p_0$, in each of the four odd sectors. Ranks modulo $p_0$ are lower bounds for ranks over $\Q(\omega)$. So, as explained above, all odd coordinates vanish on every candidate.
\item \emph{The case $N=7$.} Here the vectors $\xi_\pi$, $\pi\in\Pp(8)$, satisfy one linear relation, whose crossing part $r$ lies in $K_8(7)$ and is rotation-invariant. Replacing $x$ by $x-\lambda r$ changes neither the class of $x$ modulo $\NC(8)$ nor its half-turn invariance, so we may impose that the $r$-coordinate vanishes. Without this, every odd sector keeps a deficiency of $78$, even after $48$ further triple contractions. With it, all odd sectors are full. \qedhere
\end{itemize}
\end{proof}

\begin{table}[ht]
\centering
\begin{tabular}{@{}lll@{}}
\toprule
$N$ & machine & time per $N$\\
\midrule
$7,\dots,11$ and $32,\dots,37$ & two-core server & $22$--$30$ min\\
$12,14,\dots,30$ & laptop A (two workers) & $12$--$14$ min\\
$13,15,\dots,31$ & laptop B (two workers) & $12$--$14$ min\\
\bottomrule
\end{tabular}
\caption{The computation of Theorem~\ref{thm:halfturn}. Each $N$ is an independent job, and all $31$ values gave identical intermediate data.}
\end{table}

\begin{remark}
Some comments on this theorem.
\begin{itemize}
\item Theorem~\ref{thm:halfturn} reduces the possible new level-$8$ vectors from the $66$-dimensional space $K_8$ to its $46$-dimensional half-turn-invariant part.
\item Closing sectors $2$ and $6$ as well gives invariance under the rotation by two positions, leaving $26$ dimensions. Random triple contractions stall there, at a deficiency of $37$ at $N=11$. This is overcome in Subsection~\ref{sec:closeeight}.
\item The intermediate data are the same for all $31$ values of $N$, apart from the linear relation at $N=7$. This strongly suggests that a uniform statement holds, for all $N\ge7$ or perhaps all $N\ge8$. Proving it would require a uniform-in-$N$ version of the odd-sector certificate.
\end{itemize}
\end{remark}

\subsection{Closing level eight}\label{sec:closeeight}

Let $R$ be the rotation of the $k$ legs. It maps $\mathcal C(k)$ to itself, since $u=\bar u$, and it permutes the crossing partitions, so it acts on the crossing coordinates. Over a field containing the $k$-th roots of unity, for instance $\F_p$ with $p\equiv1\bmod k$, we can split everything according to the eigenvalues $\omega^j$ of $R$.

\begin{proposition}[rotation trick]\label{prop:rotation}
The space $K_k$ of the linear step is invariant under $R$. On the $\omega^j$-eigenspace of $R$, the conditions coming from all leg-sums $d_i$ and all merges $m_i$ are equivalent to those coming from $d_0$ and $m_0$ alone.
\end{proposition}

\begin{proof}
Removing or merging the legs $i,i+1$ of $Rx$ amounts to removing or merging the legs $i+1,i+2$ of $x$, followed by a rotation at level $k-1$. Since rotations preserve the noncrossing span, the conditions of index $i+1$ on $x$ are those of index $i$ on $Rx$. This gives the invariance of $K_k$, and on an eigenvector $Rx=\omega^jx$ the conditions of index $i$ are those of index $0$ on $R^ix=\omega^{ij}x$.
\end{proof}

So $K_k$ comes directly with an eigenbasis, computed in blocks of size about $1/k$ of the crossing part. This reproduces $\dim K_k=5,14,66$ at $k=6,7,8$, and makes level $9$ accessible. It also explains why the random triple contractions of Proposition~\ref{prop:tripleeight} stall. After splitting by characters, a first ribbon at $(s,t)$ is equivalent to one at $(0,t-s)$, and when $x$ is invariant under $R^m$ only $t-s$ modulo $m$ matters. So the random choices covered few inequivalent contractions, and missed the second ribbons anchored at $s_2\ne0$ on $\rho_1(x\otimes x)$.

\begin{proposition}[computational, $8\le N\le37$]\label{prop:closeeight}
Let $8\le N\le37$, and let $y$ be coordinates adapted to the rotation on the $46$-dimensional half-turn-invariant part of $K_8$, given by Theorem~\ref{thm:halfturn}: there are $14$, $10$, $12$, $10$ of them in the sectors $j=0,2,4,6$. Modulo $p_0=1048433$:
\begin{enumerate}
\item The quadrics put $y_a^3$ in the ideal of constraints for the $20$ coordinates of characters $2$ and $6$. So every candidate is invariant under $R^2$, and $26$ coordinates remain.
\item The triple contractions put $y_a^3$ in the ideal for the $12$ coordinates of character $4$. So every candidate is invariant under $R$, and $14$ coordinates remain. In this ring, the complete family of triple contractions has $2\times64=128$ members.
\item In the $14$ rotation-invariant coordinates, the products of the quadrics with the quadratic monomials, together with the products of the quadrics and the triple contractions with the linear forms, span all $2380$ quartic monomials.
\end{enumerate}
Hence no candidate survives. The chain $46\to26\to14\to0$ is the same for all these $N$, and was confirmed modulo a second prime, $1048361$, for $8\le N\le15$.
\end{proposition}

The status of the three steps differs. Step $3$ is a full-rank statement, so it is certificate-type: it lifts to $\Q$ as soon as a nonzero minor is exhibited, as explained in Subsection~\ref{sec:certify}. Steps $1$ and $2$ are memberships modulo $p$, which do not lift automatically, and this is why the result is labeled computational.

\begin{corollary}[computational]\label{cor:closeeight}
Let $8\le N\le37$ and $S_N\subsetneq G\subset S_N^+$. Then $\Fix(u^{\otimes8})=\spn\{\xi_\pi:\pi\in\NC(8)\}$, and so $\dim\Fix(u^{\otimes k})=C_k$ for all $k\le8$.
\end{corollary}

\begin{proof}
By Theorem~\ref{thm:levelseven}, $G$ agrees with $S_N^+$ through level $7$. So either it agrees at level $8$ too, or it first differs there, and then Theorem~\ref{thm:halfturn} and Proposition~\ref{prop:closeeight} leave no candidate.
\end{proof}

\begin{table}[ht]
\centering
\begin{tabular}{@{}lll@{}}
\toprule
$N$ & machine & time per $N$\\
\midrule
$9,11,\dots,37$ & laptop A (two workers) & $17$ min\\
$8,10,\dots,36$ & laptop B (two workers) & $17$ min\\
$8,\dots,15$, second prime & laptops A and B & $17$ min\\
$11$, check & two-core server & --\\
\bottomrule
\end{tabular}
\caption{The computation of Proposition~\ref{prop:closeeight}, done with the self-contained script \texttt{tyson.py} (Python, with numpy and scipy), run overnight on the two laptops.}
\end{table}

\subsection{Positivity certificates}\label{sec:positivity}

The positivity method of Section~\ref{sec:pairings}, which settles the levels of $O_N^*\subset O_N^+$, transports to the present setting, and gives a much cheaper route to most of Theorem~\ref{thm:halfturn} and Proposition~\ref{prop:closeeight}. Let $S_N\subsetneq G\subset S_N^+$, with category $\mathcal C$ first differing from $\NC$ at level $k\le N$, and let $W=\mathcal C(k)\ominus\NC(k)\ne0$ be the orthogonal complement of the noncrossing span inside $\mathcal C(k)$. Besides the rotation $R$ of the legs, we use the antilinear involution $J=\mathrm{conj}\circ\mathrm{rev}$, conjugation of the coordinates composed with the reversal of the legs, so that $J\xi_\pi=\xi_{\mathrm{rev}\,\pi}$.

\begin{lemma}\label{lem:jreal}
$R$ and $J$ preserve $\mathcal C(k)$, $\NC(k)$ and $W$, and $JRJ=R^{-1}$. Hence each $W_\omega=\ker(R-\omega)\cap W$ is $J$-invariant, and every nonzero $W_\omega$ contains a nonzero $J$-real vector $x=Jx$.
\end{lemma}

\begin{proof}
The coordinates $u_{ij}$ are projections, so the adjoint of $u_{i_1j_1}\cdots u_{i_kj_k}$ is the same monomial read backwards: the entrywise conjugate of $u^{\otimes k}$ is $u^{\otimes k}$ with its legs reversed. Since conjugation maps the fixed vectors of a representation $v$ to those of $\bar v$, it maps $\mathcal C(k)$ onto $\mathrm{rev}(\mathcal C(k))$, i.e.\ $J\mathcal C(k)=\mathcal C(k)$. Rotation invariance is Frobenius reciprocity. Both operations permute the $\xi_\pi$ and preserve crossings, and are (anti)unitary, so they preserve $\NC(k)$ and $W$. The relation $JRJ=R^{-1}$ holds on legs. For $x\in W_\omega$, $RJx=JR^{-1}x=J(\bar\omega x)=\omega Jx$, so $J$ is an antilinear involution of $W_\omega$, with nonzero fixed vectors $x+Jx$ or $i(x-Jx)$.
\end{proof}

\begin{proposition}[exclusion by positivity]\label{prop:posS}
Let $U_\omega$ be the $\omega$-eigenspace of $R$ on $K_k(N)$, $U_\omega^J=\{x\in U_\omega:Jx=x\}$ its real form, and let $\mathrm{Re}\,q,\mathrm{Im}\,q$ be the real quadratic forms on $U_\omega^J$ obtained from the quadrics $q(x)=\mathrm{coeff}_\tau(\rho(x\otimes x))$ of the quadratic step. If for every $\omega$ some real combination of these forms is positive definite on $U_\omega^J$, then $\mathcal C$ cannot first differ from $\NC$ at level $k$.
\end{proposition}

\begin{proof}
Otherwise some $W_\omega$ contains a $J$-real $x\ne0$, at which all the quadrics, hence all the real forms, vanish, contradicting positivity. The crossing part of $x$ lies in $K_k(N)$, and is again $J$-real and an $R$-eigenvector, since $R$ and $J$ commute with the crossing part and with the orthogonal projection onto $\NC(k)$, so the argument runs in the crossing coordinates.
\end{proof}

On a $J$-real vector the sesquilinear forms $\mathrm{coeff}_\tau(\rho(x\otimes Jx))$ of Section~\ref{sec:pairings} coincide with the quadrics, so no new constraint is used. What changes is the question: whether the quadrics have a common \emph{real} zero on $U_\omega^J$, rather than a complex one on $U_\omega$, and for this a positive definite combination is a certificate, lifting to $\Q$ as in Section~\ref{sec:certify}, and allowing an $N$-dependent treatment as in Theorem~\ref{thm:ten}. Positivity is sufficient, not necessary. Only ribbons anchored at one position are needed, a rotation multiplying the restriction of a quadric to $U_\omega$ by a unimodular scalar.

\begin{proposition}[computational, $N=11$]\label{prop:poseight}
At level $8$, with the $1154$ quadrics of Proposition~\ref{prop:eight}, the real forms on $U_j^J$ admit a positive definite combination for every rotation character $j\ne0$: for $j$ odd ($\dim U_j^J=5$) they span all $15$ quadratic forms, for $j=2,6$ ($\dim10$) and $j=4$ ($\dim12$) the best combination has smallest eigenvalue $0.031$, resp.\ $0.019$, with $\|c\|=1$ and Frobenius-orthonormal forms. For $j=0$ ($\dim14$) only $18$ of the $105$ forms are seen, there is no positive definite combination, and a numerical search finds no real common zero either (\texttt{positivity.py}).
\end{proposition}

So the quadrics and positivity alone give the half-turn theorem and the steps $46\to26\to14$ of Proposition~\ref{prop:closeeight}, in seconds and by certificates, and the remaining block $j=0$ is the one closed by the quartic step $14\to0$, itself a full-rank statement. Together these make level $8$ certificate-type, as asked in Question~\ref{qu:eightexact} below. At level $9$ the same test applies to the full $307$-dimensional kernel, character by character, with no slices, the blocks having dimensions $33$ to $37$, and it closes the level, all nine blocks included, see Theorem~\ref{thm:nine} below.

\subsection{Level nine: first steps}

At level $9$ we have $|\Pp(9)|=21147$ and $|\NC(9)|=4862$, so the crossing part has dimension $16285$. The linear step, the quadrics and the slices below are computed modulo a prime, as at level $8$, while the positivity certificates of Theorem~\ref{thm:nine} at the end of the section are computed in floating point, from the same kernel recomputed over $\mathbb R$, character by character (\texttt{nine.py}).

\begin{proposition}[computational, $N=11,13$]\label{prop:nine}
Work modulo $p=1048573\equiv1\bmod9$, with the same script \texttt{tyson.py}.
\begin{enumerate}
\item $\dim K_9=307$, with contributions $(37,33,33,36,33,33,36,33,33)$ from the rotation characters $j=0,\dots,8$, and the kernel was checked.
\item The quadrics of all $90$ ribbon configurations landing at levels $\le8$ leave gaps of
\[
(4465,\ 4446,\ 4446,\ 4464,\ 4446,\ 4446,\ 4464,\ 4446,\ 4446)
\]
in the spaces of quadratic forms of each character, which have dimension about $5250$.
\end{enumerate}
The results are identical at $N=11$ and $N=13$.
\end{proposition}

So the sequence of $\dim K_k$ at levels $5,\dots,9$ is $0,5,14,66,307$, and at level $9$ the quadrics remove only about $15\%$ of the quadratic forms, much as at level $8$ and unlike level $7$. The level-$8$ experience suggests the next steps: restrict to a symmetric part, then cubic memberships, then a full-rank statement. Since $9$ is odd there is no half-turn, so the first step has to be different.

A natural replacement is to look at the supports of the candidates. Write $E_j\subset K_9$ for the $\omega^j$-eigenspace of the rotation $R$, over a field containing the $9$-th roots of unity. For a nonempty set $S\subset\Z_9$, the \emph{slice} $S$ consists of the candidates lying in $\bigoplus_{j\in S}E_j$. We say that the slice \emph{closes} if it contains no nonzero candidate, and then all the slices $T\subset S$ close too. There are $511$ slices, but fewer up to symmetry.

\begin{lemma}[Galois symmetry]\label{lem:galois}
Let $a\in\Z_9^*$. The slice $S$ closes if and only if the slice $aS$ closes.
\end{lemma}

\begin{proof}
The rotation $R$ and all the constraints (the linear step, the quadrics and the triple contractions) have rational coefficients. Let $\sigma$ be an automorphism of $\C$ with $\sigma(\omega)=\omega^a$, acting on vectors coordinatewise. If $x$ is a candidate then so is $\sigma(x)$, and $Rx=\omega^jx$ gives $R\sigma(x)=\omega^{aj}\sigma(x)$. So $\sigma$ maps the slice $S$ bijectively onto the slice $aS$.
\end{proof}

The orbits of $\Z_9^*=\{\pm1,\pm2,\pm4\}$ on the characters are $\{0\}$, $\{3,6\}$ and $\{1,2,4,5,7,8\}$, and the $511$ slices form $103$ orbits. Up to symmetry, there are $3$ slices made of one character ($j=0,1,3$) and $8$ made of two.

\begin{proposition}[computational, $N=11,13$]\label{prop:ninekitty}
Work modulo $p=1048573$, with the script \texttt{kitty.py}, starting from the data of Proposition~\ref{prop:nine}.
\begin{enumerate}
\item In none of the $511$ slices do the restricted quadrics span all the quadratic forms. So the quadrics alone close no slice, not even a single character.
\item For a random subspace $H\subset K_9$ of dimension $117$, spanned by $13$ random vectors in each $E_j$, the restricted quadrics span all the quadratic forms on $H$. So the common zeros of the quadrics form a cone of dimension at most $307-117=190$.
\item On the slices $j=1$ and $j=2$, which have $33$ coordinates each, all the cubic monomials have the single character $3j$, and there are $6545$ of them. The quadrics times the linear forms leave a deficiency of $2684$ in this sector. The first $40$ triple contractions, computed in $2.5$ hours, bring it down to $1749$, and no cube $y_a^3$ lies in the ideal yet.
\end{enumerate}
The results are identical at $N=11$ and $N=13$.
\end{proposition}


\begin{remark}
The computation was done with \texttt{kitty.py}, the successor of \texttt{tyson.py}, overnight on two laptops, item 3 taking $2.5$ hours per slice. At level $8$ the same script reproduces the chain $46\to26\to14\to0$ of Proposition~\ref{prop:closeeight}, the quadrics alone closing there every slice made of odd characters; at level $9$ they are much weaker. The slices $j=1,2,4$ are conjugate by Lemma~\ref{lem:galois}, and indeed gave the same deficiencies, step by step, as did the two values of $N$. Pushed further (\texttt{tiger.py}, $14$ hours and $300$ triple contractions per slice), the deficiencies stall, at $1344$ on the slice $j=1$ at $N=11$, and at $2884$ on the slice $j=3$ at $N=13$, flat over the last $100$ contractions. So the degree-$3$ route, which closed level $8$, does not close even a single-character slice at level $9$, and this is what led us to positivity.
\end{remark}

\subsection{Outlook}

Three things remain at levels $8$ and $9$. First, Proposition~\ref{prop:closeeight} should become exact.

\begin{question}\label{qu:eightexact}
Prove that if $S_N\subsetneq G\subset S_N^+$ first differs from $S_N^+$ at level $8$, then $\dim\Fix_G(u^{\otimes8})=C_8=1430$, by certificates which lift from $\F_p$ to $\Q$, at each $8\le N\le37$, and ideally for all $N\ge8$.
\end{question}

Step $3$ of Proposition~\ref{prop:closeeight} is a full-rank statement, which lifts to $\Q$ as soon as a nonzero minor is exhibited. Steps $1$ and $2$ are memberships modulo $p$, and there are two ways out: write the memberships with explicit rational coefficients, found by rational reconstruction and checked exactly, or replace steps $1$--$2$ by full-rank statements in each rotation sector, avoiding memberships altogether. The uniformity of the chain $46\to26\to14\to0$ in $N$ suggests that the coefficients are rational functions of $N$, as at levels $6$ and $7$, and this would also give the statement for all $N\ge8$. Second, the small values of $N$.

\begin{question}\label{qu:eightsmall}
Settle level $8$ at $N=6$ and $N=7$.
\end{question}

For $N\le7$ the partition vectors at level $8$ satisfy linear relations, so the coordinates are not the same. The half-turn theorem covers $N=7$, and the closure should follow by a direct adaptation of the computation, with the relations quotiented out first. Third, and most importantly, level $9$.

\begin{question}\label{qu:nine}
Show that $\Fix_G(u^{\otimes9})$ is spanned by the noncrossing partition vectors, for $S_N\subsetneq G\subset S_N^+$ first differing from $S_N^+$ at level $9$, for all $N\ge9$, and by exact certificates.
\end{question}

The sequence of $\dim K_k$ at levels $5,\dots,9$ is $0,5,14,66,307$, and at level $9$ the quadrics alone are much weaker than at level $8$, with $9$ odd there is no half-turn, and the cubic memberships stall (Proposition~\ref{prop:ninekitty}). What works is positivity, as in Subsection~\ref{sec:positivity}, on the $J$-real forms of the nine rotation blocks of $K_9$, which at level $8$ did in seconds what the cubics did in nights, and which at level $9$ answers the question at eight values of $N$.

\begin{theorem}[level $9$, numerical]\label{thm:nine}
Let $N\in\{9,\dots,13,17,19,23\}$. Then no $S_N\subsetneq G\subset S_N^+$ first differs from $S_N^+$ at level $9$. Together with Corollary~\ref{cor:closeeight}, the maximality conjecture holds through level $9$ at these $N$: any such $G$ has the fixed vectors of $S_N^+$ at all levels $k\le9$.
\end{theorem}

\begin{proof}
By Proposition~\ref{prop:posS} it suffices to find, for each rotation character $j\in\Z_9$, a positive definite combination of the real forms $\mathrm{Re}\,q,\mathrm{Im}\,q$ on $U_j^J$, where $\dim U_j^J=37,33,33,36,33,33,36,33,33$ for $j=0,\dots,8$, by Proposition~\ref{prop:nine}. We used the quadrics of the $90$ ribbon configurations of Proposition~\ref{prop:nine}, anchored at one position ($70$ of them at $N=12$, $67$ at $N=19$). On each $U_j^J$ their real forms span a space of dimension $117$ only, out of $\binom{d+1}{2}$ with $d=\dim U_j^J$, but this is enough: the best combination, with $\|c\|=1$ and Frobenius-orthonormal forms, has smallest eigenvalue at least $5.5\cdot10^{-4}$, $1.7\cdot10^{-3}$, $3.1\cdot10^{-3}$, $4.5\cdot10^{-3}$, $5.9\cdot10^{-3}$, $1.0\cdot10^{-2}$, $1.2\cdot10^{-2}$, $1.4\cdot10^{-2}$ at $N=9,\dots,13,17,19,23$, against a largest eigenvalue of about $0.45$, on every block (\texttt{nine.py}, about two hours per $N$ on a laptop, in floating point). The margins increase with $N$, and the smallest one is each time on the block $j=0$, the one which at level $8$ needed the quartic step.
\end{proof}


%======================================================================
% Section: orthogonal
%======================================================================

\section{Orthogonal seeds, duals and twists}\label{sec:orthproblem}

In this section we discuss the orthogonal analogue of the maximality conjecture: the chain
\[
O_N\subset O_N^*\subset O_N^+
\]
is conjectured to be complete. That is, $O_N^*$ should be the only quantum group strictly between $O_N$ and $O_N^+$. The
half of this statement concerning the quantum groups containing $O_N^*$, namely the maximality of $O_N^*\subset O_N^+$,
appears as a conjecture in Brannan, Collins and Vergnioux~\cite[Sec.~5]{bcv}, with a partial result there. We
set up the problem and reformulate it in terms of ``seeds''. Then we study the two most natural families of seeds, group
duals and twists. The main result is about the anticommutative twist $O_N^{-1}$: the quantum group $\langle O_N,O_N^{-1}\rangle$
agrees with $O_N^+$ up to level $\max(20,2N)$, and it equals $O_N^+$ for all $N$ provided that a single statement holds at
$N=3$, namely $\langle O_3,O_2^+\rangle=O_3^+$. The induction on $N$ behind this uses the generation result
$\langle O_N,O_{N-1}^+\rangle=O_N^+$ for $N\ge4$ of~\cite{bcv}, for which we include a short proof. The level-by-level
analysis with pairings is the subject of Section~\ref{sec:pairings}.

\subsection{Setting}

Let $N\ge 2$. For a pairing $\pi$ of $\{1,\dots,k\}$ let $\xi_\pi=\sum_{i:\ \ker i\ge\pi}e_{i_1}\otimes\dots\otimes e_{i_k}$.
For a closed subgroup $G\subset O_N^+$ write $\mathcal C_G(k)=\Fix_G(u^{\otimes k})$. By Frobenius reciprocity, the fixed
vectors determine all Hom spaces. We use:
\begin{itemize}[leftmargin=2em]
\item $\TL(k)=\spn\{\xi_\pi:\pi\in NC_2(k)\}=\mathcal C_{O_N^+}(k)$, of dimension $\mathcal C_{k/2}$ for $N\ge2$;
\item $\HL(k)=\spn\{\xi_\pi:\pi\in P_2^*(k)\}=\mathcal C_{O_N^*}(k)$, where $P_2^*$ are the pairings joining odd to even points. Here $O_N^*$ is the half-liberated orthogonal group~\cite{bsp}, defined by the relations $abc=cba$ between the coordinates; see Bichon and Dubois-Violette~\cite{bdv} for the half-commutative Hopf algebras in general.
For $N\ge m$ the $m!$ vectors $\xi_\pi$, $\pi\in P_2^*(2m)$, are linearly independent;
\item $\Br(k)=\spn\{\xi_\pi:\pi\in P_2(k)\}=\mathcal C_{O_N}(k)$.
\end{itemize}
When the dimension matters we write $\TL_N$, $\Br_N$. Intermediate quantum groups $O_N\subset G\subset O_N^+$ correspond to
categories $\TL\subset C\subset\Br$. Odd levels vanish for every $G\supset O_N$, and $\mathcal C_G$ is closed under tensor products,
rotation (since $u=\bar u$), contraction of adjacent points, and the antilinear involution $J\xi=\overline{\mathrm{rev}(\xi)}$.
As in Section~\ref{sec:method}, $G$ \emph{first differs at level $k$} from $O_N^+$ if $\mathcal C_G(l)=\TL(l)$ for $l<k$ and $\mathcal C_G(k)\ne\TL(k)$.

\begin{remark}[unitary picture]\label{rem:unitarypic}
By the $2\times2$ model of $O_N^*$ with $U_N$-valued entries and Schur--Weyl duality,
$\HL(2m)\simeq\mathrm{End}_{U_N}((\C^N)^{\otimes m})$, the image of $\C[S_m]$. Under this identification
$\TL(2m)$ becomes the span of the $C_m$ noncrossing permutations $NC(m)\subset S_m$. Since $-1\in O_N$,
the map $G\mapsto PG$ from intermediate $O_N\subset G\subset O_N^+$ to intermediate
$PO_N\subset H\subset PO_N^+=PU_N^+$ is injective, with $PO_N^*=PU_N$. It is not onto in general: its image consists of the
$H$ whose category is closed under rotation by one point. See Section~\ref{sec:related}.
\end{remark}

\begin{remark}[$N=2$]\label{rem:ntwo}
For $N=2$ we have $\dim\HL(2m)=\sum_{\lambda\vdash m,\ \ell(\lambda)\le2}f_\lambda^2=C_m=\dim\TL(2m)$, hence
$O_2^*=O_2^+$. Since $O_2\subset O_2^*$ is maximal~\cite{bbcc} (this case goes back to Podle\'s), every intermediate
$O_2\subsetneq G\subset O_2^+$ equals $O_2^+$. From now on the interesting range is $N\ge3$.
\end{remark}

\subsection{Seeds}

The maximality of $O_N\subset O_N^*$ is known~\cite{bbcc}. It reduces completeness to a statement about generation.

\begin{proposition}\label{prop:seeds}
The chain $O_N\subset O_N^*\subset O_N^+$ is complete if and only if
\[
\langle O_N,K\rangle=O_N^+\quad\text{for every closed subgroup }K\subset O_N^+\text{ with }K\not\subset O_N^*.
\]
\end{proposition}

\begin{proof}
Let $G$ be intermediate with $G\notin\{O_N,O_N^*,O_N^+\}$. If $G\subset O_N^*$, then $G=O_N$ by~\cite{bbcc}, a
contradiction. So $G\not\subset O_N^*$, and $K=G$ violates the condition. Conversely, if $K\not\subset O_N^*$ and
$\langle O_N,K\rangle\ne O_N^+$, then $\langle O_N,K\rangle$ is intermediate and contains $K$, so it is neither $O_N$ nor
$O_N^*$.
\end{proof}

We call such a $K$ a \emph{seed}. Note that $\langle O_N,K\rangle$ only depends on the $O_N$-conjugacy class of $K$.

\begin{remark}[seeds ruled out in advance]\label{rem:ruledout}
\leavevmode
\begin{itemize}
\item \emph{Quantum permutation groups.} If $K\subset S_N^+$, up to $O_N$-conjugation, then $\langle O_N,K\rangle\supset
\langle S_N,K\rangle$. So $K$ can only be an orthogonal counterexample if $\langle S_N,K\rangle\ne S_N^+$, i.e.\ if it is already
a counterexample to the maximality of $S_N\subset S_N^+$, provided that $\langle O_N,S_N^+\rangle=O_N^+$. We do not know a
proof of this last statement in general. Since $\langle O_N,S_N^+\rangle$ contains the corner of $\langle O_{N-1},S_{N-1}^+\rangle$,
Proposition~\ref{prop:chir} below reduces it by induction to the single case $N=4$, that is, to $\langle O_4,S_4^+\rangle=O_4^+$.
Alternatively, $S_N^+$ contains the dual of $\Z_2*\Z_2$ on $4$ points, which is conjugate under $O_N$ to $\mathrm{diag}(s,t,1,\dots,1)$,
so that $\langle O_N,S_N^+\rangle=O_N^+$ would also follow from Conjecture~\ref{conj:torus} below. Granting it, this excludes graph
quantum symmetries $G^+(X)$, Hadamard models, and so on.
\item \emph{Easy inputs.} If the hull of $K$ with $S_N$ or $H_N$ is easy, then so is $\langle O_N,K\rangle$, and the
classification of easy orthogonal quantum groups~\cite{bcs} leaves only $O_N,O_N^*,O_N^+$. In the unitary setting the same
argument shows that every easy $U_N\subset G\subset U_N^+$, in particular every example of Mang--Weber~\cite{mw}, has
$PG\in\{PU_N,PU_N^+\}$.
\item \emph{Local inputs.} Since $O_2^*=O_2^+$, an exotic $G$ contains no non-classical quantum subgroup acting on $2$
coordinates. This is the orthogonal analogue of the restriction proposition of Section~\ref{sec:exotic} (Proposition~\ref{prop:restriction}).
\end{itemize}
\end{remark}

\subsection{Group duals}\label{sec:gd}

Up to $O_N$-conjugation, a group dual $\widehat\Gamma\subset O_N^+$ is given by $a$ involutions $s_1,\dots,s_a$ and $b$
elements $g_1,\dots,g_b$ generating $\Gamma$, with $N=a+2b$. Each involution acts on one coordinate, and each pair
$g_j,g_j^{-1}$ acts on a $2\times2$ rotation block. In the complex basis $f_1,\dots,f_N$ where $u$ is diagonal, we get an
alphabet $A$ of $N$ letters with an involution $r\mapsto\bar r$: involution letters are self-inverse, and the two letters of a
block are inverse to each other. So $\Gamma$ is a quotient of $\Z_2^{*a}*F_b$. Unlike in Section~\ref{sec:duals}, generators of order
$>2$ are allowed.

\begin{proposition}\label{prop:gdbasic}
\leavevmode
\begin{enumerate}
\item $\mathcal C_{\widehat\Gamma}(k)=\spn\{f_w:w\in A^k,\ w_1\cdots w_k=1\text{ in }\Gamma\}$.
\item In the basis $(f_w)$, $\xi_\pi=\sum f_w$, summed over the words $w$ \emph{compatible} with $\pi$: for each pair
$(s,t)$ of $\pi$, $w_t=\overline{w_s}$. Hence the category of $\langle O_N,\widehat\Gamma\rangle$ is
\[
C(k)=\{x\in\Br(k):x\text{ is supported on trivial words}\},
\]
and it only depends on the laws of $\Gamma$ of length $k$.
\item $\widehat\Gamma\subset O_N$ iff $\Gamma$ is abelian. $\widehat\Gamma\subset O_N^*$ iff $\Gamma$ is
\emph{half-abelian}, meaning $xyz=zyx$ for all letters $x,y,z\in A$.
\end{enumerate}
\end{proposition}

\begin{proof}
(1) and (2) are direct; for example, $\sum_ie_i\otimes e_i=\sum_rf_r\otimes f_{\bar r}$ on each block. For (3), $O_N$ and
$O_N^*$ are generated over $\TL$ by the crossings $(13)(24)$ and $\kappa=(14)(25)(36)$, and their compatible words are
$xy\bar x\bar y$ and $xyz\bar x\bar y\bar z$.
\end{proof}

\begin{proposition}[reduction]\label{prop:jnha}
If group duals give a counterexample, then some \emph{just non-half-abelian} group gives one: a non-half-abelian $\Gamma$
all of whose proper quotients are half-abelian. Equivalently, the normal closure of the elements $xyz(zyx)^{-1}$ is the
unique minimal normal subgroup of $\Gamma$.
\end{proposition}

\begin{proof}
If $\Gamma\to\Gamma'$ is a quotient preserving the letters, then $\Gamma'$ has more laws, so
$\langle O_N,\widehat{\Gamma'}\rangle\subset\langle O_N,\widehat\Gamma\rangle$, and $\widehat{\Gamma'}$ is still a seed
when $\Gamma'$ is not half-abelian. Half-abelianity is a finite set of relations, so Zorn's lemma applies, as in the
just-non-abelian reduction of Section~\ref{sec:duals}.
\end{proof}

Let $D(\Gamma)=\{\pi\in P_2:\ \xi_\pi\in\mathcal C_{\widehat\Gamma}\}$ be the set of pairings whose vectors are fixed by $\widehat\Gamma$. By Proposition~\ref{prop:gdbasic}, the category of $\langle O_N,\widehat\Gamma\rangle$ is $\Br\cap\mathcal C_{\widehat\Gamma}$, and it contains $\spn\{\xi_\pi:\pi\in D(\Gamma)\}$.

\begin{definition}\label{def:easygd}
We say that $\widehat\Gamma\subset O_N^+$ is \emph{easy over $O_N$} if
\[
\Br\cap\mathcal C_{\widehat\Gamma}=\spn\{\xi_\pi:\pi\in D(\Gamma)\},
\]
that is, if the category of $\langle O_N,\widehat\Gamma\rangle$ is spanned by pairings.
\end{definition}

This is a strong condition: a priori, $\Br\cap\mathcal C_{\widehat\Gamma}$ could contain linear combinations of pairing vectors supported on trivial words, without any of the pairings in it being in $D(\Gamma)$.

\begin{theorem}[trichotomy]\label{thm:gdtri}
\begin{enumerate}
\item $D(\Gamma)$ is a category of pairings containing $NC_2$. So it is $NC_2$, $P_2^*$ or $P_2$, according to whether $\Gamma$ is not half-abelian, half-abelian but not abelian, or abelian.
\item If $\widehat\Gamma$ is easy over $O_N$, then $\langle O_N,\widehat\Gamma\rangle$ is $O_N^+$, $O_N^*$ or $O_N$, in the three cases above.
\end{enumerate}
\end{theorem}

\begin{proof}
(1) Composites of pairing maps are multiples $N^{\ell}\ne0$ of pairing maps, so $D(\Gamma)$ is closed under the category operations. The classification of categories of pairings~\cite{bcs,bsp} and Proposition~\ref{prop:gdbasic}(3) give the rest. (2) is then immediate from Definition~\ref{def:easygd}.
\end{proof}

\begin{proposition}[computational]\label{prop:gdcomp}
The group duals below are easy over $O_N$ at the levels shown. That is, for every generating tuple of permutations listed, taken up to simultaneous conjugation and grouped by their sets of trivial words, the joint dimensions are exactly those of $O_N$, $O_N^*$ or $O_N^+$, according to whether $\Gamma$ is abelian, half-abelian or neither:
\begin{center}
\begin{tabular}{@{}ccccc@{}}
\toprule
letters & $N$ & permutations in & distinct law sets & levels \\
\midrule
$1$ involution, $1$ rotation pair & $3$ & $S_2,\dots,S_6$ & $31$ (in $S_6$) & $4,6,8$ \\
$2$ rotation pairs & $4$ & $S_2,\dots,S_6$ & $305$ (in $S_6$) & $4,6,8$ \\
$2$ involutions, $1$ rotation pair & $4$ & $S_2,\dots,S_5$ & $595$ (in $S_5$) & $4,6,8$ \\
$1$ involution, $2$ rotation pairs & $5$ & $S_2,\dots,S_5$ & $2339$ (in $S_5$) & $4,6$ \\
$3$ involutions, $1$ rotation pair & $5$ & $S_2,\dots,S_4$ & $917$ (in $S_4$) & $4,6$ \\
\bottomrule
\end{tabular}
\end{center}
About $4500$ law sets in all, with no exception. For residually finite $\Gamma$, such as $\Z_3*\Z_3$, $\mathrm{SL}_2(\Z)$ or free groups, the laws of a given length are seen by finite quotients, so these searches cover such groups level by level.
\end{proposition}

\begin{conjecture}\label{conj:easiness}
Every group dual $\widehat\Gamma\subset O_N^+$ is easy over $O_N$. Hence, by Theorem~\ref{thm:gdtri}, group duals are never seeds of counterexamples.
\end{conjecture}

A natural approach to Conjecture~\ref{conj:easiness} would be a M\"obius-type argument, as in the classification of the orthogonal easy quantum groups by Raum and Weber~\cite{rw}. A special case is Conjecture~\ref{conj:torus} below, coming from $\Gamma=\Z_2*\Z_2$.

\begin{remark}[unitary group duals]
Seeds need not live in $O_N^+$: any $G\subset U_N^+$ gives $PG\subset PU_N^+=PO_N^+$, and one asks whether
$\langle PO_N,PG\rangle=PU_N^+$. For $G=U\widehat\Gamma U^*$ the data is a pair $(\Gamma,M)$ with $M=U^tU$ symmetric
unitary, a continuous family. We checked (computational) all $119$ law sets at $N=3$ with generators in $S_2,\dots,S_5$
and all $611$ law sets at $N=4$ with generators in $S_2,S_3,S_4$, at levels $4,6,8$ and several positions $M$. In all cases
$\langle PO_N,PG\rangle=PU_N^+$ when $PG$ is non-classical, and it is $PO_N$ or $PU_N=PO_N^*$ otherwise. These questions
are taken up in Section~\ref{sec:related}.
\end{remark}

\subsection{The anticommutative twist}\label{sec:twist}

Let $O_N^{-1}$ be the twist of $O_N$ in which $u_{ij}$ and $u_{kl}$ anticommute for $i\ne k$, $j\ne l$ and commute otherwise.
It is the quantum automorphism group of the $N$-hypercube~\cite{bbc}. It is not half-commutative for $N\ge3$, so it is a
seed. For a multi-index $i\in[N]^k$ let
\[
c(i)=(-1)^{\#\{s<t\,:\,i_s<i_t\}},\qquad \Theta e_i=c(i)e_i .
\]

\begin{lemma}[signs]\label{lem:signs}
Let $\pi\in P_2(k)$ and $\ker i\ge\pi$. Two blocks of $\pi$ with different labels contribute $-1$ to $c(i)$ if they cross,
and $+1$ otherwise, whatever the labels. Pairs of positions with equal labels contribute nothing. Consequently
$\Theta\xi_\pi$ is the twisted pairing vector of~\cite{bbc}, and
\[
\mathcal C_{O_N^{-1}}(k)=\Theta\,\Br_N(k),\qquad \mathcal C_{\langle O_N,O_N^{-1}\rangle}(k)=\Br_N(k)\cap\Theta\,\Br_N(k).
\]
\end{lemma}

\begin{proof}
For blocks $\{s_1<s_2\}$ labeled $a$ and $\{t_1<t_2\}$ labeled $b>a$, count the pairs ``$a$-position before
$b$-position''. The six relative orders $aabb,abba,bbaa,baab,abab,baba$ give $4,2,0,2,3,1$, which is odd exactly in the two
crossing cases.
\end{proof}

We write $G_N=\langle O_N,O_N^{-1}\rangle$. The most direct use of Lemma~\ref{lem:signs} gives the low levels.

\begin{proposition}[exact]\label{prop:twoN}
For every $k\le2N$ we have $\mathcal C_{G_N}(k)=\TL_N(k)$.
\end{proposition}

\begin{proof}
Let $x=\sum a_\pi\xi_\pi=\Theta\sum b_\pi\xi_\pi$. As $k/2\le N$, the coefficients are well defined. On $i$ with
$\ker i=\pi$ and distinct labels, Lemma~\ref{lem:signs} gives $a_\pi=(-1)^{cr(\pi)}b_\pi$. Now let $\pi$ be crossing, pick two
crossing blocks with union $Q$, and evaluate on $i$ whose kernel merges them, with distinct labels elsewhere. The pairings
below $\ker i$ are the three $\sigma$ that agree with $\pi$ outside $Q$, and $\pi$ is the only one with crossing chords
inside $Q$. A block outside $Q$ crosses the two chords of $\sigma$ in $Q$ with a total parity independent of $\sigma$, and the
two chords inside $Q$ carry equal labels. So $c(i)(-1)^{cr(\sigma)}=(-1)^{[\sigma=\pi]}$, and the equation at $i$ reads
$\sum_\sigma a_\sigma=\sum_\sigma(-1)^{[\sigma=\pi]}a_\sigma$, i.e.\ $a_\pi=0$.
\end{proof}

To go beyond $2N$ we induct on $N$, through corners. For a closed subgroup $H\subset O_{N-1}^+$ with fundamental
representation $v$, its \emph{corner} is the closed subgroup of $O_N^+$ with fundamental representation $\mathrm{diag}(v,1)$.
For $x\in(\C^N)^{\otimes k}$ and $S\subset\{1,\dots,k\}$, the \emph{slice} $x_S\in(\C^{N-1})^{\otimes(k-|S|)}$ collects the
coordinates of $x$ with label $N$ exactly on $S$. Since $\mathrm{diag}(v,1)^{\otimes k}=\bigoplus_Sv^{\otimes(k-|S|)}$,
\begin{equation}\label{eq:corner}
x\text{ is fixed by the corner of }H\iff x_S\in C_H(k-|S|)\text{ for all }S.
\end{equation}

\begin{lemma}[corners]\label{lem:corner}
For $N\ge3$, $G_N$ contains the corner of $G_{N-1}$. More precisely, if $x\in \mathcal C_{G_N}(k)$, then $x_S\in \mathcal C_{G_{N-1}}(k-|S|)$
for every $S$.
\end{lemma}

\begin{proof}
The slice of $\xi_\pi$ is $\xi_{\pi\setminus S}$ (in dimension $N-1$) if $S$ is a union of blocks of $\pi$, and $0$ otherwise.
So slices map $\Br_N$ to $\Br_{N-1}$. If the label $N$ sits exactly on $S$, the pairs of positions involving it contribute
$\#\{p\notin S,\ q\in S,\ p<q\}$ to the exponent of $c(i)$, a number depending only on $S$. So $(\Theta_Ny)_S=\pm\Theta_{N-1}y_S$.
Hence $x=\Theta_Ny$ with $x,y\in\Br_N$ gives $x_S\in\Br_{N-1}\cap\Theta_{N-1}\Br_{N-1}$. Conclude by~\eqref{eq:corner}.
\end{proof}

The second ingredient is a generation result of Brannan, Collins and Vergnioux~\cite[Thm.~4.1]{bcv}, the orthogonal
analogue of Chirvasitu's $U_N^+=\langle U_N,U_{N-1}^+\rangle$ for $N\ge3$~\cite{ch1,ch2}. Chirvasitu's argument uses
non-orthogonal decompositions, which are harmless for $U_N$, whose invariants are $GL_N$-invariants, but not for $O_N$.
We include a short proof in the language of slices used above; it moves a unit vector into a hyperplane, and this is
where the restriction $N\ge4$ enters.

\begin{proposition}[Brannan--Collins--Vergnioux]\label{prop:chir}
Let $N\ge4$. If $x\in\Br_N(k)$ is fixed by the corner of $O_{N-1}^+$, then $x\in\TL_N(k)$. Hence
$\langle O_N,O_{N-1}^+\rangle=O_N^+$.
\end{proposition}

\begin{proof}
Since $x$ is $O_N$-invariant, it is fixed by every conjugate of the corner, i.e.\ by $O^+(e^\perp)\oplus1$ for every unit
vector $e$. For $W=e^\perp$ and $S\subset\{1,\dots,k\}$, let $x^e_S$ be the slice of $x$ in a basis
containing $e$, viewed as a multilinear form on $W$. Here a vector $y$ is read as the multilinear form
$y(w_1,\dots,w_k)=\sum_iy_i\,(w_1)_{i_1}\cdots(w_k)_{i_k}$, with no complex conjugation; all the vectors $e,w_j$ below are real. By~\eqref{eq:corner}, $x^e_S\in\TL(W)$.
\begin{enumerate}
\item \emph{Normalize.} Fix $e$ and write $x^e_\emptyset=\sum_\pi\alpha_\pi\xi^W_\pi$, $\pi\in NC_2(k)$. Put
$y=\sum_\pi\alpha_\pi\xi_\pi\in\TL_N(k)$ and $z=x-y$. Then $z\in\Br_N(k)$, $z$ is fixed by the corners, and $z$ vanishes on
$W^k$. By $O_N$-invariance, $z$ vanishes on $H^k$ for every hyperplane $H$.
\item \emph{The coefficients do not depend on $e$.} Write $z^e_S=\sum_\sigma\beta^{e,S}_\sigma\xi^{e^\perp}_\sigma$, with
$\sigma\in NC_2$ of the remaining points. This expansion is unique, because $\dim e^\perp=N-1\ge2$ and the Temperley--Lieb
vectors are independent at loop parameter $\ge2$. If $g\in O_N$ maps $e$ to $e'$, then $z^{e'}_S\circ g=z^e_S$ and
$\xi^{e'^\perp}_\sigma\circ g=\xi^{e^\perp}_\sigma$, so $\beta^{e,S}=\beta^{e',S}=:\beta^S$.
\item \emph{Move $e'$ into $e^\perp$.} Take unit vectors $e\perp e'$ and evaluate $z^{e'}_S$ on $U=e^\perp\cap e'^\perp$. All
the arguments, $e'$ included, lie in the hyperplane $e^\perp$, so the value is $0$. Thus
$\sum_\sigma\beta^S_\sigma\xi^U_\sigma=0$, and since $\dim U=N-2\ge2$ we get $\beta^S=0$.
\item \emph{Conclude.} All slices $z^e_S$ vanish. Writing each argument as $w_j=\lambda_je+w_j'$ with $w_j'\perp e$ gives
$z(w)=\sum_S\lambda^Sz^e_S(w')=0$. So $x=y\in\TL_N(k)$.
\end{enumerate}
The argument works at each level separately: if $x\in\Br_N(k)$ has all its slices in $\TL_{N-1}$, then $x\in\TL_N(k)$.
\end{proof}

\begin{theorem}\label{thm:reduction}
If $\langle O_3,O_2^+\rangle=O_3^+$, with $O_2^+$ embedded as a corner, then $\langle O_N,O_N^{-1}\rangle=O_N^+$ for every
$N\ge2$. More precisely, if $\langle O_3,O_2^+\rangle$ agrees with $O_3^+$ up to level $k$, then so does $G_N$ for every $N\ge3$.
\end{theorem}

\begin{proof}
For $N=2$, $O_2^{-1}$ is not classical, so $G_2=O_2^+$ by Remark~\ref{rem:ntwo}. By Lemma~\ref{lem:corner}, $G_3$ contains
the corner of $G_2=O_2^+$, so $\mathcal C_{G_3}(l)\subset \mathcal C_{\langle O_3,O_2^+\rangle}(l)=\TL_3(l)$ for $l\le k$. For $N\ge4$, if
$\mathcal C_{G_{N-1}}(l)=\TL_{N-1}(l)$ for $l\le k$, then every $x\in \mathcal C_{G_N}(l)$, $l\le k$, has its slices in $\TL_{N-1}$ by
Lemma~\ref{lem:corner}, hence $x\in\TL_N(l)$ by Proposition~\ref{prop:chir}.
\end{proof}

So everything rests on $N=3$, where step~3 of Proposition~\ref{prop:chir} breaks down: $U$ is a line, and the
Temperley--Lieb vectors collapse. This is also the case that required extra care in the unitary setting~\cite{ch1,ch2},
and the case left open in~\cite{bcv}, where it is expected to hold, with numerical evidence, see the discussion
following~\cite[Thm.~4.1]{bcv}. We now reduce it further, to a question about words, using the torus of $O_2^+$.

Let $\Gamma=\Z_2^{*3}$ with generators $a_1,a_2,a_3$, and identify $i\in[3]^k$ with the word $a_{i_1}\cdots a_{i_k}$. For
$d=1,2,3$ let $\varphi_d:\Gamma\to\Z_2*\Z_2$ kill $a_d$. Call a word \emph{pairwise reducible} if $\varphi_d(w)=1$ for
$d=1,2,3$, and \emph{exotic} if it is pairwise reducible but $w\ne1$ in $\Gamma$. Let $T_N^+=\widehat{\Z_2^{*N}}$ be the free
torus of $O_N^+$, acting diagonally. Chirvasitu proved $O_N^+=\langle O_N,T_N^+\rangle$ for $N\ge2$~\cite[Theorem~3.3]{ch2}.
Level by level, this says that a vector of $\Br_N(k)$ supported on words trivial in $\Z_2^{*N}$ lies in $\TL_N(k)$.

\begin{proposition}\label{prop:torus}
\leavevmode
\begin{enumerate}
\item Every $z\in \mathcal C_{\langle O_3,O_2^+\rangle}(k)$ lies in $\Br_3(k)$ and is supported on pairwise reducible words.
\item If moreover $z(w)=0$ for every exotic word $w$ of length $k$, then $z\in\TL_3(k)$.
\item \emph{Rotation relations.} For $z\in\Br_3(k)$, a word $u$ and letters $x\ne y$,
\[
\sum_{p:\ u_p=y}z(u^{p\to x})=\sum_{p:\ u_p=x}z(u^{p\to y}),
\]
where $u^{p\to x}$ is $u$ with its $p$-th letter replaced by $x$.
\end{enumerate}
\end{proposition}

\begin{proof}
(1) $O_2^+$ contains $T_2^+$, so $\langle O_3,O_2^+\rangle$ contains $\mathrm{diag}(s,t,1)$, with $s,t$ free involutions.
Fixed vectors of this group dual are supported on words with $\varphi_3(w)=1$. Since $z$ is invariant under the coordinate
permutations in $O_3$, the same holds for $\varphi_1,\varphi_2$. (2) $z$ is then supported on trivial words, so it is fixed by
$T_3^+$, and we use~\cite[Theorem~3.3]{ch2}. (3) $z$ is fixed by the rotations in the plane $(e_x,e_y)$. Differentiating at
the identity gives $\sum_pA^{(p)}z=0$, with $A=E_{xy}-E_{yx}$ acting on the $p$-th factor. The $u$-coordinate of this
identity is the claimed relation.
\end{proof}

For a word $w=1$ in $\Gamma$ there is nothing to prove, and every pairwise reducible word of length $\le8$ is trivial. The
shortest exotic words have length $10$, for instance $ababcbabac=[(ab)^2,c]$. Here is how to kill them. Pick a position
$p$, let $x=w_p$, pick a letter $y\neq x$, and apply the rotation relation to $u=w^{p\to y}$. The term with $p$ is $z(w)$.
Every other term is $z(v)$ with $v$ obtained from $u$ by switching one more letter between $x$ and $y$. If none of these $v$ is
pairwise reducible, then $z(w)=0$ by Proposition~\ref{prop:torus}(1).

\begin{theorem}[exact]\label{thm:sixteen}
For every exotic word $w$ of length $k\le20$ there exist $p$ and $y$ such that every word $v$ above which is pairwise
reducible is exotic, and killed in the same way. For $k\le18$ there is no such $v$ at all. Consequently:
\begin{enumerate}
\item $\langle O_3,O_2^+\rangle$ agrees with $O_3^+$ up to level $20$;
\item $\langle O_N,O_N^{-1}\rangle$ agrees with $O_N^+$ up to level $\max(20,2N)$, for every $N\ge2$.
\end{enumerate}
\end{theorem}

\begin{proof}
The first statement is a finite check, done by exhaustive enumeration (\texttt{words.py}): there are $10$, $240$, $3654$, $45360$, $502482$ and $5186470$ exotic words
of lengths $10,12,\dots,20$ beginning with the letter $a$, and the others follow by permuting letters. At $k=20$, exactly $240$ of them need the second step, always with $p=1$ and $y=b$, through $1$ to $5$ exotic words $v$. Then (1) follows from
Proposition~\ref{prop:torus}, and (2) from Theorem~\ref{thm:reduction} and Proposition~\ref{prop:twoN}.
\end{proof}

\begin{remark}[what the check sees, and where it must stop]
\leavevmode
\begin{itemize}
\item Read $w$ as a path in the Cayley tree of $\Gamma$, from $1$ to the element $w$. A computation with the three
projections $\varphi_d$ shows that a word $v$ above is pairwise reducible only when the step $q$ ``looks like'' the step
$p$ in all three projections to $\Z_2*\Z_2$. For instance, this happens when $p$ and $q$ cross the same edge of the tree in
opposite directions. So $p$ must be a letter of the reduced form whose edge is crossed only once, and in all cases found
the working positions are of this type.
\item The single-step argument first fails at length $20$, for instance for $acccbaacbaaabcaabcbb$, whose reduced form $acbcbabcbc$ has length $10$. It cannot work for words such as $rr^{-1}r$ with $r=[(ab)^2,c]$, of length $30$, in
which every edge of the geodesic is crossed three times. A proof for all $k$ needs an induction. The natural candidate
uses the leftover pairwise reducible $v$, which replace an excursion $xCx$ by $yCy$, keeping the group element.
\item Independently, the conclusion of Theorem~\ref{thm:sixteen}(1) was confirmed by direct linear algebra at levels
$\le10$, without using the torus.
\end{itemize}
\end{remark}

\begin{conjecture}\label{conj:torus}
Every $z\in\Br_3$ supported on pairwise reducible words vanishes on exotic words. Hence
$\langle O_3,O_2^+\rangle=\langle O_3,\mathrm{diag}(s,t,1)\rangle=O_3^+$, and $\langle O_N,O_N^{-1}\rangle=O_N^+$ for all $N\ge2$.
\end{conjecture}

Note that $\langle O_3,\mathrm{diag}(s,t,1)\rangle$ is generated by $O_3$ and a group dual, with $\Gamma=\Z_2*\Z_2$ and
letters $s,t,1$. So Conjecture~\ref{conj:torus} is a case of Conjecture~\ref{conj:easiness}: it says that this group dual is easy over $O_3$.

As for a proof, the argument behind Theorem~\ref{thm:sixteen}, with one step up to length $18$ and two steps at length $20$, must fail in its single-step form for words like $rr^{-1}r$, with $r=[(ab)^2,c]$, which have length $30$. A natural approach is by induction on the backtracking of words in the Cayley tree of $\Z_2^{*3}$, using the leftover pairwise reducible words, which replace an excursion $xCx$ by $yCy$. A next step would be the exhaustive check at length $22$, about ten times the cost of length $20$, to see whether two steps still suffice.

\subsection{Other twists and seeds}\label{sec:othertwists}

Some further natural seeds, all computational, support the same picture.

\paragraph{Sign twists of $O_N$.} Twisting $O_N$ by its diagonal subgroup $\Z_2^N$, with signs given by a graph $\epsilon$ on $N$ vertices, makes rows $i,j$ anticommute iff $ij$ is an edge. Then $\mathcal C_{O_N^\epsilon}=\Theta_\epsilon\Br$ for a diagonal sign operator $\Theta_\epsilon$, and the complete graph gives $O_N^{-1}$. Here $O_N^\epsilon\subset O_N^*$ iff every triangle has an even number of edges, iff $\epsilon$ is complete bipartite, possibly empty. For all graphs at $N=3,4$ (levels $4,6,8$) and $N=5$ (all $34$ graphs, levels $4,6$), $\langle O_N,O_N^\epsilon\rangle$ is $O_N$ for the empty graph, $O_N^*$ for the complete bipartite ones, and $O_N^+$ otherwise.

\paragraph{$\theta$-twists.} For the Connes--Landi type deformations of $O_N$ by its maximal torus, at $N=4,6$ with $\theta\in\{\frac12,\frac13,\frac14,\frac15,\frac16,\frac25,\sqrt2-1\}$ and a random value, we always get $O_N^+$, except $\theta=\frac12$ at $N=4$, which is a complete bipartite sign twist, and gives $O_N^*$.

\paragraph{Projectively classical examples.} The $K$-half-liberated groups $U_{N,K}^*$ and $U_{N,\infty}$ of Banica and Bichon satisfy $PG=PU_N$, since in their model in $C(U_N^K)\rtimes\Z_K$ the products $u_{ij}u_{kl}^*$ become diagonal with commuting entries. So they lift to $O_N^*$, and together with the easy case (Remark~\ref{rem:ruledout}) this disposes of the whole Mang--Weber landscape~\cite{mw}.

So every natural non-easy construction we tried is an easy object seen through a change of basis, a twist, or a product construction, and obeys the rule ``not half-commutative $\Rightarrow O_N^+$''.



%======================================================================
% Section: pairings
%======================================================================

\section{Pairings and uncappable vectors}\label{sec:pairings}

We study the intermediate quantum groups $O_N\subset G\subset O_N^+$ level by level, with the notation of
Section~\ref{sec:orthproblem}, in the spirit of McCarthy~\cite{mc2} for $S_N\subset S_N^+$ and of Section~\ref{sec:method}. We first treat
$O_N^*\subset G\subset O_N^+$, i.e.\ categories $\TL\subset C\subset\HL$. The first level where $C$ can differ from $\TL$ is
$2m=6$, and there a difference forces $G=O_N^*$. If $G$ first differs from $O_N^+$ at level $2m\ge8$, the new vectors are
\emph{uncappable}, and gluing such a vector to itself produces quadratic constraints coming from the lower levels. A
positivity certificate for these constraints excludes the level. Levels $8$ and $10$ are excluded exactly for all
$N\ge3$, and levels $12$ and $14$ numerically. We then treat the full Brauer problem $\TL\subset C\subset\Br$, which gives
completeness of $O_N\subset O_N^*\subset O_N^+$ through $10$ points for several $N$. Finally we compare with the
permutation case of Section~\ref{sec:method}.

\subsection{Level 6}

\begin{proposition}\label{prop:6}
$P_2^*(6)=NC_2(6)\sqcup\{\kappa\}$ with $\kappa=(14)(25)(36)$, the rotated half-commutation crossing.
Hence if $\mathcal C_G(6)\neq\TL(6)$, then $\mathcal C_G(6)=\HL(6)\ni\xi_\kappa$, so $abc=cba$ holds in $G$ and $G=O_N^*$.
\end{proposition}
\begin{proof}
$|P_2^*(6)|=3!=6$ and $|NC_2(6)|=5$, and all of these vectors are independent for $N\ge3$. Since
$\HL(6)/\TL(6)$ is one-dimensional, any extra vector gives $\xi_\kappa$ modulo $\TL$.
\end{proof}
So if $G\ne O_N^*$, then $G$ can only first differ from $O_N^+$ at level $8$ or higher.

\subsection{The level where $G$ first differs}

From now on let $m\ge4$ be minimal with $C(2m)\ne\TL(2m)$, where $C=\mathcal C_G$, i.e.\ $G$ first differs from $O_N^+$ at level $2m$. Let
$W=C(2m)\ominus\TL(2m)$, the orthogonal complement for the standard scalar product.

\begin{definition}
For $1\le i\le 2m$, let $\mathrm{cap}_i:(\C^N)^{\otimes 2m}\to(\C^N)^{\otimes 2m-2}$ contract the points $i,i+1$,
cyclically, so that $\mathrm{cap}_{2m}$ contracts the last point with the first. A vector is \emph{uncappable} if all
$2m$ caps kill it. Let $U(2m)\subset\HL(2m)$ be the space of uncappable vectors.
\end{definition}

\begin{lemma}\label{lem:uncap}
$W\subset U(2m)$. Moreover $W$ is invariant under the one-click rotation $\rho$ and under $J$, and
$J\rho J=\rho^{-1}$. Each eigenspace $W_\omega=\ker(\rho-\omega)$ is $J$-invariant and therefore contains
a nonzero $J$-real vector $x=Jx$.
\end{lemma}
\begin{proof}
For $x\in W$, $\mathrm{cap}_i x\in C(2m-2)=\TL(2m-2)$ by minimality. For $y\in\TL(2m-2)$ we have
$\langle \mathrm{cap}_ix,y\rangle=\langle x,\mathrm{cup}_iy\rangle=0$, since $\mathrm{cup}_i y\in\TL(2m)$. So
$\mathrm{cap}_ix=0$. Rotation and $J$ preserve both $C$ and $\TL$, and they are (anti)unitary, so they preserve $W$. Finally
$\rho Jx=J\rho^{-1}x=J(\bar\omega x)=\omega Jx$ for $x\in W_\omega$, since $|\omega|=1$.
\end{proof}

\begin{definition}
For pairings $p,q$ of $2m$ points and $a\le 2m$, let $p\circ_a q$ be the planar gluing of the last $a$ points of $p$
to the first $a$ points of $q$ (rainbow), with a factor $N^{\#\text{loops}}$. Extend bilinearly.
\end{definition}

\begin{lemma}\label{lem:glue}
Let $x\in W$ be $J$-real and $3\le j<m$. Then $x\circ_{2m-j}x\in\TL(2j)$. When $N\ge j$, this means that the
coefficients of $x\circ_{2m-j}x$ on the crossing pairings of $P_2^*(2j)$ all vanish. In terms of sesquilinear forms:
for such a crossing pairing $\tau$, put $B_\tau(x,y)=\mathrm{coeff}_\tau(x\circ_{2m-j}Jy)$, which is linear in $x$ and antilinear
in $y$, and $q_\tau(x)=B_\tau(x,x)$. Then $\overline{q_\tau(x)}=q_{\mathrm{rev}\,\tau}(x)$, so the real parts $\mathrm{Re}\,q_\tau$ are real
quadratic forms on the real space of $J$-real vectors, and they all vanish at $x$.
\end{lemma}
\begin{proof}
$x\circ_{2m-j}x=x\circ_{2m-j}Jx$ is obtained from $x\otimes Jx\in C(4m)$ by noncrossing contractions, so it lies in
$C(2j)=\TL(2j)$. Reading $x$ as an operator $T_x$ from $2m-j$ to $j$ strands, $x\circ_{2m-j}Jx$ is the vector of $T_xT_x^*$.
Finally $J(a\circ_{2m-j}b)=Jb\circ_{2m-j}Ja$, so $z=x\circ_{2m-j}Jx$ satisfies $Jz=z$, and since $J\xi_\tau=\xi_{\mathrm{rev}\,\tau}$,
this gives $\overline{q_\tau(x)}=q_{\mathrm{rev}\,\tau}(x)$.
\end{proof}

\begin{theorem}[Exclusion criterion]\label{thm:crit}
Let $U_\omega$ be the eigenspaces of $\rho$ on $U(2m)$, with their real structures $U_\omega^J$, and let
$q_1,q_2,\dots$ be the real quadratic forms $\mathrm{Re}\,q_\tau$ of Lemma~\ref{lem:glue}, restricted to $U_\omega^J$. Suppose that for every
$\omega$ some real combination $\sum_kc_kq_k$ is positive definite on $U_\omega^J$. Then no $G$ with
$O_N^*\subset G\subset O_N^+$ first differs from $O_N^+$ at level $2m$.
\end{theorem}
\begin{proof}
Otherwise $W_\omega\neq0$ for some $\omega$. A $J$-real $0\ne x\in W_\omega$ satisfies $q_k(x)=0$ for all $k$,
hence $\sum c_kq_k(x)=0$, which contradicts positivity.
\end{proof}

For $N<m$ the same argument works modulo the Gram kernel. ``Lies in $\TL$'' then means that the projection on
$\HL\ominus\TL$, taken for the Gram scalar product, vanishes.

\subsection{Level 8, exact for all $N$}

$\dim\HL(8)=24$ and $\dim\TL(8)=14$. The $10$ crossing pairings in $P_2^*(8)$ form two rotation orbits: an orbit
$\mathcal O_8$ of size $8$, containing $(12)(36)(47)(58)$, and an orbit $\mathcal O_2=\{(14)(27)(36)(58),\ (16)(25)(38)(47)\}$.

\begin{proposition}[exact, all $N$]\label{prop:eightpair}
Let $\mathrm{Cap}(N)$ be the $48\times24$ matrix of the $8$ cyclic caps, from the coefficients on $P_2^*(8)$ to
those on $P_2^*(6)$.
\begin{enumerate}
\item A $22\times22$ minor of $\mathrm{Cap}(N)$ equals $-16N^2(N^2-1)(N^2+1)(N^2-2)$, so $\mathrm{rank}\,\mathrm{Cap}(N)\ge22$ for all integers $N\ge2$.
\item There are vectors $X_\pm(N)$ with entries in $\Z[N]$ of degree $\le3$ such that $\mathrm{Cap}(N)X_\pm(N)=0$ and
$\rho X_\pm=\pm X_\pm$ identically in $N$. Normalized by $x_+=X_+/(N-2)$ and $x_-=X_-/(N+2)$, they have coefficient $1$ on
$(12)(34)(56)(78)$. Their crossing parts are
\[
(N-2)\,x_+\big|_{\rm cross}=-(N^2+N-1)\Big(\sum_{\mathcal O_8}\xi_\pi-(N+2)\sum_{\mathcal O_2}\xi_\pi\Big),
\]
\[
(N+2)\,x_-\big|_{\rm cross}=(N^2-N-1)\Big(\sum_{\mathcal O_8}\epsilon(\pi)\xi_\pi-(N-2)\sum_{\mathcal O_2}\epsilon(\pi)\xi_\pi\Big),
\]
where $\epsilon=\pm1$ alternates along the rotation, with $\epsilon=1$ on $(12)(36)(47)(58)$ and on $(14)(27)(36)(58)$.
\item The coefficient $h_\pm(N)$ of $\xi_\kappa$ in $x_\pm\circ_5Jx_\pm\in\HL(6)$ is
\[
h_+(N)=-\frac{4N(N+3)(N^2+N-1)^2}{(N-2)^2},\qquad h_-(N)=-\frac{4N(N-3)(N^2-N-1)^2}{(N+2)^2}.
\]
\end{enumerate}
\end{proposition}

\begin{proof}
All entries involved are polynomials of known degree in $N$: degree $\le1$ for $\mathrm{Cap}$, $\le3$ for $X_\pm$,
and $\le2\cdot3+2$ for $(N\mp2)^2h_\pm$. Each identity was therefore verified by exact evaluation at more points than
its degree, and interpolated and rechecked on further values.
\end{proof}

\begin{theorem}[exact, level $8$, all $N\ge3$]\label{thm:eight}
Let $N\ge3$ and $O_N^*\subset G\subset O_N^+$. If $\mathcal C_G(6)=\TL(6)$, then $\mathcal C_G(8)=\TL(8)$.
\end{theorem}

\begin{proof}
Suppose $W=C(8)\ominus\TL(8)\neq0$.
\begin{itemize}
\item \emph{$N\ge4$.} The $24$ vectors are independent, and by Proposition~\ref{prop:eightpair}(1,2) we get $U(8)=\C x_+\oplus\C x_-$,
with distinct rotation eigenvalues. Since $W\subset U(8)$ is rotation-invariant (Lemma~\ref{lem:uncap}), it contains $x_+$
or $x_-$. Then $x_\pm\circ_5Jx_\pm\in C(6)=\TL(6)$, so $h_\pm(N)=0$. But $h_\pm(N)\ne0$ for every integer $N\ge4$.
\item \emph{$N=3$.} The Gram matrix of $P_2^*(8)$ has rank $23$. Exactly, $x_-$ has Gram norm $0$, i.e.\ it is the
linear relation, while $x_+$ has Gram norm $237600$. Since all vectors at level $6$ are independent at $N=3$, the
uncappable space is $\C x_+$, and $h_+(3)=-8712\neq0$. \qedhere
\end{itemize}
\end{proof}

\begin{remark}
Two comments on these formulas.
\begin{itemize}
\item The factor $N-3$ in $h_-$ is no accident: at $N=3$ the vector $x_-$ is itself the linear relation among the
partition vectors. This is the orthogonal analogue of the survival locus in Section~\ref{sec:method}, where the only
survivor at $N=5$ is the relation among the vectors of $\mathcal P(6)$.
\item The roots of $N^2\pm N-1$ are $\pm$ the golden ratio and its conjugate, and those of $N^2-2$ are $\pm\sqrt2$. These are
Beraha-type numbers, as in Section~\ref{sec:method}, and they never occur at integer $N$.
\end{itemize}
\end{remark}

\begin{corollary}
For $N\ge3$, every quantum group $O_N^*\subset G\subsetneq O_N^+$ whose fixed vectors differ from those of $O_N^+$
at some level $\le8$ is equal to $O_N^*$.
\end{corollary}

\subsection{Level 10, exact for all $N$}

$\dim\HL(10)=120$ and $\dim\TL(10)=42$. For $N\ge5$ the uncappable space $U(10)$ has dimension $13$, and the rotation
$\rho$ (of order $10$) splits it over $\Q(N)$ into three blocks: $\ker(\rho-1)$, $\ker(\rho+1)$ and $\ker\Phi_5(\rho)$.

\begin{remark}[what the constraint means]
An uncappable $x\in\HL(10)$, read as a map from $6$ strands to $4$ strands, takes values in the $O_N^+$-irreducible $r_4$.
Restricted to $O_N^*$, $r_4$ is multiplicity-free with $5$ constituents, the $PU_N$-irreducibles
$(2|2)$, $(11|11)$, $(2|11)$, $(11|2)$ and the adjoint $(1|1)$. For instance at $N=8$ we get $3905=1232+720+945+945+63$.
So $f_4\HL(8)f_4$ is commutative of dimension $5$, and Lemma~\ref{lem:glue} says exactly that $xx^*$ is a scalar on $r_4$:
$x$ must reach the five constituents with equal weights.
\end{remark}

\begin{lemma}[real parts]\label{lem:realpart}
Let $V\subset\HL(2m)$ be a $\rho$-invariant subspace defined over $\mathbb R$, with real basis $v_1,\dots,v_d$, and let
$x=\sum y_bv_b\in W\cap V_\C$. Write $y=u+iw$ with $u,w\in\mathbb R^d$. For $0\le r<2m$ and every crossing pairing $\tau$ at
level $2j<2m$, the real part of the $\tau$-coefficient of $\rho^rx\circ_{2m-j}J\rho^rx$ equals $S_{\tau,r}(u)+S_{\tau,r}(w)$,
where $S_{\tau,r}$ is the real symmetric form $S_{\tau,r}(u)=\mathrm{coeff}_\tau\big(\rho^r\textstyle\sum u_bv_b\circ_{2m-j}\mathrm{rev}(\rho^r\sum u_bv_b)\big)$.
So if some real combination $\sum c_{\tau,r}S_{\tau,r}$ is positive definite on $\mathbb R^d$, then $W\cap V_\C=0$.
\end{lemma}

\begin{proposition}[exact, all $N\ge5$]\label{prop:ten}
\begin{enumerate}
\item \emph{Blocks.} On $\ker\Phi_{10}(\rho)$ the caps are injective: a full-rank minor equals $4N^6(N^2+1)^2$. The block
$\ker(\rho-1)\cap U(10)$ has dimension $4$, $\ker(\rho+1)\cap U(10)$ has dimension $1$, and $\ker\Phi_5(\rho)\cap U(10)$ has dimension $8$.
The corresponding rank minors are $-2N^2(N^8-13N^6+62N^4-128N^2+96)$ (computed on the $17$ dihedral orbit sums), $-12N$ and $-8N^{11}$.
Each block has an explicit basis with entries in $\Z[N]$ of degree $\le9$, $\le3$ and $\le4$, and the bases are independent for $N\ne0$.
\item \emph{$\rho=-1$.} For the basis vector $x$, the first crossing coefficient of $x\circ_6x$ is
$-N(N^6-18N^4+52N^2-39)$, which never vanishes for $N\ge5$.
\item \emph{$\rho=1$.} Every rotation orbit in $P_2^*(10)$ is reversal-invariant, so $J$ acts by complex conjugation on this
block. Let $Q_0,\dots,Q_{10}$ be the real forms coming from the $10$ crossing pairings of $P_2^*(8)$, in our fixed
ordering, and from the crossing $\kappa$ of $P_2^*(6)$. Then
\begin{gather*}
\sum_kc_k(N)Q_k,\\
c(N)=(-3N^2,-3N^2,30N,20N,-3N^2,20N,-3N^2,30,-18N,-18N,-18),
\end{gather*}
is positive definite for all real $N\ge7$. Its leading principal minors are polynomials of degrees $21,42,57,70$, all of
whose coefficients are positive after the substitution $N\mapsto N+7$. For $N=5,6$ there are exact rational certificates.
\item \emph{$\Phi_5(\rho)=0$.} With the rotation-averaged forms $\bar Q_k=\sum_{r=0}^4S_{k,r}$, the combination with
\[
c(N)=(-2N,-2N,0,N^2,-2N,N^2,-2N,-N,-12,-12,0)
\]
is positive definite for all $N\ge5$. Its $8$ leading principal
minors have degrees $8,18,\dots,64$ and positive coefficients in $N-5$.
\end{enumerate}
\end{proposition}

\begin{proof}
All the data are polynomial in $N$ with explicit degree bounds: degree $\le1$ for the caps, the basis degrees above, and
at most $3$ loops in a gluing. Every identity and every minor was obtained by exact evaluation at more points than its
degree bound, and checked on further values. Positivity then follows from Descartes' rule of signs.
\end{proof}

\begin{theorem}[exact, level $10$, all $N\ge3$]\label{thm:ten}
Let $N\ge3$ and $O_N^*\subset G\subset O_N^+$. If $\mathcal C_G(8)=\TL(8)$, then $\mathcal C_G(10)=\TL(10)$.
\end{theorem}

\begin{proof}
For $N\ge5$: $W\subset U(10)$ is $\rho$-invariant, so it is the sum of its intersections with the three blocks of
Proposition~\ref{prop:ten}(1). Each intersection vanishes, by (2) for $\rho=-1$, and by Lemma~\ref{lem:realpart} together
with (3) and (4) for the other two blocks. For $N=3,4$ we work modulo the Gram kernel. There $\dim U(10)=6$ and $12$
respectively, and ``lies in $\TL(8)$'' is tested by Gram projection onto $\HL(8)\ominus\TL(8)$. For both values exact
rational positive-definite certificates were found, combining all rotations.
\end{proof}

\begin{corollary}
For $N\ge3$, every quantum group $O_N^*\subset G\subsetneq O_N^+$ whose fixed vectors differ from those of $O_N^+$
at some level $\le10$ is equal to $O_N^*$.
\end{corollary}

\begin{remark}
Two comments on the proof.
\begin{itemize}
\item The $\rho=1$ block is the delicate one. With a natural normalization, the best certificate has margin of
order $N^{-3}$, and its direction is asymptotically the difference of the forms attached to the two large
constituents of $r_4$. So at $N=\infty$ there is a direction in which all five weights agree to leading order. The
exclusion is a finite-$N$ phenomenon, made uniform by the $N$-dependent certificate $c(N)$.
\item The holomorphic constraints alone (gluings $x\circ x$, without $J$) do not suffice here. On the $\rho=1$ block they span
only $3$ of the $10$ quadratic forms in $4$ variables, and have $8$ common complex zeros. Positivity is essential, in
contrast with level $7$ for $S_N$ in Section~\ref{sec:method}.
\end{itemize}
\end{remark}

\subsection{Completeness: the Brauer version}

We now drop the assumption $G\supset O_N^*$. Intermediate quantum groups $O_N\subset G\subset O_N^+$ correspond to
categories $\TL\subset C\subset\mathrm{Br}$, where $\mathrm{Br}(k)=\spn\{\xi_\pi:\pi\in P_2(k)\}$, and $C\cap\HL$ is the
category of $\langle G,O_N^*\rangle$. The machinery of the previous subsections applies verbatim with $P_2^*$ replaced by $P_2$:
uncappable vectors, rotation blocks, the real-parts lemma, and Gram projections for small $N$. There is one new feature:
at level $6$ the answer is not ``$W=0$'', since $O_N^*$ itself first differs there.

\begin{proposition}[level $4$]\label{prop:brfour}
$P_2(4)=NC_2(4)\sqcup\{(13)(24)\}$. So if $C(4)\ne\TL(4)$, then $C$ contains the basic crossing and $G=O_N$.
\end{proposition}

\begin{proposition}[level $6$]\label{prop:brsix}
Let $N\in\{3,4,5,6,7,8,9,10,12,16\}$ and $C(4)=\TL(4)$. Then $C(6)\subset\HL(6)$ (\texttt{pairlevels.py}). Consequently, if $C(6)\ne\TL(6)$, then $C\supset\HL$, so
$G\subset O_N^*$, and $G\in\{O_N,O_N^*\}$ by~\cite{bbcc}.
\end{proposition}

\begin{proof}
The uncappable space in $\mathrm{Br}(6)$ has dimension $4$. It splits into the $\rho$-invariant block, of dimension $2$, and
the block of eigenvalues $e^{\pm2\pi i/3}$, of real dimension $2$. The latter is excluded by a positive-definite certificate.
On the $\rho$-invariant block, the only constraint from lower levels is the crossing coefficient $q(x)$ of
$x\circ_4x\in\TL(4)$. Exactly, this quadric has two zero lines, $\C h$ and $\C y$.
\begin{itemize}
\item $h$ is supported on $P_2^*(6)$: it is the uncappable part of $\xi_\kappa$.
\item $y$ is not supported on $P_2^*(6)$, and it fails closure: the projection of $y\circ_3y$ onto $\mathrm{Br}(6)\ominus\TL(6)$
is not proportional to $y$.
\end{itemize}
Every vector of $W$ satisfies $q=0$, so $W\cap(\rho\text{-invariant block})$ is contained in one of the two lines. If it were
$\C y$, then $W=\C y$, and closure would require $y\circ_3y\in\TL(6)\oplus\C y$. Hence $W\subset\C h$, and $C(6)\subset\TL(6)+\C h=\HL(6)$.
\end{proof}

\begin{proposition}[levels $8$, $10$ and $12$, numerical]\label{prop:breightten}
Let $C$ first differ from $\TL$ at level $2m\in\{8,10,12\}$. Then all rotation blocks of the uncappable space of $\mathrm{Br}(2m)$
are excluded by positive-definite certificates, using the deepest gluing $j=m-1$. This was checked at
\begin{itemize}
\item level $8$, where $\dim U=31$ (and $17$ at $N=3$), for $N=3,4,5,6,7,8,9,10,12,16$;
\item level $10$, where $\dim U=293$ and the blocks have real dimensions up to $68$, for $N=5,6,7,8,9,10,12$; and with Gram projections at $N=3$ ($\dim U=91$) and $N=4$ ($\dim U=251$). At $N=16$ the rotation-invariant block is left with a positive but uncertified value, as discussed below.
\item level $12$, where $\dim U=3326$ and the blocks have real dimensions up to $588$, for $N=6,7,8$ (\texttt{evander.py}).
\end{itemize}
\end{proposition}

\begin{theorem}[completeness through $10$ and $12$ points, numerical]\label{thm:complete}
Let $3\le N\le10$ or $N=12$. Then every compact quantum group $O_N\subset G\subset O_N^+$ whose fixed vectors differ from those of
$O_N^+$ at some level $\le10$ is equal to $O_N$ or to $O_N^*$. For $N\in\{6,7,8\}$ the same holds at some level $\le12$.
\end{theorem}

\begin{proof}
Let $2m$ be the level at which $G$ first differs from $O_N^+$. If $2m=4$, then $G=O_N$. If $2m=6$, use Proposition~\ref{prop:brsix}. Levels $8$,
$10$, and $12$ where stated, are impossible, by Proposition~\ref{prop:breightten}.
\end{proof}

\begin{remark}
Two comments on this result.
\begin{itemize}
\item The case analysis is natural. Level $4$ detects $O_N$, level $6$ detects $O_N^*$, and from level $8$ on nothing
survives. The level-$6$ step is the only place where closure, rather than positivity, is needed. It is the orthogonal
counterpart of Step~3 in Section~\ref{sec:method}.
\item The Brauer problem is much larger than the half-liberated one: $|P_2(10)|=945$ against $5!=120$, and $293$ uncappable
vectors against $13$. Exactness for all $N$ is within reach at levels $6$ and $8$ by the method of Theorem~\ref{thm:ten}.
\end{itemize}
\end{remark}

\subsection{Computational results}

``Excluded'' means that a positive-definite
combination was found on every eigenspace. We used only the deepest gluing $j=m-1$: it suffices in all cases,
while $j=3$ alone does not.

\begin{center}
\begin{tabular}{@{}cccl@{}}
\toprule
level $2m$ & $\dim\HL$ & $\dim U(2m)$ ($N\ge m$) & $N$ tested: all excluded \\
\midrule
6 & 6 & 1 & all $N\ge3$ (Prop.~\ref{prop:6}) \\
8 & 24 & 2 & all $N\ge3$, exact (Thm.~\ref{thm:eight}) \\
10 & 120 & 13 & all $N\ge3$, exact (Thm.~\ref{thm:ten}) \\
12 & 720 & 80 & $3,4,5,6,8,12$ \\
14 & 5040 & 579 & $7,8,9,10,12$ \\
\bottomrule
\end{tabular}
\end{center}

For $N<m$ the dimensions are smaller. For example, at $N=3$ and level $12$, $\dim\HL=513$ and $\dim U=23$. These small
cases lie well beyond the range $k\le 2N$ where Brauer relations start, and they are excluded just as cleanly.

\begin{theorem}[numerical]
For the values of $N$ in the table, every quantum group $O_N^*\subset G\subset O_N^+$ whose fixed-vector
spaces first differ from those of $O_N^+$ at $\le 2m_{\max}$ points is equal to $O_N^*$. Here $2m_{\max}=14$ for
$N\in\{7,8,9,10,12\}$, and $2m_{\max}=12$ for $N\in\{3,4,5,6\}$. (Up to level $10$ this holds for all $N\ge3$, exactly.)
\end{theorem}

Some caveats and questions are in order. Levels $8$ and $10$ are exact for all $N\ge3$, while the certificates at levels $12$ and $14$ were computed in floating point, at sampled values of $N$, and at level $14$ with $40$--$60$ random real combinations of the $588$ crossing coordinates instead of all of them, which is still a valid sufficient condition; at $N=16,20$ the rotation-invariant block is left there with a best value which is positive but of order $10^{-6}$, below what floating point can certify. The same happens in the Brauer problem at level $12$, where $|P_2(12)|=10395$ and $\dim U=3326$ at $N=8$: all seven rotation blocks are excluded at $N=6,7,8$, with smallest margins of order $10^{-5}$, which gives the $12$-point statement of Theorem~\ref{thm:complete}, while at $N=9,10,12,16$ the block $k=0$, of dimension $300$, is left with a best value between $6\cdot10^{-6}$ and $4\cdot10^{-8}$, decreasing with $N$ (\texttt{evander.py}). As at level $10$, the rotation-invariant block is the delicate one, with a margin vanishing as $N\to\infty$. The main questions are to make all this exact, for all $N$, by the method of Theorem~\ref{thm:ten}, with polynomial bases, minors and $N$-dependent certificates, first at the Brauer levels $6$ and $8$, and at the half-liberated level $12$, which is better conditioned than level $10$, with margins of order $1/N$ against $N^{-3}$, and then to reach level $16$, where $|P_2^*(16)|=8!=40320$, the gluing tables become too large, and one should rather work in the $\C[S_m]$ picture of Subsection~\ref{sec:shift}.

\subsection{Comparison with Section~\ref{sec:method}: a shift by one}\label{sec:shift}

The two problems are of the same nature. Under Jones' fattening~\cite{jo2}, $NC(K)$ corresponds to $NC_2(2K)$, and the operations
by noncrossing partitions become the Temperley--Lieb operations. The ambient spaces differ, though:
\begin{itemize}
\item for $S_N$ it is $\spn\mathcal P(K)$;
\item for $O_N^*$ it is $\HL(2K)\simeq\C[S_K]$, by Schur--Weyl duality.
\end{itemize}
A partition becomes the permutation whose cycles are its blocks, so $\mathcal P(K)\subset S_K$, and the permutation
problem sits inside the half-liberated one, with fewer extra elements. Counting the extra basis elements gives a striking
shift by one level.

\begin{center}
\begin{tabular}{@{}cccccc@{}}
\toprule
\multicolumn{3}{c}{$S_N\subset S_N^+$, level $K$} & \multicolumn{3}{c}{$O_N^*\subset O_N^+$, level $2(K-1)$}\\
\cmidrule(r){1-3}\cmidrule(l){4-6}
$K$ & $B_K-C_K$ & first-step kernel & $2(K-1)$ & $(K-1)!-C_{K-1}$ & $\dim U$ \\
\midrule
$4$ & $1$ & $1$ & $6$ & $1$ & $1$ \\
$5$ & $10$ & $0$ & $8$ & $10$ & $2$ \\
$6$ & $71$ & $5$ & $10$ & $78$ & $13$ \\
$7$ & $448$ & $14$ & $12$ & $588$ & $80$ \\
$8$ & $2710$ & $66$ & $14$ & $4611$ & $579$ \\
\bottomrule
\end{tabular}
\end{center}

Here $B_K$ and $C_K$ are the Bell and Catalan numbers, ``first-step kernel'' is $\dim K_K(N)$ from Table~\ref{tab:sizes}, and
$\dim U$ is the dimension of the uncappable space, for large $N$. So $S_N$ at level $K+1$ behaves like $O_N^*$ at level
$2K$.
\begin{itemize}
\item The trivial levels match: $S_N$ at level $4$ (one basic crossing) and $O_N^*$ at level $6$ (one half-commutation
crossing).
\item The numbers of extra elements agree exactly at $S_N$ level $5$ and $O_N^*$ level $8$, both $10$, and they stay close
afterwards. The reason is that $\mathcal P(K)$ contains every permutation for $K\le3$, while $S_3$ already contains a
permutation which is not a partition.
\item The steps match too. Step~1 of Section~\ref{sec:method} (leg-sums and merges landing in $\mathcal{NC}$) is the uncappable step
here. Step~2 (quadrics from ribbons) is Lemma~\ref{lem:glue}. Step~3 (closure) appears only in the Brauer problem at
level $6$ (Proposition~\ref{prop:brsix}).
\item At matching levels the first step is stronger for $S_N$, because singletons and $3$-blocks give more linear
constraints. This is why $S_N$ level $5$ dies by a single determinant, while $O_N^*$ level $8$ keeps a $2$-dimensional
uncappable space and needs one quadric.
\item Here we use positivity, through the $C^*$-structure ($x\circ Jx=xx^*$). This is essential at level $10$, where the
holomorphic quadrics alone leave complex solutions, while Sections~\ref{sec:method}--\ref{sec:higher} never use it.
\item As in Section~\ref{sec:method}, the next level is better conditioned than the hard one. Level $12$ has positivity margins of
order $1/N$ on its tight block, against $N^{-3}$ at level $10$, just as $S_N$ level $7$ was easier than level $6$.
\end{itemize}



%======================================================================
% Section: related
%======================================================================

\section{Projective versions, seeds and levels}\label{sec:related}

We discuss the projective version of the orthogonal problem, namely the intermediate quantum groups
$PO_N\subset H\subset PO_N^+$. Projective easiness~\cite{bme} says nothing about these beyond
the orthogonal problem, and the half-liberated part of the question is the maximality of $PU_N\subset PO_N^+$. This section mostly sets up problems and collects evidence.

\subsection{Projective versions}

For a closed subgroup $G\subset O_N^+$, the projective version $PG$ is given by the subalgebra of $C(G)$ generated by the
entries $u_{ia}u_{jb}$ of $v=u\otimes u$. We have $PO_N^*=PU_N$ and $PO_N^+=PU_N^+$. The intermediate objects
\[
PO_N\subset H\subset PO_N^+
\]
are described by their fixed vectors $C_H(m)=\Fix_H(v^{\otimes m})\subset(\C^N)^{\otimes2m}$, with
$\TL(2m)\subset C_H(m)\subset\Br(2m)$. These spaces are closed under tensor products, under all Temperley--Lieb maps between
even numbers of strands (in particular the caps at positions $i,i+1$ for $1\le i\le 2m-1$), and under the rotation $\rho^2$ by
two points. Since $\bar v=FvF$, with $F$ the flip inside each pair, and since taking adjoints reverses the order of products,
they are also closed under the antilinear involution
\[
x\mapsto F^{\otimes m}\,\mathrm{rev}_m(\bar x)=\overline{\mathrm{rev}(x)}=Jx,
\]
where $\mathrm{rev}_m$ reverses the order of the $m$ pairs, and $J$ is the involution of Section~\ref{sec:orthproblem}. For a classical group one could
drop $\mathrm{rev}_m$, but not in general, $\mathrm{rev}_m$ not being among the available operations.
What is missing, compared with Section~\ref{sec:pairings}, is the rotation $\rho$ by one point, and with it the cyclic cap.

\begin{proposition}\label{prop:inj}
The map $G\mapsto PG$, from intermediate quantum groups $O_N\subset G\subset O_N^+$ to intermediate quantum groups
$PO_N\subset H\subset PO_N^+$, is injective. Its image consists of the $H$ whose categories are closed under $\rho$.
\end{proposition}

\begin{proof}[Sketch]
Since $-1\in O_N\subset G$, the odd levels of $\mathcal C_G$ vanish, and $\mathcal C_G(2m)=\mathcal C_{PG}(m)$. So $\mathcal C_G$, hence $G$, is determined by
$PG$, and $\mathcal C_{PG}$ is $\rho$-invariant. Conversely, if $C_H$ is closed under $\rho$, then $C(2m)=C_H(m)$ and $C(2m+1)=0$
define a family of fixed-vector spaces closed under tensor products, adjacent contractions, rotation and $J$, the last
one being the conjugation axiom, as noted above. Hence this is the category of some $O_N\subset G\subset O_N^+$, and $PG=H$.
\end{proof}

So the problem splits in two:
\begin{enumerate}
\item \emph{Liftable $H$}, i.e.\ $H=PG$. These are governed by Sections~\ref{sec:orthproblem} and~\ref{sec:pairings}: conjecturally $H\in\{PO_N,PU_N,PO_N^+\}$.
\item \emph{Non-liftable $H$}. These are new. Their categories are not closed under one-point rotation.
\end{enumerate}

\begin{remark}[projective easiness]
Banica and M\'esz\'aros define projective categories of pairings $NC_2(2k,2l)\subset E(k,l)\subset P_2(2k,2l)$, with the
axiom $\sigma\in E\Rightarrow|\sigma|\in E$, and show that $PU_N$ is the only intermediate projectively easy
$PO_N\subset H\subset PO_N^+$~\cite[Theorem~4.6]{bme}. Their axiom is the easy version of closure under $\rho$, and their
Proposition~4.5 is the easy version of Proposition~\ref{prop:inj}. So their theorem is the affine classification of~\cite{bcs},
transported. Dropping the axiom gives the natural question: classify the ``weak'' projective categories of pairings,
closed under the operations above but not necessarily under $\rho$. This is a finite combinatorial problem at each level.
\end{remark}

\begin{proposition}[$N=2$]
There is no quantum group strictly between $PO_2$ and $PO_2^+$, other than $PU_2$.
\end{proposition}

\begin{proof}
Since $O_2^*=O_2^+$ (Remark~\ref{rem:ntwo}), we have $PO_2^+=PO_2^*=PU_2=SO_3$, a classical group, and $PO_2\subset PU_2$ is maximal
by~\cite{bbcc}.
\end{proof}

\subsection{The inclusion $PU_N\subset PO_N^+$}

Since $PO_N^*=PU_N$ and $PO_N^+=PU_N^+$, the projective chain reads $PO_N\subset PU_N\subset PO_N^+$, and the inclusion $PO_N\subset PU_N$ is maximal~\cite{bbcc}. The other half is the following.

\begin{conjecture}\label{conj:mn}
The inclusion $PU_N\subset PO_N^+$ is maximal.
\end{conjecture}

Some remarks on this statement:
\begin{itemize}
\item For liftable $H\supset PU_N$, Conjecture~\ref{conj:mn} is the half-liberated problem of Section~\ref{sec:pairings}, which is solved there through level $10$ for all $N\ge3$.
\item In the level picture, the problem lives inside $\HL(2m)\simeq\C[S_m]$ (Remark~\ref{rem:unitarypic}), but without the one-point rotation. This is related to the ``shift by one'' of Section~\ref{sec:shift}.
\item Natural non-classical quantum subgroups of $PO_N^+$, not containing $PU_N$, come from the Weyl models of Banica and Nechita~\cite{bn}, already met in Section~\ref{sec:twists}.
\end{itemize}

\begin{remark}[a quantum set interpretation]
One can view $PU_N=\Aut(M_N)$ and $PO_N^+=\Aut^+(M_N,\mathrm{tr})$, the quantum automorphism group of the matrix algebra with its trace~\cite{ba,bqp,wa}. Conjecture~\ref{conj:mn} then reads as the analogue, for the ``quantum set'' $M_N$, of the maximality of $S_N=\Aut(\C^N)$ in $S_N^+=\Aut^+(\C^N)$, and $\Aut^+(M_N,\mathrm{tr})$ has the same Temperley--Lieb category as $S_{N^2}^+$, at loop parameter $N$. This is a pleasant reformulation, but it brings nothing new to the problem.
\end{remark}

\subsection{Seeds, projective case}

As in Proposition~\ref{prop:seeds}, by the maximality of $PO_N\subset PU_N$, completeness of
$PO_N\subset PU_N\subset PO_N^+$ is equivalent to $\langle PO_N,K\rangle=PO_N^+$ for every closed $K\subset PO_N^+$ with
$K\not\subset PU_N$. The new seeds are those $K$ which are not projective versions of subgroups of $O_N^+$. Natural
candidates:
\begin{itemize}
\item projective versions $PG$ of $G\subset U_N^+$. The unitary group duals of Section~\ref{sec:gd}, given by pairs $(\Gamma,M)$,
belong here. We checked $119$ law sets at $N=3$ and $611$ at $N=4$, levels $4,6,8$ (computational, \texttt{projective.py}), with the outcome:
$PO_N^+$ when $PG$ is non-classical, and $PO_N$ or $PU_N$ when $PG$ is classical, $56$ times out of $611$ at $N=4$;
\item quantum automorphism groups of quantum graphs on $M_N$;
\item Weyl models, as in Section~\ref{sec:flat};
\item projective versions of the twists of Sections~\ref{sec:twist}--\ref{sec:othertwists}, and the corners of $PO_N^+$: the corner method of Section~\ref{sec:twist} (Lemma~\ref{lem:corner}, Proposition~\ref{prop:chir}) should transport to the
projective setting, with $PO_{N-1}^+$ embedded as a corner of $PO_N^+$.
\end{itemize}

\subsection{Levels, projective case}\label{sec:projlevels}

We now transport the level method of Section~\ref{sec:pairings}. The first question is which operations the categories $C_H(m)=\Fix_H(v^{\otimes m})\subset(\C^N)^{\otimes 2m}$ are closed under. Since $v=u\otimes u$ and $u^{\otimes 2a}=v^{\otimes a}$ as matrices over $C(PO_N^+)$, every Temperley--Lieb diagram between $2a$ and $2b$ points is an intertwiner of $H$. So $C_H$ is closed under tensor products, under all adjacent caps and cups, at any position, under the rotation $\rho^2$ by two points, and under $J$. What is missing is the one-point rotation $\rho$, and with it the cyclic cap: conjugating the cap at $(2m-2,2m-1)$ by $\rho^2$ gives the cyclic cap composed with $\rho$, which kills the same vectors. We call a family of subspaces $\TL(2m)\subset C(m)\subset\mathrm{Br}(2m)$ closed under these operations a \emph{$\rho^2$-category}. The categories $C_H$ are exactly the $\rho^2$-categories, by Tannaka duality for $PO_N^+$, and the liftable ones are those closed under $\rho$, by Proposition~\ref{prop:inj}. The method then goes as before, with one change: $W=C(m)\ominus\TL(2m)$ is only $\rho^2$-invariant, so the rotation blocks of eigenvalues $\omega$ and $-\omega$ merge, and only even rotations of the quadratic forms may be used in the certificates. The uncappable space $U(2m)$ is that of the Brauer problem.

\begin{proposition}[level $4$]
If $C_H(2)\neq\TL(4)$ then $H\in\{PO_N,PU_N\}$.
\end{proposition}
\begin{proof}
$\mathrm{Br}(4)=\TL(4)\oplus\C\xi_{(13)(24)}$, and $\xi_{(13)(24)}$ is the flip of the two copies of $v$. So $v_{ij}v_{kl}=v_{kl}v_{ij}$, $H$ is classical, hence $H\subset PU_N$, and $PO_N\subset PU_N$ is maximal~\cite{bbcc}.
\end{proof}

At level $6$ the extra vector of $\HL(6)$ appears, and the question is whether it forces $H=PU_N$, i.e.\ whether $\xi_\kappa$ generates $\HL$ without the help of $\rho$. Note that $\kappa=(14)(25)(36)$ is $\rho$-invariant, but this does not make the generated $\rho^2$-category $\rho$-invariant, since the tensor product does not commute with $\rho$.

\begin{proposition}[generation, computational]\label{prop:gen}
The $\rho^2$-category generated by $\TL$ and $\xi_\kappa$ coincides with $\HL$ at levels $2m\le 12$: its dimensions are $1,2,6,24,120,720$, against $|P_2^*(2m)|=m!$, at $N=5,7$ (\texttt{projlevels.py}). Consequently, if $C_H(3)\neq\TL(6)$, then $C_H(m)\supset\HL(2m)$ for $m\le6$.
\end{proposition}

The closure is computed within the pairing coordinates of $\mathrm{Br}(2m)$, by iterating tensor products, cups, caps, gluings, $\rho^2$ and $J$ until the dimensions stabilize, which takes four rounds through level $10$, and six through level $12$. The generated spaces never leave $\HL$, as they must.

\begin{proposition}[level $6$]\label{prop:projsix}
Let $N\in\{3,4,5,6,7,8,9,10,12,16\}$ and $C_H(2)=\TL(4)$. Then $C_H(3)\subset\HL(6)$.
\end{proposition}
\begin{proof}
The eigenvalues of $\rho$ on $U(6)$ are $1$, twice, and $e^{\pm2\pi i/3}$ (\texttt{projlevels.py}). So the eigenspaces of $\rho^2$ on $U(6)$ are those of $\rho$, and on the second block $\rho=\rho^4$ is an even rotation. Hence the proof of Proposition~\ref{prop:brsix} applies word for word: the second block is excluded by its certificate, and on the invariant block the quadric has the two zero lines $\C h$ and $\C y$, with $y$ failing closure.
\end{proof}

From level $8$ on, we use the certificates of Section~\ref{sec:pairings} on the merged blocks, for the two problems $PU_N\subset H$ (half-liberated pairings) and $PO_N\subset H$ (all pairings). The half-liberated level $8$ is exact.

\begin{proposition}[level $8$, half-liberated, exact for all $N\ge3$]\label{prop:projeight}
Let $PU_N\subset H\subset PO_N^+$ with $C_H(3)=\TL(6)$. Then $C_H(4)=\TL(8)$.
\end{proposition}
\begin{proof}
Here $U(8)=\C x_+\oplus\C x_-$, the two vectors of Proposition~\ref{prop:eightpair}, with $\rho x_\pm=\pm x_\pm$, and they merge into one block. Both are $J$-real, so the $J$-real vectors are $x=\alpha x_++\beta x_-$ with $\alpha,\beta$ real, and the $\kappa$-coefficient of $x\circ_5Jx$ is $\alpha^2h_+(N)+\beta^2h_-(N)+2\alpha\beta h_{+-}(N)$. The cross term $h_{+-}(N)$, the $\kappa$-coefficient of $x_+\circ_5Jx_-$, vanishes identically in $N$ (exactly, at $N=3,\dots,20$, degree bound $8$; \texttt{projeight.py}). So the quadric is $\mathrm{diag}(h_+,h_-)$, negative definite for $N\ge4$ by the formulas of Proposition~\ref{prop:eightpair}, and for $N=3$ the argument of Theorem~\ref{thm:eight} applies, $x_-$ being the Gram relation.
\end{proof}

The cases not marked exact are numerical, with the random crossing combinations and the positivity search of \texttt{evander.py}, families \texttt{PH} and \texttt{PR}. ``Excluded'' means a positive-definite combination on every merged block.

\begin{center}
\begin{tabular}{@{}clccll@{}}
\toprule
problem & $2m$ & $\dim U$ & merged blocks & $N$ & status \\
\midrule
$PU_N\subset H$ & 6 & 1 & $1$ & all $N\ge3$ & $C_H(3)=\HL(6)$, Prop.~\ref{prop:gen} \\
& 8 & 2 & $2$ & all $N\ge3$ & excluded, exact (Prop.~\ref{prop:projeight}) \\
& 10 & 13 & $5,4,4$ & all $N\ge5$ & excluded, exact (\texttt{projtencert.py}) \\
\midrule
$PO_N\subset H$ & 4 & -- & -- & all $N$ & $H\in\{PO_N,PU_N\}$ \\
& 6 & 4 & $2,2$ & $3$--$10,12,16$ & $C_H(3)\subset\HL(6)$ (Prop.~\ref{prop:projsix}) \\
& 8 & 31 & $13,8,10$ & $4$--$10,12,16$ & excluded, margins $1.7\cdot10^{-3}$ to $1.4\cdot10^{-5}$ \\
& 10 & 293 & $61,116,116$ & $5$--$10,12,16$ & not excluded, best $-3\cdot10^{-5}$ to $-3\cdot10^{-6}$ \\
\bottomrule
\end{tabular}
\end{center}

The dimensions are those of Section~\ref{sec:pairings}; the merged blocks are $\dim U_\omega+\dim U_{-\omega}$. At the half-liberated level $10$ the $\Phi_5$ blocks are unaffected, since the $\Phi_{10}$ blocks are empty, and $\rho^5=1$ there, so Proposition~\ref{prop:ten}(4) applies. Only $4+1$ merge, into a $5$-dimensional block with genuine cross terms between $\ker(\rho-1)$ and $\ker(\rho+1)$. Still, the $c(N)$ of Proposition~\ref{prop:ten}(3) works on it for $N\ge7$, with leading minors of degrees $21,36,57,70,77$, positive in $N-7$, and $N=5,6$ have integer certificates (\texttt{projtencert.py}, with $Q_2,Q_3$ and $Q_6,Q_7$ swapped). The Brauer level $10$ is the one new phenomenon: every $\rho$-block was excluded in Proposition~\ref{prop:breightten}, but no common certificate exists on the merged blocks, at any $N$ tried. A vector $x_1+x_{-1}$, with $x_{\pm1}$ in the two $\rho$-blocks, can satisfy the quadratic constraints by cancellation between the blocks, so merging weakens the degree-$2$ certificate. Whether the best values near $0$ signal a weakness of the method or a non-liftable $H$ at $10$ points is open, and the same-level gluing modulo $W$ and the cubic gluings are the next constraints to add.

\begin{question}\label{qu:projective}
Develop the projective version of the level method further. More precisely:
\begin{enumerate}
\item determine whether the categories $C_H$ have symmetries beyond tensor products, caps, cups, $\rho^2$ and $J$;
\item prove that $\xi_\kappa$ generates $\HL$ under these operations at all levels, as Proposition~\ref{prop:gen} suggests;
\item settle the Brauer level $10$, with the same-level gluing modulo $W$ or the cubic gluings, and make the half-liberated level $10$ exact at $N=3,4$, modulo the Gram kernel;
\item clarify the literature: is the maximality of $PU_N\subset PU_N^+$ known, and is every $H\subset PU_N^+$ the projective version of some $G\subset U_N^+$?
\end{enumerate}
\end{question}

\begin{theorem}[projective versions, numerical]\label{thm:projective}
Let $PO_N\subset H\subset PO_N^+$, with fixed-vector spaces $C_H(m)\subset(\C^N)^{\otimes 2m}$.
\begin{enumerate}
\item If $PU_N\subset H$ and $H$ first differs from $PO_N^+$ at $2m\le10$ points, then $2m=6$ and $C_H(m)=\HL(2m)$ for $m\le6$: $H$ coincides with $PU_N$ through $12$ points. This holds exactly, for all $N\ge3$ at $2m\le8$, and for all $N\ge5$ at $2m=10$.
\item If $N\in\{4,\dots,10,12,16\}$ and $H$ first differs from $PO_N^+$ at $2m\le8$ points, then either $H\in\{PO_N,PU_N\}$, or $2m=6$ and $C_H(m)\supset\HL(2m)$ for $m\le6$: $H$ coincides with $PU_N$ through $12$ points.
\end{enumerate}
\end{theorem}
\begin{proof}
(1) Level $4$ is impossible, since $PU_N$ is not classical. Level $6$ forces $C_H(3)=\HL(6)$, and Proposition~\ref{prop:gen} gives $C_H(m)\supset\HL(2m)$ for $m\le6$, while $C_H\subset\HL$ since $H\supset PU_N$. Level $8$ is excluded by Proposition~\ref{prop:projeight}, and level $10$ by the table. (2) Level $4$ gives $H\in\{PO_N,PU_N\}$, level $6$ gives $C_H(3)=\HL(6)$ by Proposition~\ref{prop:projsix}, then Proposition~\ref{prop:gen}, and level $8$ is excluded by the table.
\end{proof}

\begin{thebibliography}{99}
\raggedright
\hypersetup{urlcolor=black}

\bibitem{ba} T.~Banica, Symmetries of a generic coaction, \emph{Math.\ Ann.}\ \textbf{314} (1999), 763--780.

\bibitem{bah} T.~Banica, Homogeneous quantum groups and their easiness level, \emph{Kyoto J.\ Math.}\ \textbf{61} (2021), 171--205.

\bibitem{bha} T.~Banica, Invitation to Hadamard matrices, book, \href{https://arxiv.org/pdf/1910.06911}{arXiv:1910.06911} (2024).

\bibitem{bqp} T.~Banica, Quantum permutation groups, book, \href{https://arxiv.org/pdf/2012.10975}{arXiv:2012.10975} (2024).

\bibitem{bb} T.~Banica and J.~Bichon, Quantum groups acting on $4$ points, \emph{J.\ Reine Angew.\ Math.}\ \textbf{626} (2009), 75--114.

\bibitem{bbrw} T.~Banica and J.~Bichon, Random walk questions for linear quantum groups, \emph{Int.\ Math.\ Res.\ Not.\ IMRN} \textbf{2015} (2015), 13406--13436.

\bibitem{bbc} T.~Banica, J.~Bichon and B.~Collins, The hyperoctahedral quantum group, \emph{J.\ Ramanujan Math.\ Soc.}\ \textbf{22} (2007), 345--384.

\bibitem{bbcc} T.~Banica, J.~Bichon, B.~Collins and S.~Curran, A maximality result for orthogonal quantum groups, \emph{Comm.\ Algebra} \textbf{41} (2013), 656--665.

\bibitem{bbcu} T.~Banica, J.~Bichon, S.~Curran, Quantum automorphisms of twisted group algebras and free hypergeometric laws, \emph{Proc.\ Amer.\ Math.\ Soc.}\ \textbf{139} (2011), 3961--3971.

\bibitem{bbn} T.~Banica, J.~Bichon and S.~Natale, Finite quantum groups and quantum permutation groups, \emph{Adv.\ Math.}\ \textbf{229} (2012), 3320--3338.

\bibitem{bbs} T.~Banica, J.~Bichon and J.-M.~Schlenker, Representations of quantum permutation algebras, \emph{J.\ Funct.\ Anal.}\ \textbf{257} (2009), 2864--2910.

\bibitem{bc} T.~Banica and B.~Collins, Integration over the Pauli quantum group, \emph{J.\ Geom.\ Phys.}\ \textbf{58} (2008), 942--961.

\bibitem{bcs} T.~Banica, S.~Curran and R.~Speicher, Classification results for easy quantum groups, \emph{Pacific J.\ Math.}\ \textbf{247} (2010), 1--26.

\bibitem{bme} T.~Banica and S.~M\'esz\'aros, Uniqueness results for noncommutative spheres and projective spaces, \emph{Illinois J.\ Math.}\ \textbf{59} (2015), 219--233.

\bibitem{bn} T.~Banica and I.~Nechita, Flat matrix models for quantum permutation groups, \emph{Adv.\ in Appl.\ Math.}\ \textbf{83} (2017), 24--46.

\bibitem{bsp} T.~Banica and R.~Speicher, Liberation of orthogonal Lie groups, \emph{Adv.\ Math.}\ \textbf{222} (2009), 1461--1501.

\bibitem{bv} T.~Banica and R.~Vergnioux, Invariants of the half-liberated orthogonal group, \emph{Ann.\ Inst.\ Fourier} \textbf{60} (2010), 2137--2164.

\bibitem{bi00} J.~Bichon, Quelques nouvelles d\'eformations du groupe sym\'etrique, \emph{C.\ R.\ Acad.\ Sci.\ Paris S\'er.\ I Math.}\ \textbf{330} (2000), 761--764.

\bibitem{bi04} J.~Bichon, Free wreath product by the quantum permutation group, \emph{Algebr.\ Represent.\ Theory} \textbf{7} (2004), 343--362.

\bibitem{bi} J.~Bichon, Algebraic quantum permutation groups, \emph{Asian-Eur.\ J.\ Math.}\ \textbf{1} (2008), 1--13.

\bibitem{bdv} J.~Bichon and M.~Dubois-Violette, Half-commutative orthogonal Hopf algebras, \emph{Pacific J.\ Math.}\ \textbf{263} (2013), 13--28.

\bibitem{bcf} M.~Brannan, A.~Chirvasitu and A.~Freslon, Topological generation and matrix models for quantum reflection groups, \emph{Adv.\ Math.}\ \textbf{363} (2020), 106982.

\bibitem{bcv} M.~Brannan, B.~Collins and R.~Vergnioux, The Connes embedding property for quantum group von Neumann algebras, \emph{Trans.\ Amer.\ Math.\ Soc.}\ \textbf{369} (2017), 3799--3819.

\bibitem{cha1} A.~Chassaniol, Quantum automorphism group of the lexicographic product of finite regular graphs, \emph{J.\ Algebra} \textbf{456} (2016), 23--45.

\bibitem{cha2} A.~Chassaniol, Study of quantum symmetries for vertex-transitive graphs using intertwiner spaces, \href{https://arxiv.org/pdf/1904.00455}{arXiv:1904.00455}.

\bibitem{ch1} A.~Chirvasitu, Residually finite quantum group algebras, \emph{J.\ Funct.\ Anal.}\ \textbf{268} (2015), 3508--3533.

\bibitem{ch2} A.~Chirvasitu, Topological generation results for free unitary and orthogonal groups, \emph{Internat.\ J.\ Math.}\ \textbf{31} (2020), 2050003.

\bibitem{cdpr} L.S.~Cirio, A.~D'Andrea, C.~Pinzari and S.~Rossi, Connected components of compact matrix quantum groups and finiteness conditions, \emph{J.\ Funct.\ Anal.}\ \textbf{267} (2014), 3154--3204.

\bibitem{dpr} A.~D'Andrea, C.~Pinzari and S.~Rossi, Polynomial growth of discrete quantum groups, topological dimension of the dual and $*$-regularity of the Fourier algebra, \emph{Ann.\ Inst.\ Fourier} \textbf{67} (2017), 2003--2027.

\bibitem{di} P.~Di\c{t}\u{a}, Some results on the parametrization of complex Hadamard matrices, \emph{J.\ Phys.\ A} \textbf{37} (2004), 5355--5374.

\bibitem{doi} Y.~Doi, Braided bialgebras and quadratic bialgebras, \emph{Comm.\ Algebra} \textbf{21} (1993), 1731--1749.

\bibitem{dr} V.G.~Drinfeld, Quasi-Hopf algebras, \emph{Leningrad Math.\ J.}\ \textbf{1} (1990), 1419--1457.

\bibitem{drs} J.~van Dobben de Bruyn, D.E.~Roberson and S.~Schmidt, Asymmetric graphs with quantum symmetry, \emph{Proc.\ Lond.\ Math.\ Soc.}\ \textbf{131} (2025), e70098.

\bibitem{ev} M.~Enock and L.~Vainerman, Deformation of a Kac algebra by an abelian subgroup, \emph{Comm.\ Math.\ Phys.}\ \textbf{178} (1996), 571--596.

\bibitem{eg} P.~Etingof and S.~Gelaki, A method of construction of finite-dimensional triangular semisimple Hopf algebras, \emph{Math.\ Res.\ Lett.}\ \textbf{5} (1998), 551--561.

\bibitem{fs} A.~Freslon and R.~Speicher, in preparation (2026).

\bibitem{ha} U.~Haagerup, Orthogonal maximal abelian $*$-subalgebras of the $n\times n$ matrices and cyclic $n$-roots, in \emph{Operator algebras and quantum field theory}, Int.\ Press (1997), 296--322.

\bibitem{hal} M.~Hall, Jr., Hadamard matrices of order $20$, Technical Report 32-761, Jet Propulsion Laboratory, Pasadena (1965).

\bibitem{imps} M.~Izumi, S.~Morrison, D.~Penneys, E.~Peters and N.~Snyder, Subfactors of index exactly $5$, \emph{Bull.\ Lond.\ Math.\ Soc.}\ \textbf{47} (2015), 257--269.

\bibitem{js} V.K.~Jain and R.P.~Shukla, Finite just non-Dedekind groups, arXiv:0705.1618 (2007).

\bibitem{jo1} V.F.R.~Jones, Index for subfactors, \emph{Invent.\ Math.}\ \textbf{72} (1983), 1--25.

\bibitem{jo2} V.F.R.~Jones, Planar algebras, I, \emph{New Zealand J.\ Math.}\ \textbf{52} (2021), 1--107.

\bibitem{jms} V.F.R.~Jones, S.~Morrison and N.~Snyder, The classification of subfactors of index at most $5$, \emph{Bull.\ Amer.\ Math.\ Soc.}\ \textbf{51} (2014), 277--327.

\bibitem{kr} A.~Klappenecker and M.~R\"otteler, Beyond stabilizer codes I: Nice error bases, \emph{IEEE Trans.\ Inform.\ Theory} \textbf{48} (2002), 2392--2395.

\bibitem{ks} C.~K\"ostler and R.~Speicher, A noncommutative de Finetti theorem: invariance under quantum permutations is equivalent to freeness with amalgamation, \emph{Comm.\ Math.\ Phys.}\ \textbf{291} (2009), 473--490.

\bibitem{los} P.H.J.~Lampio, P.R.J.~\"Osterg\aa rd and F.~Sz\"oll\H{o}si, Orderly generation of Butson Hadamard matrices, \emph{Math.\ Comp.}\ \textbf{89} (2020), 313--331.

\bibitem{lmr} M.~Lupini, L.~Man\v{c}inska and D.E.~Roberson, Nonlocal games and quantum permutation groups, \emph{J.\ Funct.\ Anal.}\ \textbf{279} (2020), 108592.

\bibitem{mw} A.~Mang and M.~Weber, Categories of two-colored pair partitions, Part I: Categories indexed by cyclic groups, \emph{Ramanujan J.}\ \textbf{53} (2020), 181--208.

\bibitem{mc0} J.P.~McCarthy, A state-space approach to quantum permutations, \emph{Expo.\ Math.}\ \textbf{40} (2022), 628--664.

\bibitem{mc1} J.P.~McCarthy, Tracing the orbitals of the quantum permutation group, \emph{Arch.\ Math.}\ \textbf{121} (2023), 211--224.

\bibitem{mc2} J.P.~McCarthy, Maximality levels of the classical permutation group in the quantum permutation group, \href{https://arxiv.org/pdf/2603.21759}{arXiv:2603.21759} (2026).

\bibitem{me} N.D.~Mermin, Simple unified form for the major no-hidden-variables theorems, \emph{Phys.\ Rev.\ Lett.}\ \textbf{65} (1990), 3373--3376.

\bibitem{ne1} M.F.~Newman, On a class of metabelian groups, \emph{Proc.\ London Math.\ Soc.}\ (3) \textbf{10} (1960), 354--364.

\bibitem{ne2} M.F.~Newman, On a class of nilpotent groups, \emph{Proc.\ London Math.\ Soc.}\ (3) \textbf{10} (1960), 365--375.

\bibitem{ni} D.~Nikshych, $K_0$-rings and twisting of finite dimensional semisimple Hopf algebras, \emph{Comm.\ Algebra} \textbf{26} (1998), 321--342.

\bibitem{ns} A.~Nica and R.~Speicher, Lectures on the combinatorics of free probability, Cambridge Univ.\ Press (2006).

\bibitem{nz} W.D.~Nichols and M.B.~Zoeller, A Hopf algebra freeness theorem, \emph{Amer.\ J.\ Math.}\ \textbf{111} (1989), 381--385.

\bibitem{pa} R.E.A.C.~Paley, On orthogonal matrices, \emph{J.\ Math.\ Phys.}\ \textbf{12} (1933), 311--320.

\bibitem{pe} M.~Petrescu, Existence of continuous families of complex Hadamard matrices of certain prime dimensions and related results, Ph.D.\ thesis, UCLA (1997).

\bibitem{rw2} S.~Raum and M.~Weber, Easy quantum groups and quantum subgroups of a semi-direct product quantum group, \emph{J.\ Noncommut.\ Geom.}\ \textbf{9} (2015), 1261--1293.

\bibitem{rw} S.~Raum and M.~Weber, The full classification of orthogonal easy quantum groups, \emph{Comm.\ Math.\ Phys.}\ \textbf{341} (2016), 751--779.

\bibitem{rs} D.E.~Roberson and S.~Schmidt, Solution group representations as quantum symmetries of graphs, \emph{J.\ Lond.\ Math.\ Soc.}\ \textbf{106} (2022), 3379--3410.

\bibitem{sz} J.~Schanz, Quantum symmetries of vertex-transitive graphs on $12$ vertices, \emph{J.\ Noncommut.\ Geom.}\ \textbf{20} (2026), 1029--1094.

\bibitem{sch} S.~Schmidt, The Petersen graph has no quantum symmetry, \emph{Bull.\ Lond.\ Math.\ Soc.}\ \textbf{50} (2018), 395--400.

\bibitem{sch2} S.~Schmidt, On the quantum symmetry groups of distance-transitive graphs, \emph{Adv.\ Math.}\ \textbf{368} (2020), 107150.

\bibitem{so} P.M.~So{\l}tan, Quantum families of maps and quantum semigroups on finite quantum spaces, \emph{J.\ Geom.\ Phys.}\ \textbf{59} (2009), 354--368.

\bibitem{szo} F.~Sz\"oll\H{o}si, Construction, classification and parametrization of complex Hadamard matrices, Ph.D.\ thesis, CEU Budapest (2011), \href{https://arxiv.org/pdf/1110.5590}{arXiv:1110.5590}.

\bibitem{tz} W.~Tadej and K.~\.{Z}yczkowski, A concise guide to complex Hadamard matrices, \emph{Open Syst.\ Inf.\ Dyn.}\ \textbf{13} (2006), 133--177.

\bibitem{tao} T.~Tao, Fuglede's conjecture is false in $5$ and higher dimensions, \emph{Math.\ Res.\ Lett.}\ \textbf{11} (2004), 251--258.

\bibitem{tw} P.~Tarrago and M.~Weber, Unitary easy quantum groups: the free case and the group case, \emph{Int.\ Math.\ Res.\ Not.\ IMRN} \textbf{2017} (2017), 5710--5750.

\bibitem{va} L.~Vainerman, $2$-cocycles and twisting of Kac algebras, \emph{Comm.\ Math.\ Phys.}\ \textbf{191} (1998), 697--721.

\bibitem{vdn} D.V.~Voiculescu, K.J.~Dykema and A.~Nica, Free random variables, AMS (1992).

\bibitem{wa95} S.~Wang, Free products of compact quantum groups, \emph{Comm.\ Math.\ Phys.}\ \textbf{167} (1995), 671--692.

\bibitem{wa} S.~Wang, Quantum symmetry groups of finite spaces, \emph{Comm.\ Math.\ Phys.}\ \textbf{195} (1998), 195--211.

\bibitem{wo1} S.L.~Woronowicz, Compact matrix pseudogroups, \emph{Comm.\ Math.\ Phys.}\ \textbf{111} (1987), 613--665.

\bibitem{wo2} S.L.~Woronowicz, Tannaka-Krein duality for compact matrix pseudogroups. Twisted ${\rm SU}(N)$ groups, \emph{Invent.\ Math.}\ \textbf{93} (1988), 35--76.

\end{thebibliography}

\end{document}
